% 本文件由 LaTeX 排版 Agent 根据有效 translation.json 自动生成。
% author=Gregory the Great, Pope (c. 540-604); work_id=3601; title=牧灵规则

\chapter{\WorkTitle{牧灵规则}}

% structure: p0001 source=h1 target=omitted reason=page-title-already-rendered-by-parent

\Emph{第一卷}\newline{}\Emph{第二卷}\newline{}\Emph{第三卷}\newline{}\Emph{第四卷}\par

\SourceRecord{Translated by James Barmby.；Nicene and Post-Nicene Fathers, Second Series；Vol. 12.；Edited by Philip Schaff and Henry Wace.；Buffalo, NY: Christian Literature Publishing Co.,；1895.；<http://www.newadvent.org/fathers/3601.htm>.}

% structure: p0001 source=h1 target=section reason=page-title

\section{司牧规章（卷一）}

% structure: p0002 source=h2 target=subsection reason=relative-heading-rank; base=h2

\subsection{序言}

本笃会版本从几部古代手稿中采用了\WorkTitle{司牧规章}这一书名，这正是额我略在将作品寄给其友塞维利亚的莱安德时所使用的称呼：\Quote{愿我将在我主教任期之初所写的\WorkTitle{司牧规章}一书……传送给您的圣德}（\WorkTitle{书信集}卷五，第49封）。此前更常用的\WorkTitle{司牧关怀书}可能取自书本身的开篇语：\Quote{司牧关怀的重担使我逃避，等等。}本书是在额我略主教任期之初发行的（如上述致莱安德书信引文所示），并且（如开篇所述）是写给拉文纳主教若望，以回复他收到的一封信。不过，尽管这次为特定目的而编撰成书，它必定源于长期思考，这一点从另一事实可以进一步证实：在他于君士坦丁堡居住期间开始并大部分完成的\WorkTitle{大道德论}（即\BibleBook{约伯}{记}注释）中，他已经勾画了这样一篇论文的计划，并表达了有朝一日将其编撰成书的希望。因为我们在那里已经写好了\WorkTitle{规章}第三卷的序言，以及该卷第一章所包含的大部分标题，随后是这些词句：\Quote{确实，我们应当特别指明在每一个问题上应如何安排劝诫的顺序；但担心冗长使我们却步。然而，在天主的助佑下，我们希望能在另一部作品中完成这一任务，只要我们这辛劳的生命还能剩下些许时间。}（\WorkTitle{道德论}卷三十，第12与13章）\par

此书在其作者生前似乎得到了应有的评价。正如我们所见，他将其寄给了塞维利亚的莱安德，显然是应后者之请，为了西班牙教会的益处；在书信中有一封来自该地区卡塔赫纳的博学主教李锡尼安写给额我略的信，信中高度赞扬了此书，尽管担心其中对主教职务资格的过高要求可能难以普遍达到（\WorkTitle{书信集}卷二，第54封）。皇帝莫里斯曾通过额我略在君士坦丁堡的执事阿纳托利索要并得到一本，并由安提约基雅宗主教阿纳斯塔西将其译为希腊文，后者本人也非常赞赏此书（\WorkTitle{书信集}卷十二，第24封）。此书似乎由隐修士奥斯定带到了英格兰。阿尔弗雷德大王在将近三百年后，在其神学家的帮助下，用西撒克逊语进行了翻译（或更确切地说是意译），并打算如他所说，将副本送到他王国的每一位主教手中。\par

在此之前，已有证据表明此书在高卢享有盛誉。在查理曼大帝于公元813年下令举行的一系列会议（即美因茨、兰斯、图尔和塞纳河畔沙隆会议）中，特别要求所有主教研读此书，连同新约圣经和教父教规。同样，在公元836年的亚琛会议上也是如此。此外，从兰斯总主教辛克马尔（公元845-882年）的一封信中可知，在主教祝圣时，会将此书的副本连同教规集一起放在祭台上交到主教手中，并嘱咐他们以此规范自己的生活。\par

这部作品完全配得上其古老声誉，是同类著作中的佼佼者，且对历代都有裨益。此前已有两部类似的作品。首先是纳齐安的额我略（约公元362年）的作品，即他的第二篇演说，称为\OriginalTerm{辩护词}{\GreekText{τοῦ} \GreekText{αὐτοῦ} \GreekText{ἀπολογητικός}}，与后来的额我略的作品一样，是为作者不愿接受主教职而辩护，并阐述该职务的责任。比较这两篇论文可知，前一篇启发了后一篇；事实上，教宗额我略在\WorkTitle{规章}第二卷的序言中也承认了他的借鉴。第二部类似的作品是金口若望的\WorkTitle{论司铎职}，共六卷，约公元382年。它也阐述了主教职务的可怕责任；但没有迹象表明教宗额我略从中汲取了内容。\par

值得注意的是，所有这些论文的主题都是主教职务，而不是现在通常理解的更广泛意义上的牧职或司铎职：并且值得注意的是，在额我略看来，讲道和精神指导灵魂的职责何等突出。它首先被视为治理的职务——\OriginalTerm{治理之位}{locus regiminis}、\OriginalTerm{治理之巅}{culmen regiminis}常被用来指称——因此纪律的实施占有突出地位；首席牧者也被视为在其羊群与天主之间的代祷者——参见例如卷一第十章——但通篇论述中，他尤其被描述为教师和灵魂的医生，其独特职责是通晓各种形式的精神疾病，从而能够根据不同情况施以相应的治疗，以\BibleQuote{宣讲圣言，以全备的忍耐和教训，进行责备、警戒、劝勉}，并通过言传身教引导灵魂走上救恩之路。额我略没有白白研读圣保禄的\BibleBook{保禄}{牧函}。确实，他在整部作品中展现出深刻的辨识力，洞察人性特质与动机，以及环境或气质使不同人群（包括牧者及其羊群）特别容易受到的诱惑。同样引人注目的是，他在这部作品及其他作品中，对整部圣经的熟悉程度。他确实仅通过拉丁译本了解圣经；他的文本考据知识常有失误；而且他喜爱的牵强神秘解释比比皆是。但作为圣经整体道德与宗教教义的真实阐释者，他配得上教会伟大圣师之一的称号。此外，尽管他对大公会议和教父在信仰问题上持有至高权威的尊崇，但他始终诉诸圣经作为行为与信仰的最终权威。\par

% structure: p0008 source=h2 target=subsection reason=relative-heading-rank; base=h2

\subsection{导言}

\Signature{额我略致其最可敬、最神圣的兄弟及同为主教的若望。}\par

最亲爱的兄弟，你以善意谦逊的态度责备我，因我本想藉隐退逃避牧职的重担；对此，为免某些人以为这些重担轻省，我在呈于你眼前的书中以笔墨表达了我对它们沉重性的全部评估，好使尚未背负此担者不致轻率追求，而已经追求者则因得到它而战栗。本书分为四个独立的论证部分，以便藉按序排列的陈述——犹如藉着某种阶梯——进入读者的心灵。因为，按事物的需要，我们尤当考虑：每人在何种方式下应达至最高治理权；既已适时达至，又当如何生活；生活良好，又当如何教导；教导正确，又当每日何等审慎地察觉自己的软弱；以免谦逊远离这接近，或生活与这抵达相悖，或教导有亏于生活，或自矜不当抬高教导。因此，愿恐惧来抑制欲望；此后，那未经追求而获权柄者，其生活应彰显此权柄。但继而，牧者生活中显现的善亦须藉其言语传播。最后，无论完成了何等善工，我们应因反省自身的软弱而为之自卑，以免隐密审判前骄傲的膨胀甚至将这些善工熄灭。然而，有许多像我一样不娴熟的人，他们不知如何衡量自己，却贪图教导自己所未学之事；他们愈是不知权柄之伟大压力的沉重，便愈是轻视它的重担；愿这些人从本书之初就受到责备，好使他们在无知而仓促中渴望掌握教导的高位时，被我们的论述拒之门外，远离他们鲁莽的冒险。\par

% structure: p0011 source=h2 target=subsection reason=relative-heading-rank; base=h2

\subsection{第一章}

不娴熟者不得擅进职掌权柄之位。\par

没有人敢在尚未以专心思索学习一门技艺之前，就去教授它。那么，不娴熟者擅自承担牧职权柄，岂不是鲁莽至极？因为治理灵魂是诸艺之艺！谁能不知，人心思虑的创伤比腹内的创伤更为隐晦？然而，多少对属灵诫命一无所知的人，竟无畏惧地自诩为心灵的医师，而那些不明药效的人，却羞于以肉身的医师自居！但因着上主的安排，今世所有位高权重者都倾向于尊敬宗教，于是在神圣教会内，有些人藉着统治的外表贪图尊荣的荣耀。他们渴望显得像导师，他们贪求凌驾他人，且如真理所证，他们在街市上寻求初次的问候，宴席上的首座，集会中的首位（\BibleRef{玛 23:6}）。正因他们只是出于骄矜而达到谦逊的导师地位，反而愈不能胜任其所承担的牧职。诚然，在导师的地位上，若所学与所教不同，言语本身便混乱了。对此，主藉先知抱怨说：\BibleQuote{他们没有经过我的许可就作王，没有经我的任命就作首领}（\BibleRef{欧 8:4}）。因为那些自行为王而非出于至上统治者旨意的人，他们不受任何德行的支持，也非出于天主召唤，而是被自身私欲所煽动，强夺而非达到统治高位。然而，那位内在的审判者既提升他们，又不认识他们；因为祂以容忍许可他们，却以弃绝的审判视他们为无有。因此，祂对那些即使行了奇迹才来到祂面前的人说：\BibleQuote{你们这些作恶的人，离开我，我不认识你们是谁}（\BibleRef{路 13:27}）。牧者们的无能受到真理之声的斥责，正如先知所言：\BibleQuote{牧者自己没有理解}（\BibleRef{依 56:11}）；主又谴责他们说：\BibleQuote{那些掌握法律的，没有认识我}（\BibleRef{耶 2:8}）。因此，真理抱怨自己不被他们所认识，并宣告祂不认识那些不认识祂之人的统治；因为事实上，那些不认识主之事物的人，主也不认识他们；正如保禄所证，他说：\BibleQuote{但若有人不认识，他也不被认识}（\BibleRef{格前 14:38}）。然而，牧者的这种无能无疑常与属下的功过相称，因为，虽然他们自己没有知识之光乃是自己的过错，但这也是严格审判的安排：因着他们的无知，那些跟随他们的人也会跌倒。因此，福音中真理亲自说：\BibleQuote{如果盲人领盲人，两人都掉进坑里}（\BibleRef{玛 15:14}）。因此，圣咏作者（不是表达自己的愿望，而是藉其先知职分）斥责这样的人，他说：\BibleQuote{让他们的眼睛瞎了，看不见，并使他们的脊背永远弯曲}（\BibleRef{咏 68:24}）。诚然，那些人是眼睛，他们被安置在最崇高尊严的面容前，承担了侦察道路的职责；而那些依附他们并跟随他们的人被称为脊背。所以，当眼睛瞎了，脊背就弯曲；因为，当先行者失去知识之光时，跟随者就被压弯以背负他们罪恶的重担。\par

% structure: p0014 source=h2 target=subsection reason=relative-heading-rank; base=h2

\subsection{第二章}

凡未经在生活中实践所学之事者，不应进入治理之位。\par

也有些人不懈地细心探究属神的规诫，但他们以理解力所洞察的，却在生活中践踏；他们忽然间教导那些并非经由实践而是通过研习所学到的东西；他们以言语宣讲的，却以行为加以攻击。由此便发生：当牧者行经陡峭之处时，羊群便跟随而至悬崖。因此，主通过先知抱怨牧者那可鄙的知识，说：\BibleQuote{当你们自己饮了最纯净的水，却用脚搅浑了剩余的水；我的羊吃了被你们脚践踏过的，喝了被你们脚搅浑过的水。}\BibleRef{则 34:18} 的确，牧者饮最纯净的水，是指他们以正确的理解汲取真理的溪流。但用脚搅浑这水，是指以邪恶的生活败坏圣善默想的研习。而羊喝被脚搅浑的水，是指那些属下之人不听从所听到的话，只模仿所见的坏榜样。他们渴慕所说的事，却因所见的行为而败坏，如同从污染的泉源中饮水时连同泥沙一同咽下。因此，又通过先知写道：\BibleQuote{恶司铎是我子民绊倒的罗网。}\BibleRef{欧 5:1} 再者，主通过先知论及司铎说：\Quote{他们成了以色列家的罪恶绊脚石。}因为确实，在教会中，无人比那拥有圣善之名与品位而行邪僻之事者造成更大的伤害。因为当他违犯时，无人敢责难；而由于对他的品位的尊敬，罪人受尊荣，恶行便因榜样而强力传播。但所有不配者若以心灵倾听的耳朵权衡真理之言，便会逃避如此重大罪责的负担：\BibleQuote{谁若使这些信我的小子中的一个跌倒，倒不如把一块驴拉的磨石拴在他的颈上，沉到海底里。}\BibleRef{玛 18:6} 磨石表示世俗生活的圆环与劳苦，海底则表示最终的惩罚。因此，凡已承担圣善的外表，而以言语或行为毁灭他人者，实在不如让他以公开的形貌受世俗行为压迫至死，胜过他以神圣的职务使他人效法他的恶行；因为，如果他独自跌倒，地狱的痛苦尚能以较可忍受的程度折磨他。\par

% structure: p0017 source=h2 target=subsection reason=relative-heading-rank; base=h2

\subsection{第三章}

\Emph{论治理的重担；以及要轻视一切逆境，畏惧顺境。}\par

如此，我们已简要说明治理的重担何等重大，以免任何不胜任神圣治理职务的人胆敢亵渎它们，并因贪图高位而承担导致毁灭的领导权。因此，雅各伯深情地劝阻我们说：\BibleQuote{我的弟兄们，不要多人做师傅。}\BibleRef{雅 3:1} 因此，天主与人之间的中保——祂超越甚至天上神体的知识与理解，从永远在天上为王——在地上却逃避接受王国。因为经上记载：\BibleQuote{耶稣既知他们要来强迫祂，立祂为王，就独自又退到山上去了。}\BibleRef{若 6:15} 因为谁能像祂那样无可指摘地拥有对人的统治权呢？祂本应统治祂自己所创造的人。但是，祂降生成人，正是为了不仅以苦难救赎我们，也以祂的言行教导我们，将自己作为榜样赐给追随者，因此祂不愿被立为王；反而自愿走向十字架的刑具。祂逃避了别人献上的高位荣耀，却渴望受凌辱的死亡之苦；这样，祂的肢体就能学会逃避世俗的恩惠，不畏惧任何恐怖，为真理而爱逆境，且畏惧顺境；因为顺境常因虚荣而玷污心灵，而逆境却以忧伤净化心灵；顺境中心灵自高自大，逆境中则即便曾自高，也使自己谦卑；顺境中人忘记自己，逆境中则即使勉强且违背意愿，也被迫记起自己的本相；顺境中连往昔所做的好事也常化为乌有，逆境中则连长久的过犯也被抹除。因为通常，在逆境的学校里，心灵在管教下被驯服，而一旦突然获得至高的统治权，便立即改变，因熟悉荣耀而傲慢。这样，撒乌耳先前因自觉不配而逃避，一旦执掌王国的统治权，便骄傲起来\BibleRef{列上 10:22}；因为他渴望在人民面前受尊荣，不愿公开受责，竟将那位曾为他傅油立王者也断绝于己。同样，达味在拣选他的那位眼中，几乎在所有行为中都中悦祂，但一旦压力的重担移去，便迸发为肿胀的疮口\BibleRef{列下 11:3}，先是对妇人贪欲而放纵如脱缰，后便对那人残酷无情地杀害；他先前知道怜悯地宽恕恶人，后来却毫无犹豫地渴求好人的死亡\BibleRef{列下 11:15}。的确，他先前不愿击打被捕的迫害者，后来却使疲惫的军队受损，毁灭了他忠心的士兵。事实上，他的罪行本会将他从选民中远远抛开，若不是鞭笞将他召回求恕。\par

% structure: p0020 source=h2 target=subsection reason=relative-heading-rank; base=h2

\subsection{第四章}

\Emph{论治理的职责多半会涣散心灵的坚毅。}\par

治理的重担在承担时常常使人心神涣散，既难以应对个别事务，又因思绪混乱而分心于诸多事务。因此，有一位智者明智地劝诫说：\BibleQuote{\Emph{我儿，不要管许多事}}\BibleRef{德 11:10}，因为心志若分散于各种事务，便无法集中于任何单一工作计划。当心被不寻常的忧虑牵引向外时，便失去了内在敬畏的稳固；它忙于安排外在事物，唯独对自身无知，虽能思考许多事，却不认识自己。因为，当它过分涉足外在事务时，犹如在旅途中忙于赶路而忘记目的地；如此，由于疏于自我省察，甚至不考虑自己所受的损失，也不知道损失有多大。因为\Person{希则克雅}也并不相信自己犯罪\BibleRef{列下 20:13}，他向来到他面前的陌生人展示了他的香料库；但他却因自己以为合法的事\BibleRef{依 39:4}，而落在审判者的愤怒之下，为后代子孙招致定罪。常常，当资财丰富、能为属下做出许多令人赞叹的事时，心便在思虑中自高，虽未公然行出恶行，却已充分激起审判者的愤怒。因为审判者在内，被审判者也在内。因此，我们在心中犯罪时，所行之事虽对人隐藏，但在审判者眼中我们却犯了罪。因为巴比伦王并非在说出狂傲之言时才首次显为有罪\BibleRef{达 4:16}；其实在他尚未吐出狂傲之言以前，就已从先知口中听见了谴责的判决。因为他后来宣告他所冒犯的全能天主给所有在他统治下的万民时，才洗除了自己所犯的骄傲之罪。\par

但此后，因统治的成功而得意，并因成就大事而欢喜，他先在心中把自已置于众人之上，接着仍怀虚荣地说：\BibleQuote{\Emph{这大巴比伦岂不是我用我的大能建造为王国之殿，为我威严的荣耀吗？}}\BibleRef{达 4:30}。我们看到，他的这番话公然招致了由内心骄矜所燃起的愤怒的报复。因为严厉的审判者先无形中看见，后来以公开的惩治予以责备。为此，天主甚至将他变成无理性的牲畜，使他离开人类社会，改变他的心意，使他与田野的走兽为伍，以便在明显严厉而公义的审判中，那自视高于世人的人连人性也丧失。我们提出这些事，并非指责治理，而是防范心灵因贪图治理而软弱，以免有不完备的人冒险攫取最高统治权，或已在平地上绊跌的人再踏上悬崖。\par

% structure: p0024 source=h2 target=subsection reason=relative-heading-rank; base=h2

\subsection{第五章}

\Emph{论那些能以德表在最高治理中利益他人，却为求自身安逸而逃避的人。}\par

因为有些人卓著地拥有各种德行，且因伟大的恩赐而被高举以训练他人；他们在贞洁的热忱中纯洁，在节制的力量中坚强，充满教义的盛宴，在忍耐的宽容中谦卑，在权威的刚毅中挺立，在慈爱的恩宠中温柔，在正义的严厉中严格。诚然，这样的人，若在蒙召时拒绝接受最高治理的职务，多半会剥夺自己领受的恩赐——这些恩赐并非仅为己身，也是为他人；而且，当他们只图谋自己的利益而不顾他人时，便丧失了那些他们想为自己保留的益处。因此，真理曾对祂的门徒说：\BibleQuote{\Emph{城造在山上是不能隐藏的；人点灯，不放在斗底下，而是放在灯台上，就照亮全屋子的人。}}\BibleRef{玛 5:15}。因此祂对\Person{伯多禄}说：\BibleQuote{\Emph{西满，若望的儿子，你爱我吗？}}\BibleRef{若 15:16}；当他立即回答爱时，被告知：\BibleQuote{\Emph{如果你爱我，就喂养我的羊。}}如果喂养的关怀是爱的证明，那么谁若富有德行，却拒绝喂养天主的羊群，便被证明不爱那大牧者。因此\Person{保禄}说：\BibleQuote{\Emph{如果基督替众人死了，那么众人就都死了；祂替众人死，是为使活着的人不再为自己活，而是为替他们死而复活的那位活。}}\BibleRef{格后 5:15}。因此\Person{梅瑟}说\BibleRef{申 25:5}，若兄弟死了没有儿子，他的兄弟应娶他的妻子，为他的兄弟生子立后；若他拒绝娶她，那妇人就唾在他脸上，她的亲属要脱去他一只脚上的鞋，并称他的家为脱鞋者之家。已故的兄弟就是那位在复活的光荣后说\BibleQuote{\Emph{去告诉我的弟兄们}}\BibleRef{玛 28:10}的那一位。祂死时好像没有孩子，因为祂尚未满全祂所拣选的人数。然后，命令活着的兄弟娶那妻子，因为将圣教会的关怀托付给能最好治理她的人，无疑是合宜的。但他若不愿意，那妇人就唾在他脸上，因为谁不愿用他所领受的恩赐去造福他人，圣教会就指责他所有的好东西，如同把唾沫吐在他脸上；从一只脚上脱去鞋子，因为经上记载：\BibleQuote{\Emph{用和平福音的预备穿在脚上}}\BibleRef{弗 6:15}。因此，若我们既关心自己又关心邻人，我们就两只脚都穿上了鞋。但那只图谋自己的利益而忽视邻人利益的人，就丢了一只鞋而蒙羞。因此，如我们所说，有些人拥有伟大恩赐，但他们只热心于默观的学问，却退缩不愿通过宣讲服务邻人的益处；他们喜欢安静的隐秘处，渴望默想的退隐。关于这种行为，若严格判断，他们无疑因他们的恩赐之大——本来可以公开造福他人——而更有罪。因为当至高的圣父的独生子亲自从父怀中来到我们众人中间，为造福许多人时，一个人本可显著造福邻人，却宁愿自己的独处胜过他人的益处，这是何等的心态呢？\par

% structure: p0027 source=h2 target=subsection reason=relative-heading-rank; base=h2

\subsection{第六章}

\Emph{那些因谦逊而逃避治理的重担的人，当不抗拒天主的命令时，才是真正谦逊的。}\par

也有些人是纯因谦逊而逃避，唯恐自己会被置于那些他们认为不如自己的人之上。确实，他们的谦逊，若有其他德行相配，在天主眼中才是真正的谦逊——当它不固执地拒绝那被命令去有益地承担的职务时。因为，一个人若明白上主之意的美善应当主导，却藐视其主导，他并不真正谦逊。相反，顺服于天主的安排，远离固执的恶习，若他已有恩赐可以造福他人，那么，当被命令承担最高治理时，他应在心中逃避，但违背己意而服从。\par

% structure: p0030 source=h2 target=subsection reason=relative-heading-rank; base=h2

\subsection{第七章}

\Emph{有时有些人令人称羡地渴望宣讲的职务，而另一些人则同样令人称羡地被强制拉入其中。}\par

虽然有时一些人可称赞地渴望讲道的职务，但另一些人同样可称赞地被强迫拉入其中；如果我们考虑两位先知的行为，就清楚看到这一点，其中一位自愿献身被派遣去宣讲，另一位却因害怕而拒绝前往。因为当主问祂该派遣谁时，\Person{依撒意亚}自愿献身说：\Emph{\BibleQuote{我在这里；请派遣我}} \BibleRef{依 6:8}。但\Person{耶肋米亚}却被派遣，却谦卑地恳求不要派遣他，说：\Emph{\BibleQuote{哎！我主天主！我实在不会说话，因为我还年轻}} \BibleRef{耶 1:6}。看，这两人外表发出不同的声音，却来自同一个爱泉。因为爱德有两条诫命：爱天主和爱近人。因此，\Person{依撒意亚}渴望通过活跃的生活造福近人，因而渴望讲道的职务；而\Person{耶肋米亚}渴望通过默观生活殷勤地粘着于造物主的爱，因而反对被派遣去宣讲。所以，一个可称赞地渴望的，另一个却可称赞地退缩；后者唯恐因说话而失去静默默观的收获；前者唯恐因保持沉默而因缺乏勤勉工作遭受损失。但在这两种情况下都要仔细注意：拒绝者没有坚持他的拒绝，而渴望被派遣的人看到自己先前已被祭坛上的火炭洁净；免得任何未被净化的人胆敢接近神圣职务，或任何被天上恩宠拣选的人以谦卑的外表骄傲地反驳它。因此，既然任何人很难确定自己已被洁净，那么拒绝讲道的职务是更安全的，尽管（如我们所说）当认出天主要求承担此职务时，不应固执地拒绝。\Person{梅瑟}奇妙地实现了这两个要求，他不愿意被置于如此大的群众之上，却仍服从了。因为如果他毫不颤抖地承担那无数百姓的领导，他或许会是骄傲的；而且，如果他拒绝服从主的命令，显然也会是骄傲的。因此，他两方面都谦卑，两方面都顺服：他衡量自己，不愿被置于百姓之上；然而，他信赖命令他的那一位的德能，于是同意了。所以，那么，所以让所有鲁莽的人推断他们的罪有多大，如果他们不怕自己寻求超越他人，而当圣人们即使在天主命令时也害怕承担领导人民之责。\Person{梅瑟}在天主劝说他时仍颤抖；然而每个软弱的人都渴望承担尊位的重担；一个几乎不能承受自己负荷而不跌倒的人，却欢喜地将肩膀置于他人（而非自己）的压力之下：自己的行为对他而言太重而难以承担，他却加重了负担。\par

% structure: p0033 source=h2 target=subsection reason=relative-heading-rank; base=h2

\subsection{第八章}

\Emph{论贪求高位并以宗徒的话来满足自己贪欲的人。}\par

但大多数贪求高位的人，抓住宗徒的话来满足自己的贪欲，宗徒说：\Emph{\BibleQuote{谁若渴望主教的职务，就是渴望一件善工}} \BibleRef{弟前 3:1}。但他在称赞这渴望的同时，立即将所称赞的转为恐惧，因为他紧接着补充说：\Emph{\BibleQuote{但主教必须无可指责}} \BibleRef{弟前 3:2}。随后当他列举必要的德行时，他表明了这无可指责包含什么。因此，对于他们的渴望，他予以认可，但通过他的诫命使他们警惕；好像明说：我称赞你所寻求的，但首先要了解你所寻求的是什么；免得在你忽略衡量自己时，你的可指责性因急于在最高荣誉位子上被众人看见而显得更加丑陋。因为这位治理艺术的大师，通过赞许来推动，通过警告来制止；这样，通过描述无可指责的高度，他抑制听众的骄傲，并通过称赞所寻求的职务，使他们向往所要求的生活。然而要注意，这话是在一个当时被置于人民之上的人通常首先被引向殉道之苦的时候说的。因此，那时寻求主教的职务是值得称赞的，因为通过它，一个人无疑最终会遭遇更重的痛苦。故此，甚至主教的职务本身也被定义为善工，如经上所说：\Emph{\BibleQuote{谁若渴望主教的职务，就是渴望一件善工}} \BibleRef{弟前 3:1}。所以，那寻求的不是这善工的职务，而是出人头地的荣耀，他就是自己反对自己的证人，证明他并不渴望主教的职务；因为那个人不仅根本不爱神圣的职务，甚至不知道它是什么，他渴望最高的统治，在内心隐秘的默想中以他人的服从为食，为自己的赞美而欢喜，将心高举到尊荣，因丰富的财富而狂喜。这样，世俗的利益就在那本应摧毁世俗利益的荣誉的幌子下被寻求；而当心灵想抓住谦卑的最高位置以抬高自己时，它内在地改变了它所外在地渴望的。\par

% structure: p0036 source=h2 target=subsection reason=relative-heading-rank; base=h2

\subsection{第九章}

\Emph{论想要高位的人多半以虚假的善工承诺自欺其心。}\par

但大体而言，那些贪求牧灵权位的人，心里也自忖着一些善工，尽管是出于骄傲的动机，仍会幻想自己将成就大事。于是，内心深处隐藏的动机是一回事，浮现在思想表面的又是另一回事。因为人的心对自己撒谎，假装爱那其实不爱的善工，又假装不爱那其实爱着的世俗荣耀。他们渴望统治，追求时却畏首畏尾，一旦得到便变得狂妄。因为在追求过程中，他们战战兢兢，唯恐得不到；但一旦得到，便视之为应得之物。当他们开始以世俗方式享用所得权位的职务时，便轻易忘记了先前在宗教热忱中思虑过的事。因此，当这样的思虑超乎寻常时，必须将心灵之眼召回已完成的善工，每个人都该思考自己身为下属时做了什么；这样他才能立即发现，作为上司，他是否真能行出所计划的好事。因为一个在低位时尚未停止骄傲的人，绝不可能在高位学会谦卑；一个在缺乏赞美时尚且渴望得到它的人，绝不可能在赞美满溢时躲避它；一个连自己一人所需尚且不足的人，绝不可能在要供养多人时战胜贪婪。因此，每个人都应从过去的生活发现自己究竟是怎样的人，免得在渴求高位时，被自己思想的幻影所欺骗。然而，通常的情形是：在平静中持守的善行，在治理的职责中反而丧失了；因为即使在风平浪静的海上，一个不熟练的人也能直线航行；但在风浪颠簸的海上，就连熟练的水手也会迷失。那崇高的权位，难道不就是一场心灵的风暴吗？心思之船被思想的飓风不断摇撼，被抛来抛去，以致在言语和行为的突发过犯中，如同撞上礁石而破碎。在这重重危险中，当如何行？当持守什么？无非是：那拥有丰富德行的人，应在不得已时接受治理；那毫无德行的人，即使被迫也不可接近治理。至于前者，要小心，若他全然拒绝，恐怕就像那将所得银子包在毛巾里的人，因埋藏而受审（\BibleRef{玛 25:18}）。因为将银子包在毛巾里，就是将所领受的恩赐隐藏于懒散的昏暗中。另一方面，后者若渴求治理，要当心，免得他因恶行的榜样，成为那些走向天国之门者的绊脚石，就像法利塞人那样，按主的话（\BibleRef{玛 23:13}），自己不进天国，也不让别人进去。他还当思考：当一位被选的主教承担起人民的责任时，他便如同医生去医治病人。倘若病痛仍活在他自己身上，那他便多么狂妄，竟急忙去医治受伤者，而自己脸上还带着疮痍！\par

% structure: p0039 source=h2 target=subsection reason=relative-heading-rank; base=h2

\subsection{第十章}

\Emph{什么样的人应当来治理。}\par

因此，那已经度属灵生活、死于一切肉身情欲的人，必须用绳索拉来作善生的榜样；他无视世俗的顺境，不惧怕任何逆境，只渴求内在的财富；他的身体与灵魂和谐一致，既不被身体的软弱所阻挡，也不被灵魂的傲慢所抑制；他不贪图别人的东西，却慷慨施舍自己的；他因慈悲心肠而迅速宽恕，却从不因过度宽恕而偏离正直的堡垒；他不做任何不法之事，却为别人所行的不法哀叹，如同自己的事；他因心中的爱而同情他人的软弱，为他人的益处而喜乐，如同自己的得益；他所作所为都成为别人的榜样，以致在众人中，至少对自己过去的行为无可羞愧；他努力如此生活，以致能用教义的溪流浇灌干涸的心灵；他已经通过祈祷的实践和考验，学会能从主那里获得所求，正如经验之声已向他说的：\BibleQuote{你还在说话时，我要说：我在这里}（\BibleRef{依 58:9}）。因为，若有人来找我们，请我们为他向某位与他有隙的贵要人物求情，而对那人并不熟悉，我们必会回答：\Quote{我们不能为你代求，因为与那人并不熟识。}那么，若一个人因与别人没有关系，便羞于为他代求，那又有何理由能抓住为百姓向天主代求的职务呢？他怎能知道自己因生活的功德而蒙主恩宠呢？他又怎能求主赦免别人，而不知自己是否已与主和好？此外，还有一事更令人担忧：即那被认为能平息主怒的人，反而因自己的罪过激起主的怒气。因为我们都知道，当一个失宠的人被派去向愤怒者求情时，后者的心会被激怒得更严厉。因此，那仍被世俗欲望捆绑的人要当心，免得因更严重地激怒严正的审判者，而当他喜乐于尊位时，反倒成了下属灭亡的原因。\par

% structure: p0042 source=h2 target=subsection reason=relative-heading-rank; base=h2

\subsection{第十一章}

\Emph{什么样的人不应当来治理。}\par

因此，愿每个人明智地度量自己，免得他敢于承担治理之位，而在他内恶习仍掌权以致定罪；免得一个被自身罪疚败坏的人想要成为他人过犯的代祷者。为此，有上天之声对梅瑟说：\BibleQuote{你告诉亚郎：世世代代中，凡你后裔身上有残疾的，不可向主、他的天主奉献饼。}\BibleRef{肋 21:17}。随即又补充说：\BibleQuote{如果他是瞎子，如果他是瘸子，如果他的鼻子是小的或是大而弯曲的，如果他是折脚或折手的，如果他是驼背的，如果他是\OriginalTerm{烂眼}{lippus}，如果他的眼中有\OriginalTerm{白斑}{albuginem}，如果他有慢性疥疮，如果他身上有脓疮，或者如果他患有\OriginalTerm{疝气}{ponderosus}。}\BibleRef{肋 21:18}。因为那个人确实是瞎子，他不熟悉超性默观的亮光，被现世生命的黑暗吞没，由于不因爱慕而仰望未来的光明，他不知道自己的行为脚步要迈向何处。因此，亚纳在预言时说：\BibleQuote{他必保护他圣徒的脚，恶人将在黑暗中沉默。}\BibleRef{列上 2:9}。但那个人是瘸子，他确实看见自己该走的方向，但因意志薄弱，不能完满地持守他所看见的生命之路，因为不稳定的习惯未能上升到德性的稳固状态，行为的脚步不能有效地跟随渴望的目标。因此，保禄说：\BibleQuote{你们要举起下垂的手，软弱无力的膝，为你们的脚修直道路，免得那瘸的偏离正道，反而应被治好。}\BibleRef{希 12:12}。但鼻子小的人，是不适合保持辨别尺度的人。因为我们用鼻子辨别香气和臭味：因此，鼻子恰当地表达了辨别力，我们藉此选择德行而避免罪恶。因此，在赞美新娘时也说：\BibleQuote{你的鼻子如同黎巴嫩上的塔楼。}\BibleRef{歌 7:4}。即圣教会藉辨别力窥探来自各方的攻击，并从高处察觉即将来临的恶习之战。但有些人，不愿被认为迟钝，常常不必要地忙于各种探究，因过度精巧而受骗。因此又加上：\BibleQuote{或有一个大而弯曲的鼻子。} 因为大而弯曲的鼻子是过度的辨别精巧，它变得过度增生，本身混淆了自己运作的正确性。但折脚或折手的人，是根本不能行走在天主之道上，完全无分于善行的人，甚至他不能像瘸子那样勉强持守，而是完全与之分离。但驼背的人，是受尘世挂虑的重担压弯的人，他从不仰视上面的事物，只专注于最卑微的事物中被践踏的东西。而且，他即使偶尔听到关于天上家乡的美好事物，也被邪僻习惯的重担压制，以致不举起他心灵的面孔朝向它，不能挺立他思想的姿态，因为尘世挂虑的习惯使其向下弯曲。关于这种人，圣咏作者说：\BibleQuote{我常被压弯，被降为卑微。}\BibleRef{咏 37:8}。真理亲自斥责这类人的过错，说：\BibleQuote{那落在荆棘中的种子，是指那些人，他们听了道，就走开，被挂虑、财富和生活的享乐所窒息，结不出果实。}\BibleRef{路 8:14}。但烂眼的人，是他的天然才智闪现出对真理的认识，但肉性行为却遮蔽了它。因为在烂眼的人中，瞳孔是健全的，但眼睑因体液下流而虚弱，变得粗厚；甚至瞳孔的明亮也受损，因为不断有液体流过而磨损。因此，烂眼的人是一个天性感觉敏锐，但被堕落的生活习惯所混乱的人。对这样的人，藉天使说得很好：\BibleQuote{用眼药擦你的眼睛，使你能看见。}\BibleRef{默 3:18}。因为我们可说是用眼药擦我们的眼睛以便看见，当我们以良好行为的药膏帮助我们理智的眼目感知真光的清明时。但眼中有白斑的人，是不被允许看见真理之光的人，因为他被傲慢自诩的智慧或正义所蒙蔽。因为眼睛的瞳孔，当黑色时能看见；但当它有白斑时，什么也看不见；由此我们可以理解，人的思想的感知意识，如果人自知是愚人和罪人，就能认识内在的亮光；但如果它自以为是正义或智慧的白色，就自我排除于来自上界的知识之光，并且它越因傲慢而自我抬高，就越完全不能穿透真光的清明；正如对有些人所说：\BibleQuote{他们自称聪明，反而成了愚人。}\BibleRef{罗 1:22}。但患有慢性\OriginalTerm{疥疮}{scabies}的人，是不断被肉体的淫欲所征服的人。因为在疥疮中，内脏的灼热被引到皮肤上；因此恰当指代淫欲，因为，如果心的诱惑爆发为行动，可以真实地说内在的灼热发作成皮肤的疥疮：它现在向外伤害身体，因为，当感官不受思想约束时，它也在行动上取得了控制。因为保禄关心要洗净这皮肤的瘙痒，他说：\BibleQuote{不要让你们受到考验，除非是人所能受的。}\BibleRef{格前 10:13}。仿佛要清楚地说：在内心受诱惑是人性的；但在诱惑的争战中，也在行动上被征服，则是魔鬼性的。他身上有\OriginalTerm{脓疮}{impetigo}的人，是心灵被贪婪蹂躏的人；这贪婪，若在小事上不被克制，就确实无限制地膨胀。\par

因为，正如\Emph{\OriginalTerm{脓疱疮}{impetigo}}毫无痛苦地侵入身体，蔓延时也不给患者带来烦恼，却毁坏了肢体的美丽；同样，贪恋也愉悦着被它束缚之人的心灵，同时使其溃烂。由于它向思想提供一件又一件可获取之物，便点燃了敌意的火焰，对所造成之伤不给予痛苦，因为它向狂热的心灵许诺从罪恶中获得丰裕。但肢体的美丽被毁坏了，因为其他德行的美丽也因此受损；它好似使全身溃烂，因为它以各种恶习败坏心灵；正如保禄作证说：\BibleQuote{贪爱钱财是万恶的根源}（\BibleRef{弟前 6:10}）。但破裂之人，是指那些不将卑鄙付诸行动，却因不断思虑而在心中被其过度压抑的人；他确实丝毫没有走到恶行的地步，但他的心仍毫无抵触地沉溺于淫乐的欢愉。因为破裂之病是当\Emph{\OriginalTerm{体液流向生殖器，并确实随着可耻的不适而肿胀}{humor viscerum ad virilia labitur, quæ profecto cum molestia dedecoris intumescunt}}时。因此，凡使所有思想流向淫荡，心中担负着卑鄙重负的人，可称为破裂者；尽管没有实际可耻行为，但在思想上仍未离开它们。他也无力在众人面前起身行善，因为隐藏在他内的可耻重负压制着他。因此，凡受其中任何一种疾病侵袭的人，都禁止向主献上饼，以免他仍被自己的罪恶所蹂躏，从而无法为他人赎罪。\par

如今，已简要说明了配得之人应如何获得牧者权威，以及不配之人应如何大为恐惧，现在让我们来示范，堪当获得此权威的人应如何在其位生活。\par

\SourceRecord{Translated by James Barmby.；Nicene and Post-Nicene Fathers, Second Series；Vol. 12.；Edited by Philip Schaff and Henry Wace.；Buffalo, NY: Christian Literature Publishing Co.,；1895.；<http://www.newadvent.org/fathers/36011.htm>.}

% structure: p0001 source=h1 target=section reason=page-title

\section{\WorkTitle{牧灵规程（卷二）}}

论牧者之生活\par

% structure: p0003 source=h2 target=subsection reason=relative-heading-rank; base=h2

\subsection{第一章}

\Emph{一位已按正当次序到达治理职位者，应如何在其中自处。}\par

牧者的行为应如此超越百姓的行为，正如牧人的生活常超乎羊群之上一般。因为那被众人视为牧者的人，其身份既使百姓被称为他的羊群，他就必须殷切思虑自己身负何等重大的责任，以保持正直。因此，他在思想上应当纯洁，在行动上应当卓越；在沉默上谨慎，在言语上有益；在同情上与每个人亲近，在默观上超乎众人之上；因谦逊而成为善人的密友，因正义的热忱而对恶人的恶习毫不妥协；不可因外在事务而放松对内在的关怀，也不可因内心的挂虑而忽略对外在事务的照顾。现在，我们应将这些简略提及的事更详尽地阐述和讨论。\par

% structure: p0006 source=h2 target=subsection reason=relative-heading-rank; base=h2

\subsection{第二章}

\Emph{论治理者应在思想上保持纯洁。}\par

治理者应常在思想上保持纯洁，因为那承担了为他人心灵擦去污秽职责的人，自己绝不能沾染任何不洁；因为要清洁污垢的手必须洁净，否则，如果它本身被泥垢弄脏，只会更加弄脏它所接触的一切。为此，先知说：\BibleQuote{凡抬主器皿的，要洁净自己}\BibleRef{依 52:11}。他们以自己的生活为保证，负责将邻人的灵魂引向永恒圣所的人，就是抬主器皿的人。因此，他们应在自己内心感知，那些将活生生的器皿托付于自身责任之怀中带往永恒圣所的人，应当是多么洁净。因此，天主的声音命令\BibleRef{出 28:15}，亚郎的胸前应紧贴判牌，用系带扎紧；因为司铎的心绝不应被松散的思虑所占据，而唯独理性约束它；那被立为他人表率的人，不应思量任何轻率或无益的事，而应在生活的庄重中显明自己胸中承载着何等丰富的理性。并且在这判牌上，又仔细规定要刻上十二位族长的名字。因为常将列祖的名字在胸前记载，就是无间断地默想古人的生命。当司铎不断思考先辈的榜样，不停默想圣人的足迹，抑制邪念，免得他行为的脚步超越秩序的界限时，他才能无可指责地行走。它也称为判牌，因为治理者应时常以敏锐的辨别力分辨善恶，殷勤思考何事适合于何情、何时、如何；他不应为自己寻求什么，而应将邻人的益处视为自己的利益。因此，同处写着：\BibleQuote{你要将训诲和真理放在亚郎的胸牌上，这胸牌要常在亚郎胸前，当他进到主面前时，他要常在主面前将以色列子民的判断带在胸前}\BibleRef{出 18:30}。司铎在主面前将以色列子民的判断带在胸前，意思是他只按内心审判者的心意来审查属下的案件，使他作为天主的代表行事时不沾染任何人性的成分，免得私人的恼怒刺激他惩戒的锐气。当他热忱反对他人的恶习时，也应除去自己的恶习，免得潜在的怨恨败坏他判断的平静，或轻率的怒气扰动它。但当那主宰万物的主的威仪（即内心的审判者）被考虑时，属下若不是带着极大的敬畏，就不能被治理。这样的敬畏确实洁净治理者的心，同时使他谦卑，保护他不致因精神的傲慢而高举，或因肉体的享乐而玷污，或因贪恋世俗事物而为其尘垢所遮蔽。但这些事物不可能不叩击治理者的心；必须赶紧以抵抗战胜它们，免得那以暗示引诱的恶习以享乐的甘美征服人，并且，当它迟迟不被逐出心灵时，就以同意的利剑杀害人。\par

% structure: p0009 source=h2 target=subsection reason=relative-heading-rank; base=h2

\subsection{第三章}

\Emph{论治理者应常在行动中居首位。}\par

统治者应常是行动中的首领，好藉其生活为属下指明生命之道，并使那追随牧者声音与习性的羊群，更藉榜样而非言辞学会行走。因为那因职分所需而必须宣讲至高之事的人，也因同一需要而必须显示至高之事。那被说话者生活所称赞的声音，更易穿透听者的心，因他藉说话所命令的，藉显示而助其实行。为此，先知说：\BibleQuote{你这给\Place{熙雍}报喜讯的，请登上高山！}\BibleRef{依 40:9}。意思是，从事天上宣讲的人，应已弃绝世俗工作的低水平，如同立于万事之巅，并因他生活之功绩从高处呼喊，而更易将属下引向更善之事。为此，在神圣法律下，司铎接受祭献的肩，且是右肩和分离的\BibleRef{出 29:22}；以表明他的行动不仅应有益，且应是卓越的；他不仅应在恶人中行正义，还应以其德行之卓越超越属下中的善行者，如同他在职秩尊严上超越他们。胸与肩一并分给他食用，好使他学习将属于祭献的那部分自己献给万有施与者；使他不仅以胸怀持守正念，且以工作之肩邀请观看者趋向高处；使他既不贪求现世之顺境，也不畏惧逆境；使他顾及内心的敬畏而轻看世界的诱惑，但思及内在甘美而轻看其恐怖。因此，依上主之命\BibleRef{出 29:5}，司铎以\OriginalTerm{厄弗得}{ephod}之袍束肩，为使他时时以德行装饰护佑，既不受顺境也不受逆境所伤；如此，如\Person{圣保禄}所说\BibleRef{格后 6:7}，\BibleQuote{在左右手持有正义的武器}，他只追求\BibleQuote{那前面的}，而不会偏向任何一边陷入低俗之乐。愿顺境不使他得意，逆境不使他困扰；愿顺利之事不诱使他屈服意志，艰难之事不迫他陷入绝望；如此，在他谦卑的心志不被任何情欲所屈服时，他就能显示厄弗得如何以优美覆盖双肩。这厄弗得被命令用金、蓝、紫、两次染的朱红及捻的细麻制成\BibleRef{出 28:8}，以表明司铎应如何以诸多德行而卓越。因此，在司铎的袍服上，首先闪耀着金，表明他应首先在智慧的理解上发光。与之相连的是蓝，闪亮如天空之色，表明他应凭借理解洞悉一切，超越世间的恩惠，爱慕天上之事；以免因不慎被自己的赞美所困，而失却对真理的理解。与金和蓝一起，还掺入了紫：这意味着，司铎的心在盼望他所宣讲的高超之事时，应在自身内抑制邪恶的念头，并凭借王权般的力量予以驳斥，因他常顾念内心重生的高贵，并以他的品行守护天国袍服的权利。因为这精神的高贵，经\Person{伯多禄}所说：\BibleQuote{你们是特选的种族，王家的司祭}\BibleRef{伯前 2:9}。关于这使我们制服邪淫的能力，我们被\Person{若望}的话所坚固：\BibleQuote{凡接受他的，他赐他们权能成为天主的子女}\BibleRef{若 1:12}。圣咏作者在提及这德能的尊严时说：\BibleQuote{在我眼中，你的朋友何其尊贵，天主！他们的元首何其强盛！}\BibleRef{咏 137:17}。圣者的心确实被提升至王侯的卓越，虽然外表看来他们忍受卑微。但与金、蓝、紫一起，还加入了两次染的朱红，以表明一切德行的卓越都应在内在审判者眼前以爱德装饰；凡在人前闪耀的，应在隐秘裁判之前以内心的爱火照亮。此外，这爱德，既然包含对天主和邻人的爱，便如同有双重的光泽。因此，凡如此渴慕造物主之美而忽略照顾邻人，或如此照顾邻人而在对天主的爱情上变得冷淡的人，无论他忽略哪一者，都不知道在装饰厄弗得时拥有两次染的朱红是什么意思。然而，当心灵专注于爱德的诫命时，无疑肉体应通过斋戒而受克苦。因此，与两次染的朱红一起，还加入了捻的细麻。因为细麻（\OriginalTerm{细麻}{byssus}）从地里生出，闪着白光；细麻所标志的，不就是身体在贞洁纯洁中闪耀的洁白吗？它也被捻成线，织入厄弗得的美观中，因为贞洁的习惯只有通过斋戒磨损肉体，才能达到完美的纯洁洁白。既然肉体克苦的功劳与其他德性共存，捻的细麻就在厄弗得多样的美中显现洁白。\par

% structure: p0012 source=h2 target=subsection reason=relative-heading-rank; base=h2

\subsection{第四章}

\Emph{统治者当在沉默中谨慎，在言语中有益。}\par

管理者在保持沉默时应谨慎，在言谈时应有益；免得他说出本应隐忍之事，或隐忍本应说出之事。因为，正如不经思考的言谈引人误入歧途，不经思考的沉默也使本可受教者留在错误之中。不谨慎的管理者常因惧怕失去人的好感，怯于自由地说出正确之事；按真理之声\BibleRef{若 10:12}，他们并非以牧者的热忱服务于守护羊群，而是如同佣工；因为他们若以沉默藏匿自己，便在狼来时逃遁。因此，主藉先知责备他们说：\Quote{\Emph{哑巴狗，不能叫唤}}\BibleRef{依 56:10}。祂又抱怨说：\Quote{\Emph{你们没有上去对抗敌人，也没有为以色列家筑墙，在上主的日子在战争中站立}}\BibleRef{则 13:5}。上去对抗敌人，就是以自由的声音为保护羊群而对抗现世权势；在上主的日子在战争中站立，就是出于正义之爱，在恶人攻击我们时抵抗他们。因为牧者若害怕说正确的话，除了以沉默转身逃遁，还能是什么呢？但若他为羊群挺身而出，便为以色列家筑起一道墙对抗敌人。同样，对罪恶的百姓说：\Quote{\Emph{你的先知为你见了虚假愚妄的事；他们没有揭露你的罪孽，以激发你悔改}}\BibleRef{哀 2:14}。因为在圣经语言中，教师有时被称为先知，因他们指明现世之短暂，显明来世之事。神圣之言判定他们见了虚假之事，因为他们惧怕责备过错，反倒以虚假的安全许诺奉承作恶者；他们全然不揭露罪人的不义，因他们抑制自己责备的声音。责备之言是揭露的钥匙，因它能指出连犯罪者本人也常未察觉的过错。因此保禄说：\Quote{\Emph{使他能藉健全的教理说服反对者}}\BibleRef{铎 1:9}。藉玛拉基亚说：\Quote{\Emph{司铎的唇舌应保存知识，众民应从他的口中寻求法律}}\BibleRef{拉 2:7}。藉依撒意亚，主劝告说：\Quote{\Emph{呼喊} \Emph{不要保留，提高你的声音如号角}}\BibleRef{依 58:1}。确实，凡接受司铎职的人，便承担了传报者的职责，当他自己呼喊时，行走在可怕审判者来临之前。因此，司铎若不知如何宣讲，这哑口无声的传报者能发出什么呼喊之声呢？为此，圣神以舌头之形像\BibleRef{宗 2:3}停留在最初牧者们身上；因为祂所充满的人，祂便立时使他善于言辞。为此，梅瑟受命，当司铎进入会幕时，他应佩有铃铛\BibleRef{出 28:33}；即他应有宣讲之声环绕，免得因沉默激怒从上注视他的审判者。因为经上记载：\Quote{\Emph{他进入圣所到上主面前，以及出来的时候，他的响声必被听见，免得他死亡}}\BibleRef{出 28:35}。司铎进入或出来时，若他的声音不被听见，他便死亡；因为他若不带宣讲之声行事，便激怒隐藏的审判者的义怒。巧妙的是，铃铛被描述为缝在他祭衣的周边上。司铎的祭衣若不是义行，还能是什么呢？正如先知所证：\Quote{\Emph{愿你的司铎披上正义}}\BibleRef{咏 130:9}？因此，铃铛缝在祭衣上，表示司铎的善行也应随同他舌头的宣讲之声，宣告生活之道。但管理者准备说话时，应牢记他应以何等谨慎的用心说话，免得他无节制地匆忙发言，使听者的心灵被错误的伤口击中，当他或许想要显得明智时，却不智地割断了合一的纽带。为此，真理说：\Quote{\Emph{你们当中有盐，彼此和睦相处}}\BibleRef{谷 9:49}。盐象征智慧之言。因此，凡努力智慧地说话的人，应极为小心，免得因他的口才扰乱了听者的合一。为此保禄说：\Quote{\Emph{不要超过所当有的智慧，而应明智地合乎冷静}}\BibleRef{罗 12:3}。因此，按天主的诫命，在司铎的祭衣上，除了铃铛，还加了石榴\BibleRef{出 28:34}。石榴象征什么，岂不是信德的合一？因为，如同在一个石榴内，许多籽粒被一层外皮保护，信德的合一也包含了圣教会中无数的子民，他们因不同的功绩而留在教会内。为避免管理者不智地匆忙发言，真理本人向祂的门徒宣告了我们刚才引用的这句话：\Quote{\Emph{你们当中有盐，彼此和睦相处}}\BibleRef{谷 9:49}。犹如祂借用司铎的祭服形象说：将石榴加在铃铛上，使你在所说的一切中，以谨慎的警醒保持信德的合一。管理者也应焦虑地提防，不仅不可说任何错误的话，而且连正确的话也不可过分和无度地说；因为所说之事的良好效果常因不慎的闲谈而消弱于听者的心中；这闲谈本身不知为听者谋益，也玷污说话者。因此，梅瑟说得很好：\Quote{\Emph{患漏症的人是不洁的}}\BibleRef{肋 15:2}。因为所听之话的品质是随后思想的种子；话由耳听入，思想便在心中生出。为此，这世界的智者称杰出的宣讲者为\OriginalTerm{撒种者}{seminiverbius}\BibleRef{宗 17:18}。因此，患漏症的人被宣布为不洁，因为沉溺于多言的人，玷污自己，那本可有序地射出、在听者心中产生正确思想之苗裔的东西，却因他不慎的闲谈播撒，不是为生育，而是为不洁。为此，保禄在劝勉他的弟子要致力于宣讲时，说：\Quote{\Emph{我在天主和耶稣基督面前，凭着他的显现和国度，郑重吩咐你：宣讲圣言，无论方便与不便，都要坚持}}\BibleRef{弟后 4:1}；在说\Emph{不便}之前，先说了\Emph{方便}，因为若不晓得把握时机，不便本身就以自身的廉价在听者心中败坏自己。\par

% structure: p0015 source=h2 target=subsection reason=relative-heading-rank; base=h2

\subsection{\Emph{第五章}}

\Emph{牧者应藉怜悯接近众人，藉默观超拔于众人之上。}\par

治理者应当在同情上与每个人邻近，在默观上超越众人，好能藉着慈爱的胸怀将他人的软弱转移到自己身上，并藉着默观的高超，在渴求不可见之事上甚至超越自己；免得他在追求崇高事物时轻视邻人的软弱，或在迁就邻人的软弱时放弃对崇高事物的渴求。因此，保禄被提到天堂\BibleRef{格后 12:3}，探索第三层天的奥秘，然而，当他被高举于那不可见之事的默观中时，仍将心灵的视线转向肉体的床榻，指示他们如何在隐秘的私处彼此相处，说：\Quote{但为了避免淫乱，男人当各有自己的妻子，女人当各有自己的丈夫。丈夫对妻子当尽己之责，妻子对丈夫亦当如此}\BibleRef{格前 7:2}。稍后又说：\Quote{你们彼此不可亏负，除非两厢情愿，暂时分房，为专务祈祷；以后仍要同房，免得撒殚因你们不能节制而诱惑你们}\BibleRef{格前 5:5}。看，他已参透天国的奥秘，却藉着俯就的胸怀探求肉体的床榻；同一心灵之眼，在提升时仰望不可见之事，在怜悯时俯察软弱者的秘密。他在默观中超越诸天，却在焦虑的关怀中不离开肉体的卧榻；因为他藉着爱德之链同时连接至高与至低者，虽在自身中被圣神的能力强力提升至高处，但在他人中，因着他的慈爱，他甘愿成为软弱的。因此他说：\Quote{有谁软弱，我不软弱呢？有谁跌倒，我不心焦呢？}\BibleRef{格后 11:29}。又说：\Quote{对犹太人，我就成为犹太人}\BibleRef{格前 9:20}。他如此行，并非放弃信仰，而是扩展慈爱；藉着将不信者的身份比喻性地转移到自己身上，从自身学习他们应如何受到怜悯；为能将他自己若处于他们地位时所合理希望得到的，施予他们。因此又说：\Quote{如果我们癫狂，那是为了天主；如果我们清醒，那是为了你们}\BibleRef{格后 5:13}。因为他知道如何在默观中超越自己，并在俯就中适应听众。因此，雅各伯看见，主从高处俯视，油倾注在石头上，有天使上去下来\BibleRef{创 28:12}；这表示真正的宣讲者不仅因默观而上升到教会的圣首，即主那里，也因怜悯而下达到祂的肢体。因此，梅瑟频繁出入会幕，他在会幕内被包裹于默观之中，在外则忙于处理软弱者的事务。在内，他思考天主的奥秘；在外，他担负肉体的重担。关于疑难之事，他常回到会幕，在约柜前求问上主；这无疑给治理者树立了榜样：当他们在外面世界不确定如何安排事物时，应返回自己的灵魂，如同返回会幕，在约柜前求问上主，以便在他们内心查考圣经中关于所疑惑之事的圣言。因此，真理本身藉着摄取我们的人性向我们显明，祂在山上祈祷，却在城中行奇迹\BibleRef{路 6:12}，从而为善牧指明了当行之路：即使已在默观中追求至高之事，仍应以同情与软弱者的需要相融合；因为爱德正是在怜悯地趋向邻人的卑微时，奇妙地升至崇高；它愈慈爱地降至这世上的软弱，便愈有力地重返高处。但那些治理他人者，应使自己成为属下甚至敢于向他们吐露秘密的人；好让那些被诱惑之浪骚扰的小子们，如同投奔母亲的胸怀般投奔牧者的心，并在他的劝慰和祈祷的泪水中洗去因打击他们的罪恶污秽而预见到的玷污。因此，圣殿门前，为进入者洗手所用的铜海，即洗濯盆，由十二只铜牛支撑\BibleRef{列上 7:23}，牛的脸面显露出来，而后部却被隐藏。这十二只牛所象征的，不就是整个牧者团体吗？关于这团体，法律经保禄解释说：\Quote{牛在场上踹谷的时候，不可笼住它的嘴}\BibleRef{格前 9:9}；\BibleRef{申 25:4}。他们的公开工作我们看得见，但他们在严厉审判者隐藏的报应中背后所存留的，我们却不知道。然而，当他们准备俯就的耐心，在告解中洗净邻人的罪过时，他们就如支撑着圣殿门前的洗濯盆；使任何力求进入永生之门的人，都能将他的诱惑展示给牧者的心，并在牛的洗濯盆中洗净他思想和行为的双手。而且往往发生这样的情况：当治理者的心灵藉着俯就得知他人的考验时，自己也被所听见的诱惑所攻击；因为同一个洗濯盆的水，既洗涤众多百姓，无疑也使自己受到玷污。在接纳洗涤者的污秽时，它仿佛失去了自身纯净的宁静。但牧者绝不应为此惧怕，因为在天主精妙平衡万物的眷顾下，他愈是怜悯地为他人的诱惑而苦恼，就愈容易从自己的诱惑中被解救出来。\par

% structure: p0018 source=h2 target=subsection reason=relative-heading-rank; base=h2

\subsection{第六章}

\Emph{治理者应当借着谦卑，成为善生者的同伴，但借着正义的热忱，对恶人的恶行保持严厉。}\par

治理者应当借着谦卑，成为善生者的同伴，并借着正义的热忱，对恶人的恶行保持严厉；如此，在任何事上，他都不应将自己置于善人之上，但当恶人的过犯需要时，他应立即意识到自己优先地位的权能；这样，在他那些生活良好的下属中间，他放弃自己的等级，视他们为平等，而对于悖逆者，他并不畏惧执行正义的法律。因为，正如我曾在 \WorkTitle{伦理} 书中所述（卷二十一，第10章，第22节），显然，自然使所有人平等；但因功绩次序的变化，罪过使一些人低于另一些人。而这由邪恶产生的差异，却由天主的审判所安排，既然并非所有人都能处于平等地位，一人就应受另一人管辖。因此，所有管辖他人者，应当在自己身上考虑的不是其地位的权威，而是其状况的平等，并且要乐于不是管辖人，而是造福于人。的确，我们的古代先祖被称为不是人的王，而是羊群的牧者。当主对 \Person{诺厄} 和他的儿子们说：\BibleQuote{你们要滋生繁衍，充满大地}（\BibleRef{创 9:1}），祂立即加上：\BibleQuote{地上的一切走兽都要怕你们，畏惧你们。}由此可见，既然命令恐惧和畏惧应加于地上走兽，就禁止加于人类。因为人自然优于禽兽，但不优于其他人；因此，对他说他应被走兽畏惧，而非被人畏惧；因为想被自己的同等人畏惧，就是违反自然的骄傲。然而，当统治者发现他们的臣民不畏惧天主时，他们有必要被臣民畏惧；这样，那些不畏惧天主审判的人，至少会因对人的畏惧而不敢犯罪。因为上级在寻求引起这种畏惧时，绝非表现骄傲，他们在其中寻求的不是自己的荣耀，而是下属的正义。因为当他们向那些生活悖逆的人要求对自己的畏惧时，他们仿佛不是统治人，而是统治走兽；因为，就他们的下属兽性而言，他们也应当屈服于畏惧。\par

但通常，治理者因超越他人，便为思想的得意所膨胀；当一切事物都服务于他的需要，他的命令迅速按他的意愿执行，所有下属都赞美他所行善事，却无权指责他所行不当，并且他们常常赞美本应责备的事，他的心灵被来自下面的大量奉承所引诱，便高抬自己，超越自己；而当他被无限的好感所包围时，便失去了内在的真理感；他忘记自己，将自己散布在他人的言论中，相信自己是外在听到的那样，而非他内在应当判断的那样。他轻视那些在他之下的人，不承认他们在自然秩序中与他平等；他在权势的偶然性上超过的人，他相信自己也在生命的功绩上超越了；他自以为比所有他见在权力上超越的人都聪明。的确，他在自己的心灵中自设某种崇高的高处，虽然与他人处于相同的自然状况，却不屑于从同一水平看待他人；于是他变得甚至像那经上所写的：\BibleQuote{他注视一切崇高的事物：他是所有骄傲之子的君王}（\BibleRef{约 41:25}）。不仅如此，他渴望独一的高位，鄙视天使的共融生活，说：\BibleQuote{我要将我的宝座安置在北方，我要与至高者相似}（\BibleRef{依 14:13}）。因此，借着奇妙的审判，他在自己内部发现跌落的深坑，而外在却在权力的顶峰上高抬自己。因为他确实变得像那背叛的天使，当他作为人却不屑于像人时。就这样，\Person{撒乌耳}在谦卑的功绩之后，在权力的高峰上变得骄傲自大：他因谦卑而被拣选，因骄傲而被弃绝；正如主所证实的，祂说：\BibleQuote{当你自视微末时，我不是立你为\Place{以色列}各支派的首领吗？}（\BibleRef{撒上 15:17}）\Person{撒乌耳}先前在自己眼中是微末的，但当他被现世权力支撑时，便不再自视为微末。因为他在与他人比较时，因拥有比所有人都大的权力而偏爱自己，便自以为超越众人。然而奇妙的，当他在自己眼中微末时，他在天主面前是伟大的；但当他在自己眼中显得伟大时，他在天主面前却是微末的。因此，通常当心灵因下属的丰裕而膨胀时，便腐败为骄傲的流溢，那权力的顶峰反而成了欲望的帮凶。事实上，谁能既保持权力又与之作战，谁就能妥善地管理这权力。谁能知道借它超越过犯，并知道以它平等地与他人相称，谁就能妥善地管理它。因为人心通常在毫无权柄支持时也会被高举；何况当权柄助长时，它岂不更加高抬自己！然而，唯有那知道谨慎地既从中获取有益之物，又拒绝诱惑，并借着它感到与别人平等，同时在惩罚罪人的热忱中将自己置于他们之上的人，才能妥善地运用这权柄。\par

但我们若要更充分地理解这一区别，就应察看第一位牧者的榜样。因为\Person{伯多禄}从天主接受了\Person{圣教会}的首席地位，当\Person{科尔乃略}善行而谦卑地俯伏在他面前时，他拒绝接受过度的敬拜，说：\BibleQuote{\Emph{你起来，不要这样；我自己也是个人}}\BibleRef{宗10:26}。但当他发现\Person{阿纳尼雅}和\Person{撒斐辣}的罪过时，他立刻显出他以何等大的权能被高举于众人之上。因为他以言语击杀了他们的性命，这是他藉着精神的洞察所察觉的；他忆起自己是在教会内抵挡罪恶的首领，尽管在弟兄善行而向他热切致敬时，他并不承认这一点。在一方面，圣洁的行为赢得了平等的共融；在另一方面，报应的热忱显明了权威的公正要求。\Person{保禄}也并不知道他自己高于善行的弟兄，当他说：\BibleQuote{\Emph{并不是因为我们管辖你们的信德，而是帮助你们的喜乐}}\BibleRef{格后1:23}。他立即加上：\BibleQuote{\Emph{因为你们凭信德站立}}，似乎要解释他的宣告说：因此我们没有管辖你们的信德，因为你们凭信德站立；因为我们与你们平等，在于我们知道你们站立之处。他并不知道自己高于弟兄，当他说：\BibleQuote{\Emph{我们在你们中间成了婴孩}}\BibleRef{得前2:7}；又说：\BibleQuote{\Emph{但我们是你们的仆人，因着基督}}\BibleRef{格后4:5}。但当他发现一个需要纠正的过错时，他立即忆起自己是师傅，说：\BibleQuote{\Emph{你们愿意怎样？要我拿着棍棒到你们那里去吗？}}\BibleRef{格前4:21}。\par

因此，当治理者以权位统御恶习而非统御弟兄时，至高权柄便有序得当。但是，当上司纠正下属的过失时，他们必须谨慎留意：在运用权威以适当纪律打击过犯的同时，如何因持守谦卑而承认自己与被纠正的弟兄是平等的；虽然在大多数情况下，我们静默自省时甚至应当视所纠正的弟兄优于自己。因为他们的恶习通过我们以纪律的力度被击打；而我们自己所犯的过犯，却连一句斥责都未受刺痛。因此，我们在人前犯罪不受惩罚，在主前所受的约束便越大；而我们的纪律使下属在此世不受过犯报应，从而更免于天主的审判。故此，心中应持守谦卑，行事应持守纪律。并且始终需要明智的洞察，以免因过度持守谦卑之德而松懈治理的正当要求，以免上司过于自卑而无法以纪律的约束管束下属的生活。因此，治理者应在外表持守他们为他人利益所承担的事，在内里持守那使他们自我评估时感到敬畏的事。然而，甚至他们的臣属也当通过合宜显现的某些迹象，感知他们内心是谦卑的；以致既在权威中看见可敬畏之处，又承认在谦卑方面有可效法之处。因此，凡任职治理者须不断学习，使外在看来强大的权能，在内里被抑制；使它不征服他们的思想；使心灵不被它的喜乐所夺；免得理智无法控制那在统治欲中自身所屈从的对象。因为，为防止统治者的心因个人权能的喜乐而陷入高傲，某位智慧人说得正对：\BibleQuote{他们立了你作领袖：不要心高，却要待他们如他们中的一员} \BibleRef{德 32:1}。因此伯多禄也说：\BibleQuote{不是做天主所托付给你们者的主宰，而是做群羊的模范} \BibleRef{伯前 5:3}。因此真理本身亲自激励我们追求更高的功德，说：\BibleQuote{你们知道：外邦人有首长主宰他们，有大臣管辖他们。在你们中却不可这样；谁若愿意在你们中成为大的，就当作你们的仆役；谁若愿意在你们中为首，就当作你们的奴仆；就如人子来不是受服事，而是服事人} \BibleRef{玛 20:25}。因此祂也指明那因接受治理而心高气傲的仆人所受的惩罚，说：\BibleQuote{如果那恶仆心里说：我主人必要迟延，于是开始拷打自己的同伴，且与醉汉一同吃喝。那仆人的主人要在他不期待的日子，和不察觉的时刻来到，把他斩绝，使他与假善人同受惩罚} \BibleRef{玛 24:48}。因为那以纪律为借口将治理职分转为统治目的的人，被正当地列入假善人之中。然而有时，若在悖逆之人中间平等超过纪律，则是更严重的失职。因为厄里，因被虚假情感胜过，不愿惩罚他犯罪的儿子，便在严厉的审判者面前，与儿子一同受到残酷的惩罚 \BibleRef{撒上 4:17}。因此，天主的声音对他说：\BibleQuote{你尊重你的儿子过于尊重我} \BibleRef{撒上 2:29}。同样，祂通过先知斥责牧者说：\BibleQuote{你们没有包扎那折断的，也没有领回那迷失的} \BibleRef{则 34:4}。因为凡跌倒于罪恶的人，因牧者关怀的力度而被召回正义状态时，那被抛弃的就被领回。因为当纪律制服罪恶时，束缚绑住骨折，免得因缺乏约束严厉的压挤而致命地流血。但骨折若被不慎地绑缚，往往变得更糟，因受韧带过度约束而使伤口更剧烈疼痛。因此，当用惩戒压制下属的罪伤时，甚至约束本身也须极为谨慎地节制，以便对失职者行使纪律权利时仍保留慈爱的情怀。因为当谨慎的是：治理者应显示自己如慈母般的仁爱，又如严父般的纪律。并且始终要以焦虑的谨慎注意，使纪律不严苛，仁爱不松弛。因为，正如我们在\WorkTitle{伦理论丛}中曾说过的（\Emph{Lib.} xx., \Emph{Moral} n. 14, c. 8, \Emph{et ep.} 25, \Emph{lib.} 1），若仅得其一而缺另一方，则纪律与怜悯皆多有不足。但在治理者对其臣属身上，应当既有公正周到的怜悯，又有慈爱严厉的纪律。因此，正如真理所教导的（\BibleRef{路 10:34}），那人被撒玛黎雅人的关怀带到客店，半死不活，并且将酒和油倒在他的伤口上；酒使伤口刺痛，油使伤口舒缓。因为凡负责医治伤口的，必须用酒施以疼痛的刺痛，用油施以仁爱的柔润，以便通过酒净化腐烂之处，通过油安抚可治愈之处。因此，温良应与严厉相混合；二者应合成一种复合剂，使臣属既不会因过于粗暴而破损，也不会因过于仁慈而松弛。正如保禄的话所指（\BibleRef{希 9:4}），这由会幕中的约柜很好地预表，其中除了约版外，还有杖和玛纳；因为，如果善牧的胸中有对圣经的认识，有约束的杖，也应有甘甜的玛纳。因此达味说：\BibleQuote{你的牧杖和短棒，给我安慰} \BibleRef{咏 23:4}。因为，若有击打的杖来约束，也应有扶持的棍来安慰。因此，当有爱，但不使人衰弱；当有刚毅，但不激怒；当有热诚，但不无节制地燃烧；当有怜悯，但不过分宽宥；如此，正义与怜悯在至高权柄中彼此融合，为首者既可在使人畏惧中抚慰臣属的心，又可在抚慰中约束他们达到敬畏。\par

% structure: p0024 source=h2 target=subsection reason=relative-heading-rank; base=h2

\subsection{第七章}

\Emph{治权者不得在忙于外在事务时放松对内务的照顾，也不得在挂虑内务时忽略对外务的供应。}\par

治权者不应在忙于外在事务时放松对内务的照顾，也不应在挂虑内务时忽略对外务的供应；免得他因沉溺于外在事物而远离内心关切，或因只专注于内务而未能将所欠给与身外的近人。因为常有这样的事：有些人仿佛忘记了自己被置于弟兄之上是为了他们的灵魂，竟全心致力于世俗事务；这些事务，当临近时，他们乐于处理，即便缺乏，也日夜以混乱思绪的沸腾去渴求；而且，当偶然因没有机会而得以安静时，他们反而因这安静而更加疲惫。因为他们以行动疲累为乐；他们以不在世俗事务中劳作为苦。因此，当他们以被世俗喧嚣驱赶为乐时，便对内在的事一无所知，而正是这些事他们本应教导他人。而由此，他们的属民的生活也无疑变得麻木；因为，当渴望在灵性上进步时，却因上级的榜样而在路上遇到绊脚石。因为当首领衰弱时，肢体就不能兴盛；如果行军的领导者行差踏错，军队追随敌人也是徒劳。没有任何劝勉能支撑属民的心，没有任何责备能惩罚他们的过失，因为，当灵魂的监护人执行现世法官的职务时，牧者的注意力就从看管羊群上转移了；而属民无法领悟真理之光，因为，当世俗事务占据了牧者的心思时，被诱惑之风吹起的尘土就蒙蔽了教会的眼睛。为防范这一点，人类的救赎主在阻止我们贪饕时说：\BibleQuote{\Emph{你们要谨慎，免得你们的心因贪食醉酒而沉重}}\BibleRef{路 21:34}，随即加上：\BibleQuote{\Emph{或人生的挂虑}}\BibleRef{路 21:34}，在同一处，为增加警告的恐惧，祂立刻说：\BibleQuote{\Emph{免得那日子忽然临到你们}}\BibleRef{路 21:34}；祂甚至还宣告那来临的方式，说：\BibleQuote{\Emph{因为那日子如同罗网，要临到全地面的一切居民}}\BibleRef{路 21:35}。因此祂又说：\BibleQuote{\Emph{没有人能事奉两个主人}}\BibleRef{路 16:13}。因此保禄以召唤将虔敬者的心灵从与世俗结伴中带离，甚至说是征募他们，他说：\BibleQuote{\Emph{没有人作战为天主效劳，而让自己纠缠于今世的事务，为能取悦于他所侍奉的主人}}\BibleRef{弟后 2:4}。因此，他向教会的治理者既推荐闲逸的研习，又指出劝告的补救方法，说：\BibleQuote{\Emph{如果你们有世俗的诉讼，就让那些在教会中可轻视的人来审判}}\BibleRef{格前 6:4}；这就是说，那些没有灵性恩赐装饰的人，应当投身于现世的职务。仿佛他更明白地说：既然他们不能参透内在事物，至少让他们在外在必要之事上操劳。因此，与天主说话（\BibleRef{出 18:17}）的\Person{梅瑟}，竟被外族人\Person{耶特洛}的斥责所判定，因为他以不明智的辛劳投身于百姓的世俗事务：并且随即也给他建议，要他另立他人代他解决世俗纷争，而他自己则更自由地学习属灵奥秘以教导百姓。\par

因此，属下的臣民应处理次要事务，而统治者则当思虑最高之事；如此，那被安置在高处以眺望前路的眼目，才不会被尘埃的烦扰所遮蔽。因为凡为首者，乃是其属下的头；为使双脚能行正道，头无疑应从高处前瞻，免得身体弯曲、头垂向地，双足便在路上踟蹰不前。然而，灵魂的监督者若自己沉溺于他本应谴责他人的俗世挂虑，又怎能以牧者的尊严在众人中自处呢？这诚然是主在义怒的公正报应中，通过先知所警告的，说：\BibleQuote{百姓如何，司铎也如何}\BibleRef{欧 4:9}。因为司铎与百姓无异，系指身负神职者，其行为与仍因肉欲追求而受审判的他人相同。这正是耶肋米亚先知以极大的爱德之痛，借圣殿被毁的景象所哀叹的，说：\BibleQuote{黄金何其失光！纯色何其改变！圣所的石头散落在各街头上。}\BibleRef{哀 4:1}。黄金超越万有，所表达的岂非圣德的高超？纯色所表达的，岂非围绕宗教、为人人所爱的敬畏？圣所的石头所指的，岂非圣职人员的身份？街道之名所寓意的，岂非今世生活的广阔？因为，在希腊语中，表示宽阔的词是\OriginalTerm{宽度}{\GreekText{πλάτος}}，街道（\OriginalTerm{街道}{plateae}）便因其宽度或辽阔而得名。然而真理本身说：\BibleQuote{宽大而广阔的路通向毁灭。}\BibleRef{玛 7:13}。因此，当圣洁的生活被世俗作为玷污时，黄金便失光；当那些曾被认为虔敬生活的人的美名减损时，纯色便改变。因为，当一个人习于圣洁之后，若又与世俗作为混杂，他的颜色便仿佛改变了，环绕他的敬畏在世人眼中也变得黯淡而不被尊重。圣所的石头也被抛在街上，当那些本应为教会的装饰、自由深入内在奥秘（犹如在帐幕的隐密处）的人，却在外面寻求世俗讼案的广阔道路。的确，他们成为圣所的石头，原是为在至圣所内、大司祭的衣袍上显现。但当宗教的仆役不能以他们生活的功德，从属下的人那里获得救赎者的尊荣时，他们便不是在大司祭的装饰中的圣所石头。确实，这些圣所的石头散落在街上，当身居圣职的人沉溺于自己欲望的广阔，紧紧抓住世俗的事务时。还应注意，经文说他们被分散，不是在大街上，而是在街头上；因为，即使他们从事世俗事务，也渴望显得最高；以便占据广阔的途径以享受欢愉，同时却因圣德的尊贵而居于街头上。\par

此外，我们也不妨将圣所的石块视作建造圣所本身的石头；它们散落于街头巷尾，当那些身居圣品、其职位曾似乎代表圣德荣耀的人，却偏爱上世俗事务之时。因此，世俗的事务，虽有时可因慈悲而予以容忍，但绝不应因对事物本身的喜爱而追求之；以免它们压垮喜爱之人的心思，使之心为自身重担所胜，从天上的高处沉沦至最低处。但是，另有一些人承担了牧养羊群的职责，却希望有暇专注于自身灵性事务，以致全然不涉足外在事务。这种人因疏忽所有关涉身体的事，完全不满足托付于他们的人的需要。他们的讲道自然多半被轻视；因为，他们虽指责罪人的行为，却不供给他们现世生活的必需品，人们便不愿听他们。因为教义之言无法穿透困乏之人的心灵，若慈悲之手不将其推荐到他心中。但圣言的种子易于发芽，当讲道者的仁爱滋润听者的胸怀时。因此，为使管理者能注入内在有益之事，他必须以无可指责的思虑，也提供外在事物。那么，让牧者在关怀所属之人的内心情感上如此热心，也不放弃供应他们外在的生活。因为，如我们所说，羊群的心可以说是自然地抗拒讲道，如果牧者忽略了外在援助的关怀。因此，首位牧者也忧心忡忡地劝诫说：\BibleQuote{我劝你们中间的长老们，我这同作长老的，也是基督受苦的见证人，和那将要显现的荣耀的分享者，你们要牧养你们中间天主的羊群}\BibleRef{伯前 5:1}。在此处他表明他所赞扬的是牧养心灵还是牧养身体，因为他立即补充说：\BibleQuote{供应它，不是出于勉强，而是出于情愿，按照天主，不是为不义之财，而是出于热忱}\BibleRef{伯前 5:1}。在这些话中，牧者确实被善意地预先警告，免得他们在满足属下匮乏的同时，以野心的剑自戕；免得他们借使邻人通过肉身的援助得到更新，自己却仍禁食于正义之粮。保禄激起牧者的这种关怀，他说：\BibleQuote{谁不照顾自己的亲属，尤其自己的家人，就是背弃了信德，比不信的人更坏}\BibleRef{弟前 5:8}。因此，在这所有之中，他们应当惧怕，并警醒注意，免得在忙于外在关怀时，被淹没了内在的专注。因为通常，如我们已说过的，管理者的心，在轻率地投身于暂时忧虑时，内在的爱会冷却；并且，在外四处奔忙，竟不怕忘记自己承担了照管灵魂的职责。那么，有必要将用于属下身上的关怀保持在一定限度之内。因此，对厄则克耳说得好：\BibleQuote{司铎不可剃头，也不可让头发长长，只可剪发}\BibleRef{则 44:20}。因为那些被置于信友之上给予神圣指导的人，才恰当称为司铎。但头发外在头上的是心中的思想；这些思想如不知不觉从脑际生发，象征现世生活的忧虑；由于忽略察觉，有时不合时宜地出现，几乎在我们未察觉时前进。既然所有管理他人的人确实应当有外在的焦虑，却又不应过度执着于它们，因此正确禁止司铎剃头或让头发长长；这样他们既不完全割舍为属下生活所需的肉身思想，也不让这些思想过度增长。因此在这段经文中说得好：\BibleQuote{只可剪发}\BibleRef{则 44:20}；意思是，暂时忧虑的关怀应根据需要延伸，又要很快剪短，以免生长到无节制的程度。因此，当通过对外在身体的谨慎关怀，生命得以保护（或：通过外在应用的谨慎关怀，身体的生命得以保护），同时通过适度的内心专注而不受妨碍时，司铎头上的头发既保留以覆盖皮肤，又剪短以免遮住眼睛。\par

% structure: p0029 source=h2 target=subsection reason=relative-heading-rank; base=h2

\subsection{第八章}

\Emph{管理者不应专心取悦人，但仍应留意什么当取悦他们。}\par

与此同时，牧者也必须保持警惕，免得取悦人的私欲攻击他；免得当他专心参透内心之事，并远见地供应外在之事时，他寻求被属下爱戴胜过真理；免得当他因善行支持而看似不属于世界时，私爱使他远离造物主。因为那借着所行的善工贪求被教会爱戴以取代基督的人，是救赎者的仇敌；因为新郎派仆人带着礼物送给新娘，若仆人企图取悦新娘的眼目，便是怀有叛逆的心思。事实上，这种私爱一旦占据牧者的心，有时会过分地将其引向柔软，但有时引向粗暴。因为出于私爱，牧者的心倾向于柔软：当他看到属下犯罪时，他不敢责备他们，唯恐他们对他的爱变得迟钝；甚至有时他用谄媚缓和属下的过失，而他本应斥责这些过失。因此先知说得很好：\BibleQuote{\Emph{祸哉，那些在各肘下缝制靠垫，在各人头下安放枕头，以捕获人灵魂的人}}\BibleRef{则 13:18}；因为给各肘下缝制靠垫，就是用谄媚的奉承来呵护那些从正直中跌倒、躺卧在今世享乐中的灵魂。这就像躺着的人手肘靠在靠垫上，头靠在枕头上，当责备的严厉从犯罪者身上撤去，恩惠的温柔被提供给他时，他可以安然卧在错谬中，而没有任何反对的粗鲁打扰他。但如此自爱的牧者无疑向那些他们担心可能损害他们追求现世荣耀的人显明自己。而那些他们看见无力对抗他们的人，他们总是用严厉责备的粗暴压制他们，从不温和地劝告他们，反而忘记牧者的仁慈，以统治的权力恐吓他们。天主的声音通过先知正义地斥责他们说：\BibleQuote{\Emph{但你们以严厉和权势统治了他们}}\BibleRef{则 24:4}。因为他们爱自己胜过爱造物主，便向属下傲慢地自高，不考虑他们应当做什么，只考虑他们能做什么；他们不惧怕未来的审判；他们狂妄地以现世权力为荣；他们喜欢可以自由地做甚至非法的事，并且没有属下敢反对他们。因此，那决心行恶却希望所有其他人对此保持沉默的人，他自己就是见证，证明他渴望自己被爱胜过真理，因为他不想真理被用来反对自己。确实，没有人能活得完全不负责任。因此，那希望真理被爱超过自己的人，是希望没有人因真理而饶恕他。因此，\Person{伯多禄}甘愿接受\Person{保禄}的责备\BibleRef{迦 2:11}；因此，\Person{达味}谦卑地听从臣子的责斥\BibleRef{撒下 12:7}；因为好的牧者，自己不觉察怀有偏私的爱，相信来自属下的坦诚之言是谦卑的敬意。但与此同时，治理的关怀必须用如此高超的管理技巧来调和，使属下的心智在能够正确感受某些事情时，进步到言论自由，但自由不至于爆发成骄傲；免得言语的自由或许过多地给予他们，失去生活的谦逊。也应当记住，好牧者渴望取悦人是正确的；但这是为了通过自己品格的甜美将邻舍引向对真理的爱慕；而不是他们渴望自己被爱，而是要把对自己的爱慕当作一条道路，借此将听者的心引向对造物主的爱。因为一个不被爱的宣讲者，无论他宣讲得多好，确实很难让人愿意聆听。因此，那管理他人的人应当努力被爱，以便被人聆听，但仍不寻求爱本身，免得他在思想隐秘的篡夺中被发现背叛那位他在职务上所服侍的主。\Person{保禄}很好地暗示了这一点，当他显示出对我们之爱的秘密时说：\BibleQuote{\Emph{如同我在一切事上取悦众人}}\BibleRef{格前 10:33}。然而他又说：\BibleQuote{\Emph{如果我仍取悦于人，我就不是基督的仆人了}}\BibleRef{迦 1:10}。因此\Person{保禄}取悦，又不取悦；因为他渴望取悦，不是寻求他自己取悦人，而是真理通过他取悦人。\par

% structure: p0032 source=h2 target=subsection reason=relative-heading-rank; base=h2

\subsection{第九章}

\Emph{牧者应当谨慎了解恶习如何常以美德自居。}\par

牧者也应当了解恶习如何常以美德自居。因为吝啬常以节俭之名遮掩自己，而挥霍则以慷慨之名隐藏。过分松懈常被当作仁爱，无节制的愤怒被视为神火之德。急躁冒进常被当作敏捷之效，拖延被视为稳重之思。因此，灵魂的牧者必须以警觉的细心辨别善恶，免得吝啬占据他的心，而他因在花费上看似节俭而沾沾自喜；或者，当任意挥霍时，他自夸好像富有同情心的慷慨；或者，在宽免他本应惩治之事时，他将属下引向永罚；或者，在无情击打罪过时，他自己更严重地冒犯；或者，因过早行动而破坏了本可以妥善庄严完成之事；或者，因拖延善行的功德而使其变为更糟。\par

% structure: p0035 source=h2 target=subsection reason=relative-heading-rank; base=h2

\subsection{第十章}

\Emph{牧者在纠正与纵容、热忱与温和之间应当如何辨别。}\par

也应知道，属下的恶行有时应审慎地视而不见，但通过这种视而不见而予以暗示；有时应适时容忍那些公开知晓的事，但有时虽隐藏却应严密调查；有时应温和地责备，但有时应严厉指责。因为，正如我们所说，有些事应审慎地视而不见，但通过这种视而不见而予以暗示，这样，当犯错者意识到自己已被发现并被容忍时，他会羞于增加那些他自己以为被沉默容忍的过错，并会在自己的审判中惩罚自己，因为统治者的耐心仁慈地原谅了他。主借着先知很好地责备了犹大，当他说：\BibleQuote{你撒谎了，也不记念我，也不把我放在心上，因为我缄默不作声，好像没有看见的人。}（\BibleRef{依 57:11}）于是，他既视而不见又使之知晓，因为他既对罪人沉默，又宣布了这沉默的事。但有些事，即使公开知晓，也应适时容忍；也就是说，当环境没有提供适当机会进行公开纠正时。因为伤口如果不合时宜地切开，会更加红肿，如果药物不合时机，无疑会失去疗效。但是，在寻找纠正属下的合适时机时，长上的耐心正是在他们罪过的重压之下得到锻炼。因此圣咏作者说得很好：\BibleQuote{罪人加筑在我背上。}（\BibleRef{咏 127:3}）因为我们在背上承受重担；因此他抱怨罪人加筑在他背上，好像在说：那些我无法纠正的人，我像担子一样背在身上。\par

有些隐藏的事，却应严密调查，这样，通过某些迹象的显露，长上可以发现属下心中深藏的一切，并通过及时的责备，从很小的事情察觉到更大的事。因此，对\Person{厄则克耳}说得很好：\BibleQuote{人子，你要挖墙。}（\BibleRef{则 8:8}）该先知随即补充说：\BibleQuote{我挖墙以后，看，有一道门。他对我说：你进去，看他们在这里所做的可憎的恶事。我进去一看，见墙上画着各种爬行之物、可憎的走兽，以及以色列家的偶像。}（\BibleRef{则 8:9-10}）\Person{厄则克耳}代表权威人士；墙象征属下的硬心。挖墙是什么意思？不就是通过尖锐的质问打开心的硬壳吗？当挖了这墙，就出现一道门，因为当硬心被仔细询问或适时责备所刺透时，就仿佛出现了一道门，通过它可以看到被责备者内心的思想。因此紧接着说：\BibleQuote{你进去，看他们在这里所做的可憎的恶事。}（\BibleRef{则 8:9}）他好像进去看可憎之事，就是通过某些外表迹象的审查，如此深入属下的心，以至于认识到他们所有非法的思想。因此他又说：\BibleQuote{我进去一看，见墙上画着各种爬行之物、可憎的走兽。}（\BibleRef{则 8:10}）爬行之物象征完全属地的思想；走兽则象征那些确实稍微高出地面，但仍渴望尘世回报的奖赏。因为爬行之物全身紧贴地面；而走兽身体大部分高出地面，却因贪婪而常常俯向地面。因此，墙内有爬行之物，当心中盘旋的思想从未超越尘世的欲望；墙内也有走兽，当虽有一些正义和可敬的思想，但仍从属于对暂时利益和荣誉的贪恋，虽本身可说是高出地面，但因渴望讨好，如同喉咙的贪欲，而卑躬屈膝于最底层。因此又加上：\BibleQuote{以色列家的偶像都画在墙上。}（\BibleRef{则 8:10}）因为经上记载：\BibleQuote{贪婪就是偶像崇拜。}（\BibleRef{哥 3:5}）因此，在走兽之后提到偶像，因为有些人虽通过可敬的行为似乎高出地面，却因不光彩的野心降低自己到地面。说\Quote{画在墙上}很恰当，因为当外在事物的表象被引入内心时，在想象图像中所默想的，就仿佛画在心上。因此应注意，首先看到墙上的洞，然后是一道门，最后隐藏的可憎之事才显露出来；因为事实上，每一个罪的迹象首先在外表看到，然后指出打开罪恶的门，最后一切隐藏的邪恶才被揭露。\par

有些事，然而，应当温和地责备：因为，当过错不是出于恶意，而仅仅是由于无知或软弱时，当然有必要以极大的节制来缓和责备。因为事实上，我们所有存活在这必死肉体内的人，都受制于我们腐朽的软弱。因此，每个人都当从自身体会，他应该如何怜悯他人的软弱，以免因过于急切地对邻人的软弱说出责备的话，而似乎忘记了自己的软弱。因此\Person{保禄}劝诫得好，当他说：\BibleQuote{弟兄们，如果见一个人陷于某种过犯，你们既是属神的人，就该用柔和的心将他挽回；但你们要小心，免得也被引诱。} \BibleRef{迦 6:1}；似乎是明说：当你对他人的软弱感到不满时，要思想你是什么样的人，这样，精神在责备的热忱中就能自我节制，因为它也惧怕自己所责备的。\par

有些事，然而，应当严厉责备：这样，当过错不被犯罪者所认识时，他能从责备者口中意识到其严重性；并且，当某人对自己所行的恶事自圆其说时，他能被责备者的刚硬引导，而对它的后果产生严重恐惧。因为事实上，治理者的职责是藉着宣讲的声音显示超性故乡的荣耀，揭示在这场人生的旅途中潜伏着古仇敌何等大的诱惑，并要以极大的热忱去纠正那些在他管辖之下、不应被温和容忍的邪恶；以免对过错愤怒太少，而他自己被定为在所有过错上有罪。因此向\Person{厄则克耳}说得好：\BibleQuote{你要取一块砖，放在自己面前，在上面刻上耶路撒冷城。} \BibleRef{则 4:1}。接着立即又加上：\BibleQuote{你要围攻它，筑堡垒，堆土堆，安营扎寨，四面架起攻城锤。} 并且，为他自己的防御，又立即加上：\BibleQuote{你要拿一个铁鏊子，将它放在你和城之间，作为铁墙。} 因为\Person{厄则克耳}先知所象征的不是别的，正是教师们，因他对他说：\BibleQuote{你要取一块砖，放在自己面前，在上面刻上耶路撒冷城？}\par

因为确实，圣善的教师取一块砖，当他们抓住听众属世的内心以便教导他们时：这块砖他们实际放在自己面前，因为他们以全副心神看守它；他们也被命令在这砖上刻上耶路撒冷城，因为他们竭尽全力藉着宣讲向属世的内心描绘超性平安的异象。但是，因为天上的光荣徒然被感知，除非也知道诡诈的仇敌在这里怎样大的诱惑袭击我们，所以适当地加上：\BibleQuote{你要围攻它，筑堡垒。} 因为确实，圣善的宣讲者围攻那块刻有耶路撒冷城的砖，当他们向一个属世却已寻求天上故乡的心灵显示，在这现世时期，有何等大的恶行对立布局攻击它。因为，当显示每一个不同的罪如何在我们前进的道路上困扰我们时，就好像藉着宣讲者的声音对耶路撒冷城进行围攻。但是，由于宣讲者不仅应当使人知道恶习如何袭击我们，而且也要知道受保护的德性如何坚固我们，所以又正确地说：\BibleQuote{你要筑堡垒。} 因为确实，圣善的宣讲者筑堡垒，当他显示什么德性抵抗什么恶习。并且，因为当德性增长时，诱惑的战争大多也增加，所以又正确地说：\BibleQuote{你要堆土堆，安营扎寨，四面架起攻城锤。} 因为，当某个宣讲者陈明日益增长的诱惑的规模时，他堆土堆。他对着耶路撒冷安营扎寨，当他向听众的正确意向预告狡猾仇敌那不可测、仿佛不可理解的伏击。他四面架起攻城锤，当他使人知道在这生命中从各方包围我们的诱惑之箭，并穿透我们的德性之墙。\par

但是，虽然治理者可能巧妙地暗示所有这些事情，他并不能为自己获得永远的赦免，除非他怀着热忱的精神对待每个人的过失。因此在那个地方又正确地加上：\BibleQuote{你要拿一个铁鏊子，将它放在你和城之间，作为铁墙。} 因为铁鏊子象征思想的煎熬，铁象征责备的刚硬。\par

但是，有什么比爱天主的热忱更猛烈地煎熬和折磨教师的心灵呢？因此\Person{保禄}正被这铁鏊子的煎熬所烧灼，当他说：\BibleQuote{有谁软弱，我不软弱呢？有谁跌倒，我不焦急呢？} \BibleRef{格后 11:29}。并且，因为凡是被天主的热忱所点燃的人，都持续受到护守，以免他因疏忽而被定罪，所以正确地说：\BibleQuote{你要将它作为铁墙放在你和城之间。} 因为铁鏊子被放在先知和城之间作为铁墙，是因为当治理者已经显示出强烈的热忱时，他们后来在这热忱与他们和听众之间保持这热忱作为坚固的防御，以免他们先前在责备上懈怠，而然后没有惩罚的能力。\par

但与此同时，应当记住：当规劝者的心灵因责备而激愤时，他很难避免说出一些不该说的话。而且经常发生的情况是，当属下的过错受到严厉谴责时，上司的舌头会流于过激的言辞。当责备过度激烈时，犯过者的心会沉入绝望。因此，激愤的统治者当他考虑到自己对属下的感情造成的伤害超过了应有的程度时，他必须在内心深处借助忏悔，以便在真理面前通过哀哭获得宽恕；即便是因为他因对真理的热忱而犯罪。这便是主通过梅瑟以比喻所吩咐的，说：\BibleQuote{若有人心地单纯，与朋友一同到林中去砍柴，斧头从手中飞出，铁刃从柄上脱落，击中朋友，致他死亡，这人应逃到上述的其中一座城，得以生存；以免那被杀的至亲，心中火热，追赶他，赶上他，将他击杀}（\BibleRef{申 19:4}）。事实上，我们与朋友一同进入林中，每当我们审视属下的过失时。我们心地单纯地砍柴，当我们以虔诚的意向砍除犯错者的恶习时。斧头从手中飞出，当责备变得比需要的更尖刻时。铁刃从柄上脱落，当从斥责中说出过于严厉的话时。他击中并杀死朋友，因为过度的辱骂使他断绝了爱德的精神。因为被责备者的心突然生出仇恨，如果过度的责备不当地加于他。但那个不慎砍木而杀死邻人的人必须逃到三座城里，以便在其中一座得以生存；因为如果他借忏悔的哀哭，隐藏于圣事统一中的望德与爱德之下，他就不被视为犯有谋杀罪。被杀者的至亲即使找到他也不会杀死他；因为当那严格的审判者——他因分享我们的本性而与我们结合——来临时，他无疑不会向那个被信德、望德与爱德隐藏于赦免庇护下的人索要其过犯的惩罚。\par

% structure: p0045 source=h2 target=subsection reason=relative-heading-rank; base=h2

\subsection{第十一章}

\Emph{统治者应当多么专注地默想圣经法律。}\par

但这一切若能由统治者妥善执行，如果他受天上的敬畏与爱德之神感动，每日默想圣经的训诫，那么神圣劝诫的话语就能在他内恢复对天上生活的关切与远见的力量，而这力量因与人惯常的交往而不断丧失；并且，那个因世俗交往而被拖入旧生活的人，能通过痛悔的渴望不断更新对灵性故乡的爱。因为人心在人言中大大荒废；而且，既然无疑可见，当人被外在事务的喧嚣驱赶时，他失去平衡而跌倒，就应当不断地留意，通过热切追求教导而重新站起。为此，保禄劝诫被派牧放羊群的门徒说：\BibleQuote{在我到来之前，请专心于阅读}（\BibleRef{弟前 4:13}）。为此，达味说：\BibleQuote{上主，我多么爱慕您的法律！它是我全日的默想}（\BibleRef{咏 108:97}）。为此，主吩咐梅瑟关于抬约柜的事，说：\BibleQuote{你要用金子铸四个环，安在柜的四角上；用皂荚木作两根杠，包上金，穿在柜两旁的环内，以便抬柜；这些杠要常留在环内，不可抽出}（\BibleRef{出 25:12}）。约柜所预表的不就是圣教会吗？它被命令在四角安装四个金环，因为教会向地球四方延伸，无疑表明它装备了四部圣福音书，以便旅程。又用皂荚木作杠，穿在环内以便抬柜，因为需要寻找刚毅而恒久的教师，如同不朽坏的木料，他们因恒常依附于神圣经卷的教导而宣告圣教会的合一，并如同穿入环中而抬柜。事实上，用杠抬柜，就是通过善师们的宣讲，将圣教会带到不信者粗野的心灵前。这些杠也被命令包上金，以便他们在言语上对他人宣讲时，自己也以生活的光辉闪耀。关于他们，又适当地加上：\InnerQuote{这些杠要常留在环内，不可抽出}；因为那些从事宣讲职务的人，确实不应从神圣学问的研习中退出。为此，命令杠常留在环内，以便当需要抬柜时，不会因要穿杠而延迟；这就是说，当一位牧者被属下询问任何灵性事务时，如果他这时才去学习，那是极其可耻的，因为他本应能解答问题。但让杠常留在环内，使教师们常在心里默想圣经的话语，就能毫不延迟地抬起盟约之柜；这样，一旦需要，就能立刻教导。为此，教会的第一位牧者很好地劝诫所有其他牧者说：\BibleQuote{要常常准备，回答每位问你们心中盼望之理由的人}（\BibleRef{伯前 3:15}）；他仿佛在明说：为使抬柜没有延迟，杠决不可从环内抽出。\par

\SourceRecord{Translated by James Barmby.；Nicene and Post-Nicene Fathers, Second Series；Vol. 12.；Edited by Philip Schaff and Henry Wace.；Buffalo, NY: Christian Literature Publishing Co.,；1895.；<http://www.newadvent.org/fathers/36012.htm>.}

% structure: p0001 source=h1 target=section reason=page-title

\section{《司牧训》卷三}

论善牧者当如何教导及劝诫其属下\par

% structure: p0003 source=h2 target=subsection reason=relative-heading-rank; base=h2

\subsection{序言。}

既然我们已经说明了牧者应为怎样的人，现在让我们陈述他应当如何施教。因为正如久于我们之前、可敬的纳齐安的额我略所教导的，同样的劝诫并不适合所有的人，因为众人并非都受相同性格的束缚。凡有益于某些人的事，往往伤害另一些人；因为大部分情况下，滋养某些动物的草药对另一些动物是致命的；安抚马匹的轻嘶会激怒幼犬；缓解一种疾病的药物会加重另一种疾病；强健者赖以活力的面包会杀死幼童。因此，教师的话语应根据听众的品性加以塑造，以便适合每个人各自的所需，却从不偏离共同建树的艺术。因为听众的专注心灵，就如同竖琴上绷紧的琴弦，高明的琴手为使音调和谐，以不同的方式弹拨它们。因此琴弦发出协调的旋律，是因为它们虽被同一拨片拨动，却非以同一种方式。由此，每一位教师，为在爱德这一德行中建树众人，应从同一教义出发，却不可用同一劝诫去触动听众的心。\par

% structure: p0005 source=h2 target=subsection reason=relative-heading-rank; base=h2

\subsection{第一章}

论讲道艺术中应有的多样性。\par

下列之人当以不同方式加以劝诫：\par

\begin{itemize}
\item 男人与女人。
\end{itemize}

\begin{itemize}
\item 穷人与富人。
\end{itemize}

\begin{itemize}
\item 欢乐者与忧伤者。
\end{itemize}

\begin{itemize}
\item 长上与属下。
\end{itemize}

\begin{itemize}
\item 仆人与主人。
\end{itemize}

\begin{itemize}
\item 此世的智者和愚拙者。
\end{itemize}

\begin{itemize}
\item 厚颜者与腼腆者。
\end{itemize}

\begin{itemize}
\item 冒进者与胆怯者。
\end{itemize}

\begin{itemize}
\item 急躁者与忍耐者。
\end{itemize}

\begin{itemize}
\item 友善者与嫉妒者。
\end{itemize}

\begin{itemize}
\item 纯朴者与虚伪者。
\end{itemize}

\begin{itemize}
\item 健康者与患病者。
\end{itemize}

\begin{itemize}
\item 那些因畏惧鞭笞而生活无过者；以及那些在邪恶中变得如此刚硬，以致连鞭笞也不能纠正者。
\end{itemize}

\begin{itemize}
\item 过于沉默者与多言耗时者。
\end{itemize}

\begin{itemize}
\item 懒惰者与急躁者。
\end{itemize}

\begin{itemize}
\item 温良者与暴躁者。
\end{itemize}

\begin{itemize}
\item 谦卑者与傲慢者。
\end{itemize}

\begin{itemize}
\item 固执者与多变者。
\end{itemize}

\begin{itemize}
\item 贪食者与节制者。
\end{itemize}

\begin{itemize}
\item 那些仁慈施舍己物者，以及那些渴望夺取他人财物者。
\end{itemize}

\begin{itemize}
\item 那些既不夺取他人之物也不慷慨施舍己物者；以及那些施舍己物却仍不停止夺取他人之物者。
\end{itemize}

\begin{itemize}
\item 那些纷争者与那些和睦者。
\end{itemize}

\begin{itemize}
\item 好争论者与缔造和平者。
\end{itemize}

\begin{itemize}
\item 那些不按正确方式理解神圣律法之言者；以及那些确按正确方式理解，却无谦卑宣讲者。
\end{itemize}

\begin{itemize}
\item 那些虽能堪当宣讲，却因过度谦卑而畏惧者；以及那些因不成熟或年龄而被禁止宣讲，却因鲁莽而被驱使去宣讲者。
\end{itemize}

\begin{itemize}
\item 那些在世俗之事上如愿顺遂者；以及那些确实贪求世俗之物，却被逆境的劳苦折磨者。
\end{itemize}

\begin{itemize}
\item 那些受婚姻束缚者，以及那些不受婚姻羁绊者。
\end{itemize}

\begin{itemize}
\item 那些有过肉体交合经验者，以及那些对此无知者。
\end{itemize}

\begin{itemize}
\item 那些哀叹行为之罪者，以及那些哀叹意念之罪者。
\end{itemize}

\begin{itemize}
\item 那些痛哭恶行却不肯放弃者，以及那些放弃却不为痛哭者。
\end{itemize}

\begin{itemize}
\item 那些甚至赞美自己所行非法之事者，以及那些谴责错误却仍不避开者。
\end{itemize}

\begin{itemize}
\item 那些被突发的激情所胜者，以及那些蓄意陷于罪责者。
\end{itemize}

\begin{itemize}
\item 那些所行非法之事虽属轻微却行之频繁者；以及那些远离小罪，却偶尔陷入重罪者。
\end{itemize}

\begin{itemize}
\item 那些甚至不开始行善者，以及那些开始行善却完全不能完成者。
\end{itemize}

\begin{itemize}
\item 那些暗中作恶而公开行善者；以及那些隐藏所行之善，却在某些公开所为之事上容许人对自己有恶念者。
\end{itemize}

但我们以列表方式汇总这些内容，若不尽量简要地列出对每一种人的劝诫方式，又有何益？\par

（劝诫一）因此，男人与女人当以不同方式加以劝诫；因为对前者加给更重的诫命，对后者则加给较轻的诫命，好让前者因大事而受锻炼，后者因轻事而被亲切地归化。\par

（劝诫二）青年人与老年人当以不同方式加以劝诫；因为严厉的劝诫大多引导前者走向改善，而温和的规劝使后者更易行善。因为经上记载：\BibleQuote{不可严责老年人，但要劝他如父。}\BibleRef{弟前 5:1}\par

% structure: p0046 source=h2 target=subsection reason=relative-heading-rank; base=h2

\subsection{第二章}

如何劝诫穷人与富人。\par

(\Emph{劝告}三.) 应当以不同的方式劝告穷人和富人：因为对前者，我们应当提供安慰以对抗患难，但对后者，则要激起恐惧以对抗自高。因为关于穷人，主通过先知说：\BibleQuote{不要害怕，因为你必不致蒙羞}\BibleRef{依 54:4}。不久之后，主安慰她说：\BibleQuote{你这困苦的、被飘摇的可怜人}\BibleRef{依 54:11}。祂又安慰她说：\BibleQuote{我曾在穷困的炉鼎中拣选了你}\BibleRef{依 48:10}。但另一方面，\Person{保禄}对他的弟子论及富人说：\BibleQuote{你要嘱咐今世的富人，不要自高，也不要倚靠无定的财富}\BibleRef{弟前 6:17}；这里特别值得注意的是，这位谦逊的导师在提及富人时，不说\Emph{请求}，而说\Emph{嘱咐}；因为，虽然应对软弱者施以怜悯，但对自高者却不该给予尊敬。因此，对这样的人，所说的正确之事更应受命定，因为他们被对短暂事物的高傲思想所膨胀。关于他们，主在福音中说：\BibleQuote{你们富有的人有祸了！因为你们已经得了你们的安慰}\BibleRef{路 6:24}。因为他们不知道永恒的喜乐是什么，所以从现世生活的富足中得到安慰。因此，应向在穷困的炉鼎中受试探的人提供安慰；应向被现世荣耀的安慰所抬举的人激起恐惧；这样，前者可以学到他们拥有眼不能见的财富，后者意识到他们绝不能保持眼所能见的财富。然而，大多数情况下，人的品性会改变他们的位置；因此富人可能谦卑，穷人可能骄傲。因此，讲道者的口舌应迅速适应听者的生活，以便更尖锐地击打穷人的骄傲，因为连临到他身上的穷困也未能将其压下；并更温和地鼓励富人的谦卑，因为连那提升他们的富足也未使他们自高。\par

有时，一个骄傲的富人也应在劝告中用和言悦色来安抚，因为坚硬的疮肿也常通过温和的敷料变软，疯子的狂怒也常因医生的和言而恢复健康，当他们被愉快地哄着时，他们精神错乱的疾病便减轻了。因为也不该轻视这件事：当恶灵进入\Person{撒乌耳}时，\Person{达味}拿起他的琴安抚了他的狂乱\BibleRef{撒上 18:10}。因为\Person{撒乌耳}暗示了掌权者的骄傲，而\Person{达味}则暗示了圣者的谦卑生活。因此，当\Person{撒乌耳}被邪灵所缠，他的狂乱便因\Person{达味}的歌声而平息；同样，当掌权者的感官因骄傲而陷入疯狂时，就应当由我们言语的平静，如同琴声的甜美，将他们召回健康状态。但有时，当责备世上的掌权者时，应先通过某些比喻探询他们，仿佛事不关己；当他们就他人作出了正确的判决后，再以合适的方式针对他们自己的罪责责打他们；这样，被世俗权力膨胀的心就不会抬起来反抗责备者，因为它已凭自己的判断践踏了骄傲的颈项，且被自己的口所发出的判决束缚，不再试图为自己辩护。因此，\Person{纳堂}先知来责备君王时，先问他对一个穷人与一个富人之间的案件的判断\BibleRef{撒下 12:4}，好让王先宣判，然后才听到自己的罪，这样他绝不致反驳自己对自己所作的公义判决。因此，这位圣人，既考虑到罪人也考虑到君王，以奇妙的次序先以告解捆绑那大胆的罪犯，然后以谴责刺透他的心。他一时隐藏他所指向的对象，但当他抓住对方时便猝然打击。因为如果他有意从讲话一开始就公开打击罪恶，那打击或许会力量减弱；但通过先引入比喻，他磨利了隐藏的谴责。他如同医生来到病人面前；他看到疮肿必须切开；但他怀疑病人的忍耐。于是他把医用钢刀藏在袍子下，突然抽出并刺入伤口，使病人在看见刀之前就感到切肤之痛，免得他先看见而拒绝感受。\par

% structure: p0050 source=h2 target=subsection reason=relative-heading-rank; base=h2

\subsection{第三章}

\Emph{如何劝告喜乐的人和忧愁的人。}\par

(\Emph{劝告}四.) 应当以不同的方式劝告喜乐的人和忧愁的人。也就是说，在喜乐的人面前应设置惩罚后将临的悲伤之事；但在忧愁的人面前，应设置王国应许的喜乐之事。让喜乐的人从威胁的严厉中学到该惧怕什么；让忧愁的人承担起他们可以期待的赏报之喜乐。因为对前者，经上说：\BibleQuote{你们现在欢笑的人有祸了！因为你们将要哭泣}\BibleRef{路 6:25}；但对后者，他们从同一导师的教导中听到：\BibleQuote{我要再见你们，你们的心要喜乐，并且你们的喜乐没有人能夺去}\BibleRef{若 16:22}。但有些人并非因环境而成为喜乐或忧愁，而是由气质使然。对这样的人，应该让他们知道，某些缺陷与某些气质相关；即喜乐的人易犯邪淫，忧愁的人易犯愤怒。因此，每个人必须不仅要考虑他因自己独特的气质所受之苦，还要考虑与之相关的更恶之事在压迫他；免得他完全不与已有的事斗争，反而还屈服于那些他自认为没有的事。\par

% structure: p0053 source=h2 target=subsection reason=relative-heading-rank; base=h2

\subsection{第四章}

\Emph{如何劝告属下和主教。}\par

(\Emph{劝诫} 5.) 对于属众和长上，应以不同方式加以劝诫：前者不可因服从而消沉，后者不可因高位而自傲；前者不可不完成所命之事，后者不可命令超出正义之事；前者应谦卑服从，后者应节制治理。因为这话（也可作比喻理解）是对前者说的：\BibleQuote{\Emph{孩子们，要在主内服从你们的父母}}，但对后者命令说：\BibleQuote{\Emph{你们作父亲的，不要激怒你们的子女}}\BibleRef{哥 3:20}。愿前者学习如何在隐藏审判者的眼前调整自己的内心思想；后者则学习如何为那些托付给他们的人在外表上提供善生的榜样。因为长上应当知道，若他们做出任何错事，他们便当得多少死亡，就正如他们向属众传递了多少丧亡的榜样。因此，他们必须更加谨慎地防范罪恶，因为藉着他们的恶行，他们不仅自己死亡，而且因着他们的恶表毁灭了他人，对别人的灵魂负有罪责。因此，前者应受劝诫，以免他们若只靠自身无法辩白时，被严厉公开；后者应受劝诫，以免他们因属众的过失而受审判，即使他们自身觉得安稳。前者应受劝诫，因他们不受照顾他人的牵累，当更加忧心自己；而后者应受劝诫，完成对他人之职责，同时不放弃对自己的责任，既热切关心自己，也不在照管托付给之人上懈怠。对前者，即空闲于己事之人，有话说：\BibleQuote{\Emph{懒惰人，你去察看蚂蚁，思想它的行径，学习智慧}}\BibleRef{箴 6:6}。但对后者，有可怕的劝诫说：\BibleQuote{\Emph{我儿，你若为朋友作保，你就双手许给陌生人，被自己口中的话缠住，被自己的言语捉住}}\BibleRef{箴 6:1}。因为为朋友作保，就是以自身行为的担保来照顾另一个人的灵魂。因此，手许给陌生人，因为心被以前没有的责任所束缚。但他被自己口中的话缠住，被自己的言语捉住，因为当被迫向托付给他的人说善言时，他自己首先必须遵守他所说的话。因此，他被自己口中的话缠住，被理性的要求所约束，不容他的生活松懈到与他的教导不符的程度。因此，在严厉的审判者面前，他必须在行为上成就他用声音命令别人的一切。因此，在上述引文中，接着也添加了这劝诫：\BibleQuote{\Emph{所以，我儿，你要听我的话，救自己脱身，因为你已落入邻人的手中：四处奔跑，快快行动，唤醒你的朋友；不可让眼睛睡眠，也不可让眼皮打盹}}\BibleRef{箴 6:3}。因为凡被置于他人之上作为生活榜样的人，不仅受劝诫要自己警醒，还当唤醒他的朋友。因为若他也不将那被他管理的人从罪恶的麻木中唤醒，仅凭善生而警醒是不够的。因为经上说得好：\BibleQuote{\Emph{不可让眼睛睡眠，也不可让眼皮打盹}}\BibleRef{箴 6:4}。因为给眼睛睡眠，就是停止热忱，以至于完全忽略对属下的照顾。但是当我们的思想为懒惰所压，姑息明知应在属下身上责备的事时，眼皮便打盹。因为沉睡，就是既不知道也不纠正托付给我们之人的行为。然而知道哪些事该受责备，却因心灵懒惰不以适当的斥责改正，这不是睡觉，而是打盹。然而，眼睛从打盹进入沉睡；因为通常，当管理他人者不割除他知道的恶时，他迟早会因自己的疏忽，甚至不再知道属下所做的错事。\par

为此，那些管理他人的人应受劝诫，因谨慎的热忱，他们要有眼目内里和四周警醒，并努力成为天上的活物\BibleRef{则 1:18}。因为天上的活物被描述为周围和内部都满了眼睛\BibleRef{默 4:6}。因此，那些管理他人的人应当有眼目内里和四周，既要自身努力取悦内心的审判者，也要在外表提供生活的榜样，察觉别人身上应纠正的事。\par

应劝告属下，不可轻率判断长上的生活，若偶见他们行事有可指责之处，免得他们因正确指责恶事，反被骄傲的冲动沉入更深的深渊。应劝告他们，当思虑长上的过失时，不可因此对他们过于放肆；但若长上有些行为极其恶劣，他们应在内心如此判断，以至因敬畏天主而仍不肯拒绝在长下背负敬畏之轭。我们将藉达味所行之事更好地说明这一点\BibleRef{撒上 24:4}。因为当迫害者撒乌耳进入山洞解手时，长期受他迫害的达味正与手下在山洞内。当手下怂恿他击杀撒乌耳时，达味以回答阻止他们，说他不可伸手加害上主的受傅者。然而他悄悄起身，割下了撒乌耳外氅的边缘。撒乌耳代表什么？岂不是坏统治者？达味代表什么？岂不是好臣民？撒乌耳解手，意味着统治者将心中孕育的邪恶扩展为散发恶臭的行为，并以行动彰显内在的污秽思想。然而达味不敢击杀他，因为虔诚臣民的心灵，克制自己不染全盘毁谤的瘟疫，即使用舌剑斩杀长上的生命，即便他们指责长上的不完善。当他们因软弱几乎无法忍住，即便谦卑地谈论长上某些极端明显的恶行时，他们便如悄悄割下外氅的边缘；因为即使无害而秘密地贬损长上的尊严，他们便如同丑化了置于他们之上的君王的外衣；然而他们仍回到自身，并最严厉地自责，即使是最轻微的言语毁谤。因此，那处也恰当写道：\BibleQuote{后来达味的心责打他，因为他割去了撒乌耳外氅的边缘}\BibleRef{撒上 24:6}。因为长上的行为，即使被正确判断为应受责备，也不应用口舌之剑来攻击。但若舌头哪怕稍一不慎滑入对他们的批评，心灵必须因补赎的痛苦而沮丧，为的是回到自身，并因冒犯了置于其上的权柄，而畏惧那设立权柄者的审判。因为当我们冒犯置于我们之上的人时，我们便是违抗了那设立他们于我们之上的天主的旨意。因此，梅瑟得知百姓抱怨他和亚郎时，便说：\BibleQuote{我们算什么？你们的抱怨不是针对我们，而是针对上主}\BibleRef{出 16:8}。\par

% structure: p0058 source=h2 target=subsection reason=relative-heading-rank; base=h2

\subsection{第五章}

\Emph{如何劝告仆人与主人。}\par

（劝告六）。仆人应受不同的劝告；仆人应常谨记其身份之谦卑；主人则不应忘记其本性，在此本性中他们与仆人处于平等地位。应劝告他们，不要因傲慢行事而违抗天主的命令；也要劝告他们，若因身份而视属下为不同类，便是对天主恩赐的骄傲。前者应承认自己是主人的仆人；后者应承认自己是仆人的同伴仆人。因为对他们说：\BibleQuote{仆人，你们要事事听从你们肉身的主人}\BibleRef{哥 3:22}；又说：\BibleQuote{凡在轭下作仆人的，应视自己的主人为配受一切尊敬}\BibleRef{弟前 6:1}；但对他们说：\BibleQuote{你们作主人的，也要同样对待他们，戒除恐吓，因为知道他们和你们的主在天上}\BibleRef{弗 6:9}。\par

% structure: p0061 source=h2 target=subsection reason=relative-heading-rank; base=h2

\subsection{第六章}

\Emph{如何劝告聪明人与愚钝人。}\par

（\Emph{劝言}七）。应劝戒世上智者与愚钝者，方式有别。劝戒智者，须使其放弃所知的；劝戒愚钝者，须使其寻求所不知的。对于前者，首先须摧毁他们自以为有智慧的想法；对于后者，则须将已认识的天上智慧建造起来，因为他们毫无骄傲，犹如预备好了心灵以承载建筑工程。对于前者，我们当努力使他们变得更有智慧的愚拙，放弃愚拙的智慧，学习天主的智慧的愚拙；对于后者，我们当宣讲，使他们从被视为愚拙的事，如同从近邻之处，过渡到真正的智慧。因为前者曾受训示说：\Quote{若有人在今世自以为有智慧，就让他变成愚拙，好成为有智慧的}\BibleRef{格前 3:18}；但对后者却说：\Quote{按肉身来说，有智慧的并不多}\BibleRef{格前 1:26}；又说：\Quote{天主拣选了世上愚拙的，为叫那有智慧的羞惭}\BibleRef{格前 1:27}。前者多半藉推理的辩论而悔改；后者有时藉榜样更好。无疑，前者因在己方的论断中被驳倒而获益；但对于后者，有时仅知道他人可称赞的作为便足矣。为此，那位欠了有智慧的和愚拙者债的卓越导师\BibleRef{罗 1:14}，当他劝戒某些有智慧的希伯来人，也有一些较为迟钝的时，对他们谈论旧约的应验，藉辩论战胜前者的智慧，说：\Quote{那衰朽衰老的，快要消失了}\BibleRef{希 8:13}。但当他察觉有些人只能藉榜样被吸引时，在同一封信中他又加上：\Quote{圣人们经历了戏弄、鞭打、甚至锁链和监禁；他们被石头砸死，被锯死，被诱惑，被刀剑杀死}（\BibleRef{希 11:36,37}）；又说：\Quote{你们应记念那些领导过你们、给你们讲过天主圣言的人，追随他们的信德，注视他们生活的结局}\BibleRef{希 13:7}；如此，得胜的理性可征服一类人，而榜样的温和力量则劝服另一类人攀登更高之事。\par

% structure: p0064 source=h2 target=subsection reason=relative-heading-rank; base=h2

\subsection{第 7 章}

应如何劝戒无羞耻者和腼腆者。\par

（\Emph{劝言}八）。应劝戒无羞耻者与腼腆者，方式有别。对于前者，只有严厉的斥责才能制止无羞耻的恶习；而对于后者，温和的劝戒多半能使之改过。前者若不受多人斥责，便不知自己有错；后者通常只须导师温和地提醒他们的恶行，就足以使他们悔改。对于前者，最好以直接斥责来纠正；但对于后者，若以侧面轻触的方式责备他们，则更有益处。因此，主公开斥责无羞耻的犹太人百姓，说：\Quote{你竟有娼妓的面额，不知羞耻}\BibleRef{耶 3:3}。但祂又安慰羞惭的他们，说：\Quote{你将忘记你幼年的羞耻，不再记念你寡居的羞辱，因为你的造主要作你的夫君}\BibleRef{依 54:4}。保禄也公开斥责无羞耻犯罪的迦拉达人，说：\Quote{无知的迦拉达人！谁迷惑了你们}\BibleRef{迦 3:1}？又说：\Quote{你们竟这样无知吗？你们既随从圣神开始，如今却要依靠肉身成全吗}\BibleRef{迦 3:3}？但他以同情之态责备那些羞惭之人的过错，说：\Quote{我在主内大大喜乐，因为你们对我终于又生出挂念，你们素来就挂念我，只是缺少机会}\BibleRef{斐 4:10}；这样，严厉的斥责揭露前者的过错，而温和的言辞则遮掩后者的疏忽。\par

% structure: p0067 source=h2 target=subsection reason=relative-heading-rank; base=h2

\subsection{第 8 章}

应如何劝戒傲慢者和怯懦者。\par

（\Emph{劝言}九。）应劝戒傲慢者与怯懦者，方式有别。前者过于自恃，受责备时藐视众人；后者过分意识到自己的软弱，多半陷入沮丧。前者认为自己所作的一切都出类拔萃；后者以为自己所行极其被人轻视，因而沮丧消沉。因此，责备者应细腻地筛检傲慢者的行为，使他们自喜之处显得得罪天主。\par

因为我们纠正傲慢者最好的方式，是表明他们自以为做对的事其实是做错了；这样，他们认为从中获得荣耀之处，反而变成有益的羞愧。但有时，当他们完全没有意识到自己犯了傲慢的罪时，若能旁敲侧击地举出另一个人的更明显的丑行，会使他们更快地得到纠正，从而从他们无法辩护之处，意识到他们错误地坚持自己所辩护的。因此，保禄见格林多人彼此傲慢地自夸，一个说我是属保禄的，另一个说我是属阿波罗的，又一个说我是属刻法的，还有一个说我是属基督的\BibleRef{格前 1:12}，便举出他们中间不仅发生了而且未被纠正的乱伦罪行，说：\Quote{听说在你们中间有淫乱的事，并且这样的淫乱连外邦人中也没有，就是有人收了他的继母。你们还是傲慢自大，而不该悲哀，把行这事的人从你们中间除去吗}\BibleRef{格前 5:1}。这无异于明说：你们傲慢地说自己是属于这个或那个，而在你们放纵的疏忽中，却显出你们不属于任何人。\par

然而，另一方面，若能同时寻找他们身上的某些优点，则更适于将胆怯者带回行善之路。这样，我们以责备攻击他们某些方面，又以称赞拥抱其他方面；使称赞的听闻滋养他们的软弱，而责备则惩戒他们的过错。通常，我们也能更有效地帮助他们获益：如果我们提及他们的善行，并在他们作了一些错事时，不以其已犯而责备，反而像禁止不应发生之事那样劝诫。如此，恩惠能增加我们所认可的事，而谦逊的劝诫对胆怯者对抗我们所指责的事更为有效。因此，同一保禄得知得撒洛尼人虽坚守所接受的宣讲，却因某种胆怯而困扰，仿佛世界末日即将来临。他首先称赞他们所表现出的坚强之处，随后以谨慎的劝勉坚固他们的软弱。他这样说：\BibleQuote{\Emph{弟兄们，我们应当常为你们感谢天主，这本是相称的，因为你们的信德格外增长，你们众人彼此相爱的爱德充足，以致我们在天主的各教会中也为你们的坚忍和信德而夸耀你们。}}\BibleRef{得后 1:3}但在这些赞美的颂词之后，稍后他又加上：\BibleQuote{\Emph{弟兄们，我们恳求你们，因我们主耶稣基督的来临，和我们聚集到他面前，不要轻易动摇理智，也不要惊慌，无论藉着神恩、话语或冒充我们发出的书信，好像主的日子已经近了。}}\BibleRef{得后 2:1}真正导师的教导方式应是：先让他们在受赞美时听到值得感恩的事，随后在受劝勉时看到应当追随的事；使先前的赞美安定他们的心，以免后续的劝诫动摇它。尽管他知道他们因对末日临近的疑虑而困扰，他并未责备他们已然如此，反而如同不知过去一般，禁止他们将来再困扰。这样，当他们相信自己连这种困扰的轻浮也不为导师所知时，就会像因被导师知晓而惧怕受责一样，害怕自己真正可受责备。\par

% structure: p0072 source=h2 target=subsection reason=relative-heading-rank; base=h2

\subsection{第九章}

\Emph{如何劝诫不耐烦者与忍耐者}\par

(\Emph{劝诫} 10.) 对不耐烦者与忍耐者的劝诫应有所不同。不耐烦者应被告知：他们因忽视约束自己的心灵，往往在不知不觉中陷入许多罪的险境；因为愤怒驱赶心智前往欲望不引导的地方，而心智在激动中做出事后知晓便后悔的事。不耐烦者还应被告知：当他们被情绪冲动所驱使时，有时如同丧失理智般行事，事后几乎意识不到自己所行的恶；并且，由于他们不抵抗自身的激动，甚至将心智平静时可能做好的事情也一并搅乱，在突如其来的冲动下摧毁了可能经过谨慎劳苦长久建立的一切。因为爱德这一众德之母与守护者，因不耐烦的恶德而丧失。经上记载：\BibleQuote{\Emph{爱德是含忍的}}\BibleRef{格前 13:4}因此，哪里有忍耐，哪里就有爱德；哪里没有忍耐，哪里就没有爱德。同样，因不耐烦的恶德，作为诸德之母的教导也被消散。经上记载：\BibleQuote{\Emph{人的智慧使他缓于怒怒，忍耐是他的光耀。}}\BibleRef{箴 19:11}因此，每个人愈被证明缺乏忍耐，便愈显得缺乏教导。因为如果他在生活中不懂得如何平静地忍受他人的恶行，就不能真正通过教导传授善。\par

此外，通过这急躁的恶习，骄傲的罪往往刺入心灵；因为当一个人不能忍受在此世被轻看时，他便试图炫耀自己隐藏的任何善，于是通过急躁被引向骄傲；当他不能忍受轻蔑时，就以炫耀自我展示而夸耀。为此经上记载：\Emph{\BibleQuote{忍耐的人胜于傲慢的人}}\BibleRef{训 7:9}；因为，事实上，忍耐的人宁愿忍受任何邪恶，也不愿他的隐藏的善因炫耀的恶习而被知晓。但傲慢的人相反，宁愿有人向他夸耀甚至假装的好，免得他可能遭受哪怕最小的恶。既然，当忍耐被放弃时，所有其他已行的善事也被推翻，那么正确地对\Person{厄则克耳}吩咐说要在天主的祭台上挖一个沟渠；即，在沟渠中，放在祭台上的全燔祭得以保存\BibleRef{则 43:13}。因为，如果祭台没有沟渠，飘过的风会吹散它在那里找到的任何祭品。而我们所说的天主的祭台，不就是义人的灵魂吗？它在天主面前献上多少祭品，就是行了多少善事。而祭台的沟渠不就是善人的忍耐吗？它使心灵谦卑以忍受逆境，像沟渠一样将它置于低处。因此，在祭台上要挖一个沟渠，免得风吹散放在上面的祭品：也就是说，被选者的心灵要保持忍耐，免得被急躁的风吹动，连那行得好的也失去。同样，这沟渠被吩咐要有一肘宽，因为如果忍耐不缺失，统一的尺度就得以保存。为此\Person{保禄}也说：\Emph{\BibleQuote{你们要彼此担待重担，这样就满全了基督的法律}}\BibleRef{迦 6:2}。因为基督的法律就是统一的爱德，只有那些即使在被重担压着时也没有过失的人才满全它。让急躁的人听听所记载的：\Emph{\BibleQuote{忍耐的人胜于勇士，控制自己心灵的人胜于攻取城池的人}}\BibleRef{箴 16:32}。因为攻取城池是小事，因为所征服的是外在的；但通过忍耐所征服的则大得多，因为心灵本身被自己克服，并在忍耐迫使它克制自己时降服于自己。让急躁的人听听真理对祂的被选者说的话：\Emph{\BibleQuote{你们要凭着忍耐保全你们的灵魂}}\BibleRef{路 21:19}。因为我们被造得如此奇妙：理性拥有灵魂，灵魂拥有身体。但灵魂若不被理性首先拥有，它就被剥夺了对身体所有权的权利。因此主指出忍耐是我们状态的守护者，祂教导我们在忍耐中拥有自己。这样我们就知道急躁的罪有多大，通过它我们失去了我们之所是的拥有权。让急躁的人听听再次通过\Person{撒罗满}说的话：\Emph{\BibleQuote{愚人尽显其心意，智者却延迟，保留到以后}}\BibleRef{箴 29:11}。因为一个人被急躁的冲动所驱使，以致整个心意被倾吐出来，内心的扰动更快地抛出它，因为智慧的纪律没有围绕它。但智者延迟，保留到以后。因为当受伤害时，他不愿现在报应自己，因为在他的容忍中他甚至希望人应被宽恕；但他并非不知道一切在最后的审判中都会被正义地报应。\par

另一方面，应该劝告忍耐的人，不要为外在所承受的而内心悲伤，免得他们用内在的恶意感染败坏这如此宝贵的祭献，这祭献他们完整地在外奉献；并且免得他们的悲伤之罪，虽不被人们察觉，但在天主的审视下被视为罪，就更加恶劣，因为它在外人面前自诩为德行的美丽外表。\par

因此，应劝谕忍耐者努力去爱那些他们必须容忍的人；以免爱德不随忍耐而来，所显出的德行反而转为更坏的憎恨之罪。为此，保禄说：\BibleQuote{\Emph{爱德是忍耐}}，接着立即加上：\BibleQuote{\Emph{是仁慈}}\BibleRef{格前 13:4}；这无疑表明，爱德以忍耐所容忍的人，也以仁慈不懈地爱着他们。这位卓越的教师同样，在劝勉门徒忍耐时说：\BibleQuote{\Emph{一切毒辣、怨恨、忿怒、吵闹、毁谤，都应当从你们中除掉}}\BibleRef{弗 4:31}，仿佛已将外在的一切整理就绪，转而关注内在的事，接着加上：\BibleQuote{\Emph{连同一切恶意}}\BibleRef{弗 4:31}；因为，的确，若内在的恶意——诸恶之母——掌权，则徒然去掉外在的怨恨、吵闹、毁谤；若邪恶被切断在外枝，却存于内根，必将以更多方式重新萌发。为此，真理本人也说：\BibleQuote{\Emph{要爱你们的仇敌，善待恨你们的人，为迫害你们并诬蔑你们的人祈祷}}\BibleRef{路 6:27}。因此，在人前容忍仇敌是德行；但在天主前爱仇敌才是德行；因为天主所悦纳的唯一祭献，是在祂眼前，在善工的祭台上，由爱德之火焰所点燃的。因此，祂对某些忍耐却未怀有爱德的人说：\BibleQuote{\Emph{为什么你看见你弟兄眼中的木屑，却看不见自己眼中的大梁？}}\BibleRef{玛 7:3}。因为不耐烦的纷扰是一粒木屑；而心中的恶意却是眼中的大梁。前者被诱惑的微风吹得动摇不定；后者却因根深蒂固的邪恶几乎无法移动。然而，经上接着恰当地说：\BibleQuote{\Emph{你这假冒为善的人！先去掉你眼中的大梁，然后才能看得清楚，去掉你弟兄眼中的木屑}}\BibleRef{玛 7:3}；这仿佛是对那邪恶的心灵说的：它心中忧伤，却在外表藉忍耐显出圣善，说：\Quote{先摆脱你身上恶意的重担，然后才能责备他人不耐烦的轻率；否则，你若不肯努力战胜伪装，容忍他人过错对你更为不利。}\par

因为忍耐者常发生这样的事：当他们在受苦或听到侮辱时，确实未受烦恼打击，且表现出忍耐，以致也保持了心地纯洁；但过后，当他们忆起所忍受的这些事时，便以烦恼之火燃烧自己，寻找报复的理由，并在回顾中，将忍受时的温良转化为恶意。若这种变化的原因被揭露，他们便能更快得到讲道者的帮助。因为狡猾的仇敌向两人发动战争：即煽动一人先出言侮辱，并激怒另一人在受伤害之感下还口。但大多数情况下，当他已是那被说服施害之人的征服者时，他却被那以平静心忍受伤害之人所征服。因此，他既已战胜那被他激怒而制伏的人，就全力攻击另一人，却因他的坚定抵抗和得胜而懊恼；于是，由于在公然辱骂中未能动摇他，便暂时停止公开争斗，而藉秘密暗示激发他的思想，寻求合适的时机来欺骗他。因为他在公开战斗中失败，便热切地设下隐藏的陷阱。在平静之时，他回到得胜者的心中，使他记起所受的暂时损害或侮辱的利箭，并极度夸大一切所遭受的，使之显得无法忍受：他以如此巨大的烦恼扰乱心神，以致那忍耐者，在得胜后被掳，竟因曾平静地忍受这些事而羞愧，后悔没有还口，并寻求若有机会，就回报更恶的事。那么，这些人像谁呢？不是像那些在战场上凭勇武得胜，却因疏忽随后在城门口被擒的人吗？不是像那些为猛烈疾病所袭却不致死，却因渐渐复发的热症而死亡的人吗？因此，应劝诫忍耐者：在得胜后守护自己的心；要警惕在公开战争中被击败的仇敌，他正对着他们心灵的城墙设下陷阱；更要害怕疾病再次悄悄袭来；以免狡猾的仇敌若后来欺骗了他们，就因践踏曾长期不屈于他的征服者的脖子而更加欢喜。\par

% structure: p0079 source=h2 target=subsection reason=relative-heading-rank; base=h2

\subsection{第十章}

\Emph{如何劝诫心平气和之人与心怀嫉妒之人。}\par

（\Emph{劝言} 11.）对心地善良者和嫉妒者应加以不同的劝诫。对心地善良者应劝诫他们，要如此因他人身上的善行而喜乐，以至渴望自己也拥有同样的善行；要如此以爱德赞美邻人的行为，以至也藉效法来增加这些善行，免得在这现世生命的赛场上，他们只作热心的支持者而观看他人的竞赛，却作懒散的旁观者，竞赛结束后仍两手空空，因为他们在竞赛中未曾劳苦，那时他们将悲伤地看着那些在其劳苦中他们曾袖手旁观的人获得棕榈枝。因为我们若不爱他人的善行，实在犯了大罪；但若我们所爱的善行不尽己力去效法，也得不到赏报。因此，应告诉心地善良者：若他们不急于效法自己所称赞的善行，那么，德行之美在他们心中就如愚蠢的观众喜爱杂技表演的虚荣一般：这些人对车夫和演员的表演报以掌声，却并不渴望成为他们所赞扬的人那样。他们赞赏那些人做了令人喜悦的事，自己却逃避以同样方式令人喜悦。应告诉心地善良者：当他们看到邻人的行为时，应当返回自己的内心，不要自负于不属于自己的行为，也不要赞美善行却拒绝实行。确实，那些喜爱却不愿效法的人，最终将受到更沉重的惩罚。\par

应劝诫嫉妒者，让他们看到自己的盲目是多么严重：他们因他人的进步而失败，因他人的喜乐而憔悴；他们的不幸是多么巨大：因邻人的改善而变得更糟，看到他人福乐的增长，内心便不安地苦恼，死于自己心灵的瘟疫。还有什么比这些人更不幸呢？他们一看到幸福，就因痛苦而变得更邪恶。而且，他人的善行他们虽不能拥有，但如果他们爱这些善行，就可以使之成为自己的。因为所有人在信仰中共同构成一个身体，如同一个身体上的许多肢体；这些肢体职责各异，但彼此和谐一致而成为一体。因此，脚藉着眼看，眼藉着脚行走；耳朵的听觉服务于口，口舌之音与耳朵相互配合；肚子支撑着手，手为肚子工作。所以，在身体的安排中，我们学到了在行为中应当遵守的。那么，不按我们所是去行，就太可耻了。事实上，那些我们在他人的身上所爱的善行，即使我们无法追随，也成了我们的；而那些在我们身上被爱的善行，也成了爱它们的人所有。因此，让嫉妒者思考爱德有多么大的力量：它使不属自己的劳苦工作，无需劳苦就成为我们的。所以应告诉嫉妒者，当他们不能克制恶意时，就正沉入那狡猾仇敌的古老邪恶中。因为经上论到他说：\BibleQuote{因魔鬼的嫉妒，死亡进入了世界}（\BibleRef{智 2:24}）。因为他自己失去了天国，便嫉妒那被造的人拥有天国；他自己毁灭了，便藉着毁灭他人来加重自己的惩罚。应劝诫嫉妒者，让他们知道，他们脚下滋生着多少毁灭的深渊；因为只要他们不把恶意从心中赶出，就会滑向公开的恶行。因为加音若不是嫉妒弟弟被悦纳的祭品，绝不会去杀害弟弟。因此经上记载：\BibleQuote{上主惠顾了亚伯尔和他的祭品，却没有惠顾加音和他的祭品。加音因此大怒，垂头丧气}（\BibleRef{创 4:4}）。这样，因祭品而生的嫉妒就成了杀害兄弟的温床。那个因比他好而使他烦恼的人，他把他从世上除掉了。应告诉嫉妒者，当他们用这种内在的瘟疫消耗自己时，就毁掉了自己内里似乎存有的一切善。因此经上说：\BibleQuote{心安是肉体的生命，嫉妒是骨头的腐烂}（\BibleRef{箴 14:30}）。因为肉体象征某些软弱和柔弱的行动，骨头象征勇敢的行动。而且常常发生这样的事：有些人内心纯洁，在某些行动上似乎软弱；而另一些人，虽在人前做出一些勇敢的事，内心却因嫉妒他人的善行而受着瘟疫的折磨。所以经上说得好：\BibleQuote{心安是肉体的生命}，因为如果保持心灵的纯洁，即使是外在软弱的事也会及时得到加强。同样，紧接着说\BibleQuote{嫉妒是骨头的腐烂}也是正确的，因为通过憎恶的恶习，在人眼中看似坚强的事，在天主眼中却消失了。因为嫉妒使骨头腐烂，意味着某些甚至坚强的事也会完全毁灭。\par

% structure: p0083 source=h2 target=subsection reason=relative-heading-rank; base=h2

\subsection{第十一章}

\Emph{如何劝诫单纯者和狡猾者。}\par

（\Emph{劝诫}12.）对纯朴者与不诚实者的劝诫方式应有所不同。纯朴者值得称赞，因他们力求从不言谎，但应劝诫他们有时也当对真理保持缄默。因为，正如谎言常损害说话者，同样，听闻真理有时亦伤害某些人。因此，主在门徒前以沉默调和言语，说：\BibleQuote{我还有许多事要告诉你们，然而你们现在不能担负}\BibleRef{若 16:12}。所以应劝诫纯朴者，既然他们常有益地避免欺骗，也应有益地宣示真理。应劝诫他们，在纯朴的良善上添加谨慎，好使他们既能拥有纯朴的安全，又不失谨慎的明辨。因此，外邦人的导师说：\BibleQuote{我愿你们在善上明智，在恶上纯朴}\BibleRef{罗 16:19}。同样，真理本身也劝诫祂的选民：\BibleQuote{你们要机警如同蛇，纯朴如同鸽子}\BibleRef{玛 10:16}；这是说，在选民心中，蛇的机警应激发鸽子的纯朴，鸽子的纯朴应调和蛇的机警，好使他们既不因谨慎而陷入诡诈，也不因纯朴而在悟性上变得迟钝。\par

然而，另一方面，应劝诫不诚实者，学习他们因罪过所忍受的双重劳苦是何等沉重。因为他们害怕被察觉，便不断寻求不诚实的辩护，被恐惧的猜疑所搅扰。但再没有比真诚更安全的辩护，再没有比真理更易说出口的。因为，当被迫为自己的欺骗辩护时，心灵为艰苦的劳役所疲累。因此经上记载：\BibleQuote{他们嘴唇的劳苦必遮盖他们}\BibleRef{咏 138:10}。因为那如今充满他们的，那时将遮盖他们；那如今以轻微的纷扰抬举其灵魂的，那时将以严厉的惩罚重压它。因此，通过耶肋米亚说：\BibleQuote{他们教自己的舌头说谎，劳苦作恶}\BibleRef{耶 9:5}；这如同明说：那些本可轻易成为真理之友的人，劳苦地犯罪；他们拒绝活在纯朴中，却以劳苦求死。因为，通常当他们在过失中被逮住时，因不愿被人知道自己的真实面目，便以欺骗的帷幕隐藏自己，并竭力为自己已被清楚察觉的罪行辩解；以致那本应谴责其过犯的人，常被环绕的谎言之雾所迷惑，发现自己几乎失去了刚才对某人确知的事实。因此，先知以犹大的比喻，恰当地对那犯罪并自辩的灵魂说：\BibleQuote{在那里，箭猪安了窝}\BibleRef{依 34:15}。因为箭猪之名表示不诚实且狡猾地自卫的双重心灵；因为，当捕获箭猪时，它的头被看见，脚也显露，全身暴露无遗；但一被捕获，它就蜷缩成团，缩回脚，藏起头，之前全部可见的一切，在捉拿者手中便全部消失了。确实，不诚实的心灵在罪恶中被捉拿时也是如此。因为箭猪的头可见，即看出罪人从何开端发展到他的罪行；箭猪的脚可见，即发现罪过是通过哪些步骤实行的；然而，不诚实的心灵突然提出借口，缩回它的脚，即隐藏其罪恶的一切痕迹；它收回头，即以奇怪的辩护表明它甚至没有开始任何恶事；它在捉拿者手中仿佛一个球，因为那责斥它的人突然失去刚刚才得知的一切，把卷缩的罪人握于自己的意识中，尽管他曾在捕捉时看见了整个罪人，却被不诚实的辩护的推诿所迷惑，同样对整个人一概不知。这样，箭猪在恶人中有窝，因为诡诈心灵的虚伪，在自我辩护的黑暗中自我蜷缩隐藏。\par

让不诚实的人听听所写的：\BibleQuote{行于纯朴者，行得稳妥}\BibleRef{箴 10:9}。因为行为纯朴确是极大安全的保证。让他们听听智者口中所说的：\BibleQuote{训诲的圣神必逃避欺诈}\BibleRef{智 1:5}。让他们再听听圣经见证所肯定的：\BibleQuote{主的交谊与纯朴的人同在}\BibleRef{箴 3:32}。因为天主的交谊就是祂借其临在的光照，向人心启示奥秘。因此，祂被说成与纯朴的人相交，是因为祂以眷顾的光辉照亮那些没有被任何狡诈阴影遮蔽的头脑，使他们领悟天上的奥秘。但心怀二意者的特别邪恶在于：当他们以歪曲和双重行为欺骗别人时，却自夸比别人更有远见；由于不考虑报应的严厉，这些可怜的人竟为自己的损失而欢欣。但让他们听听\Person{索福尼亚}先知如何向他们展示天主斥责的威力，说：\BibleQuote{看，主的日子来临了，伟大而可怕，那是愤怒的日子，那日子是黑暗和阴沉的日子，是云和旋风的日子，是号角和呐喊的日子，临于一切设防的城市和一切高耸的角落}\BibleRef{索 1:15}。因为所谓设防的城市，岂不是指那些多疑的头脑，它们始终被虚假的防御所包围；每当它们的过错受到攻击时，就不让真理的利箭触及它们？而高耸的角落（墙角总是双重的）又象征什么？岂不是不诚实的心：它们回避真理的纯朴，在邪恶的双重性中仿佛自我折叠，更糟的是，由于不诚实的过错，它们竟以精明的骄傲抬高自己的思想？因此，主的日子充满惩罚和斥责，临于设防的城市和高耸的角落，因为最后审判的愤怒既摧毁那些被防御封闭而抗拒真理的人心，也展开那些在双重性中卷曲的心。因为那时设防的城市倒塌了，因为天主没有穿透的灵魂将被定罪。那时高耸的角落崩溃了，因为那些在不诚实的精明中自高自傲的心，被正义的判决击倒。\par

% structure: p0088 source=h2 target=subsection reason=relative-heading-rank; base=h2

\subsection{第十二章}

\Emph{如何劝诫健康者和患病者}\par

（\Emph{劝诫} 13。）对于健康者和患病者，劝诫的方式应有所不同。对于健康者，应劝诫他们善用身体的健康以促进灵魂的健康：以免他们将所赐健康的恩宠用于不义，反而因恩赐而变得更糟，日后应受更重的惩罚，因为他们现在竟然不怕滥用天主更丰厚的恩赐。对于健康者，应劝诫他们不可轻视赢得永恒健康的机会。因为经上记载：\BibleQuote{看，如今正是悦纳的时候；看，如今正是救恩的日子}\BibleRef{格后 6:2}。应劝诫他们，以免当他们可能取悦天主时不愿去做，到太迟时想做却无能为力。因此，智慧后来抛弃那些久久拒绝她起初呼唤的人，说：\BibleQuote{我呼唤了，你们却拒绝；我伸出了手，没有人理会；你们轻视了我所有的劝告，不接受我的斥责；我也要在你们的毁灭中发笑，在你们所惧怕的来临时加以嘲笑}\BibleRef{箴 1:24}。又说：\BibleQuote{那时他们呼求我，我必不应允；他们清晨寻求我，必寻不着我}\BibleRef{箴 1:28}。同样，当为行善而赐予的身体健康遭到轻视时，一旦失去，才感到这恩赐何等珍贵；最后，当它在适当的时候被赐予却无益地享用过后，现在再徒然寻找它。因此，经由\Person{撒罗满}说得很好：\BibleQuote{不要把你的尊荣送给外方人，不要将你的岁月交给残忍者；恐怕外方人满得你的财物，你的辛劳落入外人家中；最后，当你的皮肉和身体衰残时，你必要哀叹}\BibleRef{箴 5:9}。因为谁是我们的外方人？岂不是那些分离于天上产业之份的恶神？我们的尊荣又是什么？岂不是我们虽由泥土之身造成，却是按我们造物主的肖像和模样受造？那残忍者又是谁？岂不是那个背教的天使，他因骄傲而以死亡的痛苦打击自己，并且虽已失丧，仍不惜将死亡带给人类？因此，一个按天主的肖像和模样受造的人，若将其生命的时光用于追求恶神的享乐，就是将自己的尊荣给了外方人。若他将被赐予的生命时光按照邪恶仇敌的意愿度过，就是将自己的岁月交给了残忍者。在此处，经上又善加补充：\BibleQuote{恐怕外方人满得你的财物，你的辛劳落入外人家中}\BibleRef{箴 5:9}。因为凡借所赐的身体健康或所赐的智慧，不是致力于修德，而是致力于行恶的人，他绝不是在用自己的财物装满自己的家，而是装满外方人的住所：也就是说，他增殖了不洁之神的行径，而且如此行事，在它的奢华或骄傲中，甚至因自己的加入而增加了丧亡者的数目。此外，又善加补充：\BibleQuote{最后，当你的皮肉和身体衰残时，你必要哀叹}\BibleRef{箴 5:9}。因为大多数情况下，所赐的肉身健康被用于罪恶；但当它突然被夺走，当肉身因痛苦而衰残，当灵魂已被催促离去时，那时，早已被用于作恶却已失去的健康，才被当作为了善生而重新寻求。于是，人们哀叹他们未曾事奉天主，此时却完全无法通过事奉祂来弥补过失的损失。因此，在另一处经上记载：\BibleQuote{祂击杀了他们，他们才寻求祂}\BibleRef{咏 76:34}。\par

但另一方面，应劝告病人：在\OriginalTerm{纪律的鞭策}{scourge of discipline}责打他们时，他们应感到自己是天主的子女。因为，如果祂不打算在纠正后赐予他们嗣业，就不会费心以苦难来教育他们。为此，主通过天使对若望说：\BibleQuote{凡我所爱的，我就责备管教他}\BibleRef{默 3:19}。为此又写道：\BibleQuote{我儿，不要轻视主的管教，被祂责备时也不要灰心；因为主所爱的，祂必管教，又鞭打凡收纳的儿子}\BibleRef{希 12:5}。为此圣咏作者说：\BibleQuote{义人的苦难虽多，上主却救他脱离这一切}\BibleRef{咏 32:20}。为此，有福的约伯也在痛苦中呼喊说：\BibleQuote{我若正义，也不敢抬头，因为我充满羞辱，饱尝痛苦}\BibleRef{约 10:15}。应告诉病人：如果他们相信\OriginalTerm{天上家乡}{heavenly country}是自己的，就必须在此世如同在异乡一般忍受劳苦。因此，建造上主圣殿时，石头是在外面锤打，好使它们能在没有锤声的情况下被安置在圣殿中；也就是说，我们现在在外面被鞭策锤打，为的是以后在圣殿中，无需纪律的敲击，就能被安置在各自的位置上；这样，锤打现在能削去我们身上一切多余之物，然后\OriginalTerm{爱德的和谐}{concord of charity}独自将我们在建筑中结合在一起。应劝告病人：想想我们为获得地上的产业，对属血肉的儿女施加多么严厉的纪律鞭策。那么，\OriginalTerm{天主的改正}{divine correction}之痛苦对我们来说又算什么呢？藉此我们既得到永不丧失的产业，又避免了永久的惩罚。为此保禄说：\BibleQuote{我们曾有生身的父亲管教我们，我们尚且敬重他们；何况万灵之父，我们岂不更当服从祂而生活么？他们是在短期内随己意管教我们；但祂是为我们的益处，叫我们在祂的\OriginalTerm{圣化}{sanctification}上有分}\BibleRef{希 12:9}。\par

应劝告病人：思考\OriginalTerm{身体苦难}{bodily affliction}带给内心多大的健康，它使心灵回到认识自己，更新那常被健康抛弃的软弱记忆，好让那因自高而远离自己的精神，藉着受击打的肉身所经历的痛苦，被提醒它隶属于怎样的状况。此事在巴郎受阻于旅程时（若他曾愿意顺从神的声音）已正确示意\BibleRef{户 22:23}。巴郎正要去实现他的目的，但他胯下的牲畜却挫败了他的愿望。驴因受阻止而停步，看见了天使，而人的心灵却看不见；因为肉身常因苦难而迟缓，它藉着所承受的鞭打，向心灵指明那心灵自身虽有肉身在它之下却看不见的天主，其程度足以阻碍那渴望在此世前进（如同行路）的精神之热忱，直到它使精神认识那阻挡它道路的不可见者。因此，通过伯多禄也说得很好：\BibleQuote{那匹无言的驴竟用人的声音阻止了先知的狂妄，做了这狂妄的责备}\BibleRef{伯后 2:16}。的确，当一个人被一匹无言的牲畜责备为狂妄时，正是高傲的心灵被受苦的肉身提醒它所应保持的谦卑之善。但巴郎并未获得这责备的益处，因为他去咒骂时，改变了声音，却没有改变心意。应劝告病人：思考身体苦难是多么大的恩赐，它既洗去已犯的罪，又阻止了可能犯的罪，它借由外在的鞭打给忧苦的心灵带来忏悔的创伤。因此经上写道：\BibleQuote{伤口的青紫洗净邪恶，腹中隐秘处的鞭痕也如此}\BibleRef{箴 20:30}。因为伤口的青紫洗净邪恶，意思是鞭打的痛苦洗净已思虑或已犯下的不义。而\Quote{腹}这一称谓常被理解为心灵。心灵被称为腹，这在经上那句话中得到了教导：\BibleQuote{人的灵是上主的灯，鉴察腹中的隐秘}\BibleRef{箴 20:27}。这就如同说：神圣灵感的照临进入人的心灵时，藉着光照将它显示给它自己，而在圣神来临之前，它既能怀有恶念，却不知道如何衡量它们。于是，伤口的青紫洗净邪恶，腹中隐秘处的鞭痕也是如此，因为当我们外表受击打时，我们沉默而痛苦地被召回对罪的记忆，将过去的一切恶行重新呈现于眼前，并因外表的受苦而内心更痛悔我们所做的。由此发生了：在身体公开的伤口中，\OriginalTerm{腹中的秘密鞭痕}{secret stripe in the belly}更完全地洁净我们，因为\OriginalTerm{隐藏的悲伤之伤}{hidden wound of sorrow}治愈了恶行的不义。\par

当劝诫病者，要他们恒心保持忍耐的德行，不断思虑我们的救赎者从祂所创造的人那里忍受了何等巨大的苦难；祂承受了如此多卑劣的辱骂；祂每日从古仇手中夺回被掳的灵魂，却忍受侮辱者打在脸上的巴掌；祂以救恩之水洗净我们，却没有从无信者吐唾沫的脸上隐藏自己的面容；祂凭着自己的代祷将我们从永罚中拯救出来，却默默忍受鞭打；祂赐予我们在天使群中永恒的尊荣，却承受拳击；祂将我们从罪的刺伤中解救出来，却未拒绝将祂的头颅置于荆棘之下；祂以永恒的甘甜使我们酣醉，却在口渴时接受苦胆的苦涩；祂虽与圣父在神性上平等，却为我们朝拜圣父，而当被人戏弄朝拜时，祂默不作声；祂本身是生命，却为死者预备生命，甚至走向死亡。既然如此，若天主为善行忍受了如此巨大的苦难，人因作恶而受天主的鞭打，又怎能认为是难以承受的呢？或者，一个理智健全的人，怎能不感恩自己所受的鞭打，因为那在此世生活无罪的，也未能在无鞭打中离去呢？\par

% structure: p0094 source=h2 target=subsection reason=relative-heading-rank; base=h2

\subsection{第十三章}

\Emph{如何劝诫那些惧怕鞭打者和那些藐视鞭打者。}\par

（\Emph{Admonition} 14.）劝诫惧怕鞭打因而度无罪生活者，与劝诫那些在邪恶中如此顽固、连鞭打也无法改正者，方式不同。因为惧怕鞭打者当被告知，绝不可将现世之物视为重大而渴求，因为恶人也拥有它们；也绝不该将现世的苦难视为无法忍受而逃避，因为他们并非不知道，善人大多也会遭受这些。当劝诫他们，如果他们真想脱离邪恶，就应畏惧永罚；但也不要停留在对惩罚的恐惧中，而应借爱德的养育成长到爱的恩宠。因为经上记载：\BibleQuote{圆满的爱德驱散恐惧}\BibleRef{若一 4:18}。又说：\BibleQuote{你们没有领受奴役之神，再次陷于恐惧，而是领受了义子的圣神，我们借此呼叫：\InnerQuote{阿爸，父呀！}}\BibleRef{罗 8:15}。因此同一位导师又说：\BibleQuote{主的圣神在哪里，那里就有自由}\BibleRef{格后 3:17}。所以，如果对惩罚的恐惧仍然阻止人作恶，那么那恐惧者的灵魂的确没有享有圣神中的自由。因为，他若不怕惩罚，无疑会犯罪。因此，被恐惧的束缚所捆绑的心，不认识自由的恩宠。因为善应当为自身而被爱，而非出于刑罚的强迫而追求。因为一个人行善，是出于害怕折磨的痛苦，那么他所惧怕的，正是所愿其不存在的，以便他能大胆地做不法之事。由此可知，清白无辜在天主眼中是这样丧失的，因为在天主眼中，恶念就是罪。\par

但是，另一方面，那些连鞭打都不能制止他们行不义的人，要用更严厉的斥责来打击，因为他们已变得麻木不仁而更加顽固。因为一般来说，他们应当被轻看，却不是出于轻看；被绝望，却不是出于绝望；也就是说，所表现的绝望是为了使他们感到恐惧，随后而来的劝告则带他们回到希望。因此，应当严厉地向他们陈明天主的审判，好使他们因想到永远的报应而恢复自知。因为他们该听，在他们身上应验了所记载的：\BibleQuote{你若在臼中捣碎愚人，如同用杵捣碎大麦，他的愚妄仍不会离开他。}\BibleRef{箴 27:22}。先知向主抱怨说：\BibleQuote{你击打他们，他们却不肯受管教。}\BibleRef{耶 5:3}。因此主说：\BibleQuote{我击杀、毁灭了这百姓，他们仍不回转离开他们的道路。}\BibleRef{依 9:13}。主又说：\BibleQuote{这百姓不归向那击打他们的。}\BibleRef{耶 15:6}。因此先知通过鞭打者的声音抱怨说：\BibleQuote{我们照料了巴比伦，她却没有痊愈。}\BibleRef{耶 51:9}。因为当处于作恶中混乱的心灵听见责斥的话语、感受责斥的鞭打，却仍不屑于回到救恩的正路时，巴比伦就受到了照料，却仍未康复。因此主责备被掳但尚未从邪恶中悔改的以色列子民说：\BibleQuote{以色列家在我眼中成了渣滓；他们都是炉中的铜、锡、铁、铅。}\BibleRef{厄 22:18}；仿佛在明说：我本想用磨难的烈火来炼净他们，希望他们能成为银或金；但他们在我面前在炉中变成了铜、锡、铁、铅，因为即使在磨难中，他们也没有进德，反而增长恶习。的确，铜被敲打时，发出的响声比所有其他金属都大。因此，人受打击时若发出抱怨的声音，就在炉中变成了铜。但锡经过人为加工，有银的假象。因此，在磨难中仍不脱离虚伪恶习的人，就变成了炉中的锡。再者，图谋邻舍性命的人使用铁。因此，炉中的铁就是在磨难中不失去害人之恶意的人。铅也是最重的金属。因此，人即使在磨难中仍不超越地上的欲望，就变成了炉中的铅。为此又记载着：\BibleQuote{她劳苦太多，她的极大的锈仍不出去，甚至用火也不出去。}\BibleRef{厄 24:12}。因为祂给我们降下磨难的烈火，为要清除我们恶习的锈；但我们连在火中也未失去锈，因为即使在鞭打中我们也不缺恶习。因此先知又说：\BibleQuote{熔炼者白费了力气，他们的邪恶并未耗尽。}\BibleRef{耶 6:29}。\par

然而，要知道，有时当他们在鞭打下的刚硬中仍不改正时，他们要用温和的劝告来抚慰。因为那些不被苦刑所改正的人，有时会被温柔的安抚从邪恶行为中制止。因为通常病人，烈性药未能治愈的，却被温水恢复了健康；有些伤口不能通过切开治愈，却被油敷治愈；坚硬的钻石根本不能用钢铁切开，却被温和的山羊血软化。\par

% structure: p0099 source=h2 target=subsection reason=relative-heading-rank; base=h2

\subsection{第十四章}

\Emph{如何劝告沉默寡言者和好言多语者。}\par

（\Emph{劝谕}十五。）缄默过度之人与多言多语之人，当以不同方式予以劝诫。因为对于缄默过度之人，应暗示他们：不谨慎地回避某些恶习时，不知不觉便会陷入更严重的恶习。往往因过分约束口舌，内心反而遭受更严重的多言之苦；心思因受不智沉默的强力压制而愈加沸腾。他们多半会越发泛滥，只因自以为在挑剔者面前不被看见而更加安心。由此，人的心思有时会骄傲自大，轻视那些他所听见说话的人，视他们为软弱。当他紧闭肉体的口时，却未察觉因骄傲而向恶习敞开到何种程度。他压抑口舌，却抬高心思；很少考虑自己的恶，在心中越自由地指责众人，也越隐秘。因此，应劝诫缄默过度之人：不仅殷勤学习在外表上应显示为何等人，也当在内里显示为何等人；他们更当惧怕思想暗中的审判，而非邻人对言辞的责备。因为经上记载：\BibleQuote{我儿，你要留意我的智慧，侧耳听我的明哲，为使你保守你的思想。}\BibleRef{箴 5:1}诚然，没有比人心更飘忽不定的，它每因恶念而溜走，便背离我们。因此圣咏作者说：\BibleQuote{我的心已离弃了我。}\BibleRef{咏 39:12}因此，当他回归自己时说：\BibleQuote{祢的仆人找到了他的心向祢祈祷。}\BibleRef{撒下 7:27}所以，当思想受到看守时，那颗惯于飞逝的心便被找着了。此外，缄默过度之人遭受不义时，多半因不言所受之苦而疼感更为敏锐。因为若口舌能平静倾诉所招致的烦恼，痛苦便从意识中流走。封闭的伤口更加疼痛；因为当内部炽热的腐败被排出时，伤口便敞开以便医治。因此，那些缄默超过合宜限度的人应当知道：在沉默忍受烦恼时，切勿加剧痛苦的猛烈。因为应劝告他们：若爱邻人如同自己，绝不可对他们隐瞒自己有理责难他们的理由。因为声音的良药对双方的健康有协同效果：加害者一方，恶行受制止；受害者一方，因敞开伤口而缓解痛苦的炽热。因为那些注意邻人恶行却默然不语的人，犹如对所见伤口不予医治，成为死亡的制造者，因他们本可治愈却不去治愈那毒害。因此，口舌应谨慎约束，而非紧紧捆住。因为经上记载：\BibleQuote{智者静默直到时候到来。}\BibleRef{训 20:7}当然，是为了当他觉得适宜时，放弃沉默的节制，藉说出适当的话投身于有益的服务。又记载：\BibleQuote{静默有时，言谈有时。}\BibleRef{训 3:7}诚然，当谨慎掂量变换的时机，以免口舌当受约束时却无益任其放纵，或当可有益发言时却懒惰自制。圣咏作者深明此理，说：\BibleQuote{上主，求祢在我的口上派一个守卫，在我的唇上设一道门。}\BibleRef{咏 140:3}他求的不是在唇上设一道墙，而是一道门：即能开能关之物。因此我们亦当谨慎学习：声音合宜适时地开口，合宜适时地缄默。\par

但另一方面，应劝诫多言多语之人：当他们沉溺于众多言语时，要警觉从何等正直状态中堕落。因为人心似水，被封闭时便向高处汇集，因它寻求当初下降的高处；一旦放纵，便因无益地散于最低处而坠落。因它从沉默的节制中散逸于多少多余的话，便如被多少水流从自身引开。因此，它不能回归内心认识自己，因为被多言四散，便将自己排除于内省之秘处。但它将全身暴露于埋伏的仇敌的创伤前，因它不设任何警醒的防卫。因此经上记载：\BibleQuote{如一座无围墙的敞开的城，就是那在言语上不能控制自己精神的人。}\BibleRef{箴 25:28}因为缺少沉默的墙，心灵之城便向仇敌的标枪敞开；当它以言语将自己抛出自身时，便显露自己向敌人。敌人越不费力便能战胜它，正如那被征服的心灵本身越费力地以多言与自己作战。\par

此外，由于懒散的心灵大多逐渐滑向堕落，当我们疏于防备闲言时，便会走向有害的话语；起初我们乐于谈论他人的事；随后舌头以诋毁咀嚼我们所谈论之人的生活；最终甚至爆发为公开的诽谤。因此，刺人的荆棘被播下，争吵兴起，仇恨的火炬被点燃，心灵的平安被熄灭。由此，通过撒罗满说得很好：\BibleQuote{放水的人，是争端的泉源}\BibleRef{箴 17:14}。因为放水就是放任舌头滔滔不绝地说话。因此，从另一面说，也有美好的话：\BibleQuote{人口中的言语，如同深水}\BibleRef{箴 18:4}。所以，那放水的人就是争端的泉源，因为不勒住自己舌头的人破坏了和睦。因此，另一方面经上记载：\BibleQuote{使愚人静默的人，平息了敌意}\BibleRef{箴 26:10}。再者，那好说话的人不能持守正义的直道，先知为此作证说：\BibleQuote{多言的人，在地上必不得正导}\BibleRef{咏 140:11}。因此撒罗满又说：\BibleQuote{多言难免有罪}\BibleRef{箴 10:19}。因此依撒意亚说：\BibleQuote{正义的修养是静默}\BibleRef{依 32:17}，指明：当无节制地多言时，心灵的正义就被荒废了。因此雅各伯说：\BibleQuote{若有人自以为虔诚，却不勒住自己的舌头，反而欺骗自己的心，这人的虔诚是虚妄的}\BibleRef{雅 1:26}。他又说：\BibleQuote{每个人要敏于听，缓于说}\BibleRef{雅 1:19}。再者，他描述舌头的力量时补充说：\BibleQuote{无法制伏的恶，满含致死的毒}\BibleRef{雅 3:8}。因此真理本身劝诫我们说：\BibleQuote{人所说的每句闲话，在审判之日都要交账}\BibleRef{玛 12:36}。的确，凡是缺乏正当必要的理由或敬虔有益的意图的话语，都是闲话。如果为闲谈也要交账，那么让我们仔细思量，多言中还有有害词语的罪，将有何等的惩罚在等待。\par

% structure: p0104 source=h2 target=subsection reason=relative-heading-rank; base=h2

\subsection{第十五章}

\Emph{如何劝诫懒惰和急躁的人。}\par

（\Emph{劝诫}16.）懒惰和急躁的人应受不同的劝诫。因为前者应被劝服，不要因拖延而失去他们应当行的善；但后者应受劝诫，以免因不假思索的急躁而抢先于行善的时刻，从而改变了其功绩的性质。因此，应向懒惰者指出，当我们不肯在适当的时候做我们能做的事时，不久，当我们想做时，我们却不能了。因为心灵的懈怠，当它没有被适切的热忱点燃时，就会被一种悄然滋生的麻木所切断，从而丧失对一切善事的渴望。因此，通过撒罗满明确地说：\BibleQuote{懒惰使人沉睡}\BibleRef{箴 19:15}。因为懒惰人仿佛醒着，因为他感觉正确，但无所事事而变得麻木；但懒惰被称为使人沉睡，因为当行善的热忱中断时，正确感觉的警觉也逐渐丧失。在同一处又恰当补充：\BibleQuote{懈怠的灵魂必受饥饿}\BibleRef{箴 19:15}。因为，它不把自己束向更高的事物，就放任自己在低下的欲望中不受照顾地游荡；并且，由于没有高尚目标的活力来约束，就遭受低等贪婪的饥渴之苦，并且，由于它忽略了以纪律约束自己，就在对快乐的渴望中更加饥饿地分散自己。因此，同一位撒罗满再次写道：\BibleQuote{懒惰人完全在欲望之中}\BibleRef{箴 21:26}。因此，在真理本身的宣讲中，\BibleRef{玛 12:44}，当一个魔鬼出去时，房屋确实是干净的；但是，当空着时，它就被他带着更多的魔鬼回来占据。因为懒惰人大多在忽略做必要之事的同时，却给自己摆出一些困难的事，而对其他事又不假思索地害怕；因此，仿佛找到了可以合理惧怕的东西，他满足于自己有理由继续保持麻木。通过撒罗满对他说得恰当：\BibleQuote{懒惰人因寒冷不耕种，因此夏天必乞讨而不得}\BibleRef{箴 20:4}。的确，懒惰人为寒冷而不耕种，当他找到借口不做他应当做的善事时。懒惰人为寒冷而不耕种，当他害怕微小的反对之恶而放弃最重要之事时。再者，说得很好：\BibleQuote{他夏天必须乞讨，也得不到。}因为现在不在善工中勤劳的人，夏天将乞讨而一无所获，因为当审判的烈日显现时，他将徒然恳求进入天国。同一位撒罗满对他也说得好：\BibleQuote{观察风的不播种，看云的不收割}\BibleRef{训 11:4}。因为风所表达的，不是恶灵的诱惑吗？被风吹动的云所代表的，不是恶人的反对吗？即云被风驱动，因为恶人被不洁之灵的气息激起。因此，那观察风的不播种，看云的不收割，因为凡害怕恶灵诱惑、凡害怕恶人迫害，而如今不播种善工种子的人，将来也不能收获成捆的神圣报偿。\par

但另一方面，急躁之人，因抢先于善行的时间，反而歪曲了善行的功绩，并且常常陷入邪恶，因为他们完全不能辨别何为善。这样的人在行事时全然不看自己在做什么，往往在事成之后才意识到自己本不该做。对于这种人，在撒罗满的书中以学习者的口吻说得很好：\Quote{\BibleQuote{我儿，不要在没有计谋的情况下做任何事，事后你必不后悔}} \BibleRef{德 32:24}。又说：\Quote{\BibleQuote{让你的眼睑走在脚步之前}} \BibleRef{箴 4:25}。因为我们的眼睑走在脚步之前，就是正确的谋略先行于我们的行动。因为一个人若忽略用思考预先展望他将要行的事，就等于闭着眼迈步；他向前走并完成旅程，却没有因前瞻而超越自己；因此他更快跌倒，因为他没有通过谋略的眼睑留意应将行动的脚踏向何处。\par

% structure: p0108 source=h2 target=subsection reason=relative-heading-rank; base=h2

\subsection{第十六章}

\Emph{温和者与易怒者应如何劝诫。}\par

(\Emph{劝诫}第十七) 温和者与易怒者应受到不同的劝诫。因为有时温和者在掌权时，会遭受怠惰的麻木，这是一种近似的性情，仿佛紧邻着。而且大多由于过于温和的松弛，他们将严格的力量软化到不必要的程度。但另一方面，易怒者因被愤怒冲动卷入心灵的狂乱，破坏了宁静的平和，从而扰乱了他们属下之人的生活。因为当暴怒驱使他们盲目冲动时，他们不知道自己在愤怒中做了什么，也不知道自己在愤怒中承受了什么。但有时，更严重的是，他们将愤怒的刺视为义德的热情。而当恶行被当作德行时，罪过便无所畏惧地堆积起来。于是，温和者常在无所作为的懒惰中变得麻木；易怒者常被正直的热情所欺骗。因此，前者的德行不知不觉地附着了恶行，而后者的恶行却看似炽热的德行。因此，前者应受劝诫远离紧邻自己的事物，后者应小心留意自身之内的事物；前者要辨别自己所缺乏的，后者要辨别自己所拥有的。让温和者拥抱谨慎，让易怒者摒弃烦乱。温和者应受劝诫，要努力拥有义德的热情；易怒者应受劝诫，要在他们认为拥有的热情上加上温和。因为为此原因，圣神以鸽子和火的形式显现给我们；即，凡祂所充满的人，祂都使他们显出鸽子的温和纯朴，并燃烧着热情之火。\par

那么，无论谁如果在温和的宁静中放弃了热情的热忱，或者在热情的热忱中失去了温和的德行，就绝非充满圣神。若我们引用保禄的权威，或许能更好地说明这一点。他对两位门徒同样赋予爱德，却为他们提供不同的宣讲助力。因为在劝诫弟茂德时，他说：\Quote{\BibleQuote{你要责备、恳求、劝戒，以百般的忍耐和教训}} \BibleRef{弟后 4:2}。他也劝诫弟铎说：\Quote{\BibleQuote{这些事你要讲论、劝勉、责备，以各样权柄}} \BibleRef{铎 2:15}。他为何如此巧妙地分配他的教导，在传授时向一人推荐权柄，向另一人推荐忍耐？除非他看出弟铎的性情更为温和，弟茂德的性情更为热切。前者他以热诚的热情点燃，后者他以忍耐的温柔节制。对一人他补足所缺，对另一人他削减过多。一人他竭力用马刺催逼，另一人用缰绳勒住。因为那位掌管教会的大农夫，浇灌一些嫩芽使之生长，又修剪另一些他看到生长过快的嫩芽；免得它们因不生长而不结果实，或因生长过度而失去可能结出的果实。然而，以热情为伪装悄悄潜入的愤怒，与毫无义德借口就搅乱烦乱心灵的愤怒，相去甚远。前者在不恰当的情况下过度延伸，而后者总是在不该燃起的地方燃起。应当知道，易怒者与急躁者在此不同：后者不能忍受别人加给他们的事，而前者则自己带来需要忍受的事。因为易怒者常常追逐躲避他们的人，挑起争端的由头，以争斗的劳苦为乐；然而，我们若在他们愤怒的骚动中避开他们，反而更能纠正他们。因为当他们烦乱时，不知道我们对他们说什么；但返身自省时，他们因曾被更平静地容忍而感到羞愧，从而更自由地接受劝勉之言。因此，当纳巴耳醉酒时，阿彼盖耳值得赞扬地对他所犯的过错保持沉默；但当他酒醒后，她又同样值得赞扬地告知他这事 \BibleRef{撒上 25:37}。\par

但是，当那些易怒者如此攻击他人，以至于无法彻底避开时，不应公开斥责他们，而应谨慎地、带有某种敬重的克制来打击他们。如果我们举出\Person{阿贝乃尔}所作所为，就能更好地说明这一点。因为，当\Person{阿撒耳}以不顾后果的急躁暴力攻击他时，经上记载：\BibleQuote{\Person{阿贝乃尔}对\Person{阿撒耳}说：\InnerQuote{你转开，不要追赶我，免得我被迫把你击倒在地。}但他轻视不听，不肯转开。于是\Person{阿贝乃尔}用枪的钝端刺入他的小腹，将他穿透，他就死了。}\newline{}\BibleRef{撒下 2:22}。因为\Person{阿撒耳}代表了什么样的人呢？不就是那些被愤怒猛烈攫住并冲昏头脑的人吗？这类人在同样的愤怒发作中，应按照其疯狂狂奔的程度而谨慎回避。因此，\Person{阿贝乃尔}——在我们的语言中被称为\Quote{父亲的灯笼}——逃跑了；因为当教师的口舌——它表明天主的上智之光——看到任何人的心灵被愤怒的峭壁裹挟而前时，它就像是不愿击打追赶者一样，克制着不向愤怒者抛出言语的箭矢。但是，当易怒者不肯通过任何考虑平息自己，反而像\Person{阿撒耳}那样，不停追赶、不停疯狂时，那么，那些努力压制这些狂怒者的人，绝对不可以自己陷入愤怒，而要展示出一切可能的平静；然而又要巧妙地施加影响，通过侧击刺穿愤怒者的心。因此，当\Person{阿贝乃尔}站住对抗追赶者时，他不是用直刺，而是用枪的钝端刺伤了他。因为用枪尖刺就是用公开斥责的攻击来反对；而用枪的钝端击打追赶者，就是平静地用某些打击来触动愤怒者，并仿佛因饶恕而战胜他。\Person{阿撒耳}随即倒下，因为激动的心灵，当感受到自己被饶恕，却又被平静的回答从内在触动时，会立刻从自己所升到的高处坠落。因此，那些在温和的打击下从自己热浪的冲击中弹回的人，仿佛未受刀剑而死去。\par

% structure: p0113 source=h2 target=subsection reason=relative-heading-rank; base=h2

\subsection{第十七章}

\Emph{如何劝告谦卑者与傲慢者。}\par

（\Emph{劝告}18.）谦卑者与傲慢者应受不同的劝告。对前者，应暗示他们所期望拥有的那种卓越何等真实；对后者，则应表明那暂时的荣耀何其虚无，即使他们拥抱它，也把握不住。让谦卑者听到他们所渴望的是永恒的，他们所轻视的是暂时的；让傲慢者听到他们所追求的是暂时的，他们所失去的是永恒的。让谦卑者从真理的权威之声中听到：\BibleQuote{凡自谦者，必被高举。}\newline{}\BibleRef{路 18:14}。让傲慢者听到：\BibleQuote{凡自高者，必被贬抑。}\newline{}\BibleRef{路 18:14}。让谦卑者听到：\BibleQuote{谦逊在尊荣之前；}\newline{}让傲慢者听到：\BibleQuote{心灵高傲在跌倒之前。}\newline{}\BibleRef{箴 15:33}。让谦卑者听到：\BibleQuote{我垂顾谁呢？岂非垂顾那谦卑、安静、并对我的话战栗的人吗？}\newline{}\BibleRef{依 66:2}。让傲慢者听到：\BibleQuote{为什么尘土和灰烬会骄傲？}\newline{}\BibleRef{德 10:9}。让谦卑者听到：\BibleQuote{天主垂顾谦卑的事。}\newline{}让傲慢者听到：\BibleQuote{高大的事，他远远知道。}\newline{}\BibleRef{咏 137:6}。让谦卑者听到：\BibleQuote{人子来不是受服事，而是服事人。}\newline{}\BibleRef{玛 20:28}；让傲慢者听到：\BibleQuote{一切罪恶的开端是骄傲。}\newline{}\BibleRef{德 10:13}。让谦卑者听到：\BibleQuote{我们的救赎者自谦自卑，服从至死。}\newline{}\BibleRef{斐 2:8}；让傲慢者听到关于他们首领的记载：\BibleQuote{他是所有骄傲之子的君王。}\newline{}\BibleRef{约 41:25}。因此，魔鬼的骄傲成为我们丧亡的缘由，而天主的谦卑成为我们救赎的论证。因为我们的仇敌，被造于万物之中，却想显现在万物之上；而我们的救赎主，伟大超乎万物之上，却屈尊成为万物中的微小者。\par

那么，让谦卑者明白：当他们自卑时，便上升到天主的相似；让傲慢者明白：当他们自高时，便堕入叛天使的模仿。那么，还有什么比傲慢更卑贱呢？它将自己伸展到自身之上，超越了真正高超的尺度。又有什么比谦卑更高超呢？它将自己降至最低处，却与那居于最高之上的造物主相连。然而，在这些情况中，还有另一件事应当仔细考虑：有些人常被虚假的谦卑表现所欺骗，而有些人则因对自己傲慢的无知而被蒙蔽。因为通常，一些自以为谦卑的人夹杂着不应当对人怀有的恐惧；而傲慢者却常常伴随着直言不讳的主张。当需要谴责某些恶习时，前者因恐惧而保持沉默，却自以为是出于谦卑的沉默；后者出于傲慢的不耐烦而发言，却相信自己是出于正直的自由而言。前者的懦弱缺点在谦卑的外表下阻止他们谴责错误；后者的傲慢放纵无羁，在自由的伪装下推动他们去谴责不应谴责的事，或以过度的方式去谴责。因此，应当劝告傲慢者不要过度自由，也要劝告谦卑者不要过度顺从；以免前者将正义的辩护变成骄傲的展示，或后者在过度顺从人的研究时，甚至被迫尊重人的恶习。\par

然而，应当考虑，我们责备骄傲者时，若在责备中掺入一些赞扬的香膏，通常会更有效。因为我们在责备中应提及他们的一些优点，或至少提及那些他们可能具备、尽管实际上没有的优点；然后，当我们先肯定了他们令我们喜悦的优点，使他们心平气和地听从以后，再最后割除那些令我们不悦的缺点。对于未驯服的马，我们也是先用轻柔的手触碰，以便之后甚至用鞭子驯服它们。蜂蜜的甜味被加入苦药中，以免对健康有益的苦味在品尝时感到过于苦涩；但味觉被甜味欺骗的同时，致命的体液被苦味排出。所以，对于骄傲者，我们责备的开端应用赞扬调和，以便他们在接受所爱的称赞之时，也能接受所憎恨的责备。\par

此外，我们在多数情况下能更有效地劝服骄傲者向善，如果我们说他们的改过对我们比对她们自己更有益处；如果我们请求他们改过，是为了赐予我们而非他们自己。因为骄傲很容易向善屈服，如果它的屈服被认为对他人也有益处。为此，梅瑟在云柱引导下、蒙天主指引穿越旷野时，想要将他的亲属曷巴布从与外邦人的往来中引开，使他服从全能天主的管辖，就说：\Quote{我们正往你要赐给我们的地方去；与我们同行，我们必善待你，因为上主对以色列说了好话。\newline{}当对方回答他说：我不与你同行，我要回到我本乡本土去。他立刻加上说：求你不要离开我们，因为你知道我们在旷野中应如何安营，你要作我们的向导。} \BibleRef{户 10:29} 然而，梅瑟心里并未因不知道路而窘迫，因为他与天主的交往已在他内开启了预言的知识；外在有云柱前行引导，内在有他与天主殷勤交谈中熟悉的话语教导他一切。但事实上，作为有远见的人，他对一个骄傲的听者说话，他寻求帮助以便给予帮助；他请求一位路上的向导，以便自己能成为他的生命向导。他这样做，是为了让骄傲的听者因以为自己是必需的而更加留心那劝他向善的声音，并且，因为他相信自己是他劝诫者的向导，他就能俯首听从劝诫的话语。\par

% structure: p0119 source=h2 target=subsection reason=relative-heading-rank; base=h2

\subsection{第十八章}

\Emph{如何劝诫固执者和反复者}\par

（劝诫十九）劝诫固执者和反复者应有所不同。前者当被告知，他们过于高估自己，因此不听从他人的劝告；而后者当被明白，他们太过于轻看和忽视自己，因此在不同时刻被自己的判断所左右。前者当被告知，除非他们自以为比别人更好，否则绝不可能轻视众人的劝告而只重视自己的思虑；后者当被告知，如果他们稍微留意自己的本相，变化的微风绝不会使他们卷进这么多不同的方面。对前者，保禄曾说：\BibleQuote{不要自作聪明} \BibleRef{罗 12:16}；对后者，应听到这句话：\BibleQuote{使我们不再作小孩子，被各种教义之风所飘荡} \BibleRef{弗 4:14}。关于前者，撒罗满说：\BibleQuote{他们必自食其果，饱受自己的计谋之苦} \BibleRef{箴 1:31}；关于后者，他又写道：\BibleQuote{愚人的心必改变} \BibleRef{箴 15:7}。因为智者的心始终如一，因为他安于善的劝勉，始终在善行中坚定不移。但愚人的心变化无常，因为由于善变，他永远不会保持原状。既然有些恶行如同出于自身而产生其他恶行，同样也从其他恶行产生，那么我们必须明白：当我们通过责备从滋生出它们的苦源处使其枯干时，就能更好地清除它们。因为固执生于骄傲，反复生于轻浮。\par

因此，当劝诫固执者，使他们承认自己思想的傲慢，并努力战胜自己；以免他们因不屑于在自身之外接受他人的正确劝告，反而在内里被骄傲所俘虏。他们当明智地观察，人子（其旨意始终与圣父一致）为了赐予我们顺从己意的榜样，如何说：\BibleQuote{我不寻求自己的旨意，只寻求派遣我来者的旨意} \BibleRef{若 5:30}。而且，为了更加赞许这一美德的恩宠，他预先宣告在最后的审判中也将持守此道，说：\BibleQuote{我凭自己不能做什么，我怎样听见，就怎样审判} \BibleRef{若 5:30}。那么，一个人有何良心能不屑于赞同他人的旨意呢？当人子、天主子显示他权能的荣耀时，尚且证明他凭自己什么都不审判？\par

但另一方面，易变者应受劝诫，以庄重坚固其心灵。因为他们先斩断心中轻浮之根，便干涸了自身善变的根源；正如先预备稳固之地以奠基，方能筑起坚固的建筑。除非事先防范轻浮之心，否则思想的反复无常绝难克制。保禄声称自己免于此，他说：\BibleQuote{\Emph{我曾用过轻浮吗？}或是我所打算的，是按肉体打算，在我有是非吗？}（\BibleRef{格后 1:17}）仿佛明说：我之所以不为任何变易之风所动，正是因我不向轻浮之恶屈服。\par

% structure: p0124 source=h2 target=subsection reason=relative-heading-rank; base=h2

\subsection{第十九章}

\Emph{如何劝诫饮食无度者和饮食有度者。}\par

（劝诫第二十则）贪食者与节制者应受不同的劝诫。因为前者伴随多言、轻浮之举和邪淫；而后者则常有急躁之罪，甚至骄傲之罪。若非无节制的多言困扰贪食者，那位每日奢华享乐的富翁就不会在舌头上烧得更厉害，说：\BibleQuote{亚巴郎父亲，可怜我吧！请打发拉匝禄来，用指尖蘸点水，凉凉我的舌头，因为我在这火焰中备受煎熬。}（\BibleRef{路 16:24}）这些话明明显示他因每日宴饮而常常犯罪于舌头，因为他全身燃烧，却特别要求凉却舌头。再次，神圣权威证明轻浮之举紧随着贪食，说：\BibleQuote{百姓坐下吃喝，起来玩耍。}（\BibleRef{出 32:6}）多数时候贪食甚至引我们走向邪淫，因为肚腹因饱胀而鼓胀，情欲的刺便被激起。因此，那狡猾的仇敌既以对苹果的贪欲开启了第一个人的官能，又将其束缚在罪恶的罗网中，天主的声音对他说：\BibleQuote{你要用肚腹爬行，用胸膛爬行。}（\BibleRef{创 3:14}）仿佛明白地说：你将在思想和肚腹中统治人心。先知证实邪淫随贪食而来，在揭发隐密之事时又谈及公开之事，说：\BibleQuote{厨师长拆毁了耶路撒冷的城墙。}（\BibleRef{耶 39:9}）厨师长便是肚腹，厨师们以极大细心侍奉它，为使其享用地被佳肴填满。但耶路撒冷的城墙是灵魂的德行，提升对天上平安的渴望。因此，厨师长拆毁了耶路撒冷的城墙，因为当肚腹以贪食膨胀时，灵魂的德行便因邪淫而被摧毁。\par

另一方面，若非急躁通常将节制者从安宁的怀抱中摇落，伯多禄在说\BibleQuote{你们要分外加添：在信德上加添德行，在德行上加添知识，在知识上加添节制}（\BibleRef{伯后 1:5}）之后，本不会立即警惕地加上\BibleQuote{在节制上加添忍耐。}因为他预见节制者将缺少他所劝诫的忍耐。若非有时骄傲之罪刺透节制者的思念，保禄本不会说：\BibleQuote{那不吃的，不要判断那吃的。}（\BibleRef{罗 14:3}）而且在向其他人说话时，他针对那些以节制之德夸耀者的格言，又加上：\BibleQuote{这些事在迷信和谦卑中确有智慧的表象，而且苦待身体，但对于满足肉欲毫无价值。}（\BibleRef{哥 2:25}）这里应注意，这位杰出的讲道者在其论证中将谦卑的表象与迷信相联，因为当身体因节制而过分受折磨时，外在表现是谦卑，但正因这谦卑，内心却充满严重的骄傲。除非有时心灵因节制之德而自高自大，那位骄傲的法利塞人就不会刻意将其列为自己大功劳之一，说：\BibleQuote{我一周禁食两次。}（\BibleRef{路 18:12}）\par

因此，贪食者应受劝诫，以免在纵情美味时，被邪淫之剑刺穿；并要明白多少多言、多少轻浮之心藉着饮食潜伏在他们身上；免得在温柔侍奉肚腹时，被残酷地束缚在罪恶的罗网中。因为我们越是无节制地放纵，在伸手取食时重演第一父母的堕落，就离我们的第二父母越远。但另一方面，节制者应受劝诫，要时时警惕地留意，以免在逃避贪食之恶时，仿佛从德行中生出更恶的恶，即折磨肉体时，精神却迸发为急躁；于是当肉体被制伏时，精神却被愤怒征服，从而无德可言。此外，有时节制者的心虽压制愤怒，却仿佛被外来的喜悦侵入而败坏，失去节制的一切益处，因为未能防范属灵的罪恶。因此先知恰当地说：\BibleQuote{你们禁食的日子，仍求自己喜悦。}（\BibleRef{依 58:3}）稍后又说：\BibleQuote{你们禁食，却吵闹争竞，用拳头打人。}（\BibleRef{依 58:3}）因为喜悦属于享乐，拳头属于愤怒。所以，若心灵放纵于混乱情感，被恶习驱散，那么身体徒然被节制所折磨。再者，他们应受劝诫，在坚持节制毫不松懈时，不要以为这在隐藏的审判者面前有超绝的德行；免得一旦以为功劳甚大，心便高傲起来。因此先知说：\BibleQuote{难道这就是我所选的禁食？而是你们要分饼给饥饿的人，将穷苦和漂泊的人领到家中。}（\BibleRef{依 58:5}）\par

关于此事，应思量禁食的德行何其微小，因为它若非与其它德行相联，便不被赞许。因此岳厄尔说：\BibleQuote{祝圣斋戒}。因为祝圣斋戒，就是借着加添其它善行，表明肉身的禁食堪当奉献给天主。要劝勉守斋者：他们将对自己克制的食物施与贫困者时，便是向天主献上悦乐天主的禁食。因为我们当明智地注意主借先知所谴责的话：\BibleQuote{你们在这七十年中，在五月和七月守斋哀哭，你们岂是向我守斋？你们吃，你们喝，岂不是为你们自己吃，为你们自己喝？}（\BibleRef{匝 7:5}）因为人若暂时克制肚腹，却不将食物给予需要者，而是保留日后供自己肚腹，那么他并非为天主守斋，而是为自己守斋。\par

因此，为使贪食的食欲不致松懈前一类人的警惕，也不致使受克苦的肉体因骄傲绊倒后一类人，让前一类人从真理的口中听到这话：\BibleQuote{你们应当谨慎，免得你们的心为宴饮沉醉及人生的挂虑所累}（\BibleRef{路 21:34}）。同一地方又加上了有益的警告：\BibleQuote{那日子就忽然临到你们，如同罗网，因为那日子要临到全地面上一切居住的人}（\BibleRef{路 21:35}）。让后一类人听到：\BibleQuote{不是入口的污秽人，而是出口的污秽人}（\BibleRef{玛 15:11}）。让前一类人听到：\BibleQuote{食物是为肚腹，肚腹是为食物；但天主必要使这两者同归于尽}（\BibleRef{格前 6:13}）。又说：\BibleQuote{不可在狂宴醉酒中}（\BibleRef{罗 13:13}）。又说：\BibleQuote{食物不能使我们悦乐天主}（\BibleRef{格前 8:8}）。让后一类人听到：\BibleQuote{在纯洁的人，一切都是纯洁的；但在败坏和不信的人，没有一样是纯洁的}（\BibleRef{铎 1:15}）。让前一类人听到：\BibleQuote{他们的天主是自己的肚腹，他们的光荣在于他们的羞耻}（\BibleRef{斐 3:19}）。让后一类人听到：\BibleQuote{有些人将背弃信德}，不久后又说：\BibleQuote{禁止嫁娶，禁戒食物，就是天主所造，叫相信且认识真理的人，以感恩的心领受的}（\BibleRef{弟前 4:1}）。让前一类人听到：\BibleQuote{无论是吃肉，是喝酒，或做任何使兄弟跌倒的事，最好都不做}（\BibleRef{罗 14:21}）。让后一类人听到：\BibleQuote{为了你的胃和你的屡次疾病，用点酒吧}（\BibleRef{弟前 5:23}）。这样，前一类人可以学会不过分贪求肉身的食物，后一类人也不敢谴责天主所造的、他们不贪恋的食物。\par

% structure: p0131 source=h2 target=subsection reason=relative-heading-rank; base=h2

\subsection{第二十章}

\Emph{如何劝勉那些施舍自己财物的，以及那些夺取他人财物的。}\par

（\Emph{劝诫} 21.）对已经慷慨施舍自己财物的人，与仍然贪图他人财物的人，应以不同方式加以规劝。因为对已经慷慨施舍自己财物的人，应规劝他们不要因自己将属世之物施予他人而心高气傲；不要因看到他人由自己供养而自以为比他人更优越。因为尘世家宅的主在分配仆人的职责与职务时，指定一些人管理，另一些人受管理。祂吩咐前者供应其余人所必需的，后者则领受从他人所获得的。然而，往往正是那些管理者得罪主人，而受管理者却得家主欢心。分施者招致愤怒，靠他人分施而生存者却无过犯。因此，对已经慷慨施舍自己财物的人，应规劝他们承认自己是由天上的主委派为暂时供应的分施者，并因明了所分施的并非自己的财物，而更加谦卑地施与。当他们想到自己被指派为那些领受者服务时，绝不可让虚荣心得意洋洋，反要让恐惧压抑他们。因此，他们也有必要忧虑，恐怕不配地分施了所委托的；恐怕将一些本不该花费的给了某些人，或在该花费的某些人身上什么都没给；恐怕在该少给的人身上给了很多，或在该多给的人身上给了很少；恐怕因草率而毫无益处地分散了所施的；恐怕因迟缓而恶毒地折磨请求者；恐怕图谋回报的念头潜入；恐怕对短暂赞美的渴求扑灭了施舍的光明；恐怕伴随的愁容玷污了所献的礼物；恐怕在妥善献上礼物后，心情过分欢悦；恐怕当一切做得妥当后，他们将功劳归于自己，因而在完成一切后立刻丧失一切。为使不将慷慨的德能归于自己，让他们听所记载的：\BibleQuote{若有人服事，就当按着天主所赐的力量服事}\BibleRef{伯前 4:11}。为使不因所施的恩惠而过分喜悦，让他们听所记载的：\BibleQuote{你们做完了一切所吩咐的，只当说：我们是无用的仆人，所做的本是我们应分做的}\BibleRef{路 17:10}。为使愁容不破坏慷慨，让他们听所记载的：\BibleQuote{天主喜爱乐意捐献的人}\BibleRef{格后 9:7}。为使不寻求对所施礼物的短暂赞美，让他们听所记载的：\BibleQuote{不要叫左手知道右手所做的}\BibleRef{玛 6:3}。即是，不要让现世的光荣掺杂于虔诚的施舍，也不要让恩惠的愿望知晓任何正直的事工。为使不要求所施恩惠的回报，让他们听所记载的：\BibleQuote{你设午宴或晚宴，不要请你的朋友、兄弟、亲戚和富有的邻舍，恐怕他们也回请你，你就得了报答。但你摆设筵席，要请贫穷的、残废的、瘸腿的、瞎眼的，你就有福了，因为他们没有什么可报答你}\BibleRef{路 14:12}。为使不延迟供应本该立刻供应的事，让他们听所记载的：\BibleQuote{你不可对你的朋友说：你去吧，明天再来，我必给你，其实你即刻可以给他}\BibleRef{箴 3:28}。为使不要以慷慨为借口徒然分散所拥有的，让他们听所记载的：\Quote{让你的施舍在你手中出汗。}为使在多需时却给得少，让他们听所记载的：\BibleQuote{少种的少收}\BibleRef{格后 9:6}。为使在应给得少时却给得太多，以致后来自己不堪忍受匮乏而爆发急躁，让他们听所记载的：\BibleQuote{这不是要别人轻省，你们受累，乃是要均平，乃要你们的富余补足他们的缺乏，叫他们的富余也补足你们的缺乏}\BibleRef{格后 8:13-14}。因为当施予者的灵魂不知如何忍受匮乏时，他在从自己身上多拿走后，便为自己寻找急躁的机会。因为心应先准备好忍耐，然后才慷慨地施予许多或全部，免得稍有不平忍受匮乏的侵袭，既丧失先前慷慨的赏报，随后抱怨更糟地摧毁灵魂。为使对那些本应施予的人一点都不给，让他们听所记载的：\BibleQuote{凡求你的，就给他}\BibleRef{路 6:30}。为使对那些本不应施予的人有所给，即使极少，让他们听所记载的：\BibleQuote{你要施舍给善人，不要接济罪人；要善待卑微的人，不要给不敬虔的人}\BibleRef{德 12:4}。又说：\BibleQuote{要在义人的坟墓上摆设你的饼和酒，但不可与罪人同吃同喝}\BibleRef{多 4:17}。\par

因为凡将饼和酒给罪人者，乃是因他们是罪人而给予援助。为此，世上一些富人用丰厚的赏赐供养戏子，而基督的穷人却忍饥挨饿。然而，凡将饼给贫乏者，即使他是罪人，不是因他是罪人，而是因他是人，他实际上并非供养罪人，而是供养一个贫穷的义人，因为他在那人身上所爱的，不是他的罪，而是他的人性。\par

那些已慷慨施舍自己财物的人，也应受劝诫，要小心提防，免得当他们用施舍赎回了已犯的罪时，又犯下仍需赎罪的过犯；免得他们认为天主的公义是可以买卖的，以为只要花钱赎罪，便可逍遥法外。因为，\BibleQuote{生命贵于食物，身体贵于衣服}\BibleRef{玛 6:25}。因此，谁若将食物或衣服施舍给穷人，却仍被灵魂或身体的不义所玷污，便是将较次要之物献给了正义，却将更重要之物献给了罪恶；因为他将自己的财产交给了天主，却将自己交给了魔鬼。\par

但另一方面，那些仍贪图夺取他人财物的人，应受劝诫，要留心听主在审判时所说的话。因为祂说：\BibleQuote{我饿了，你们没有给我吃的；我渴了，你们没有给我喝的；我作客，你们没有收留我；我赤身露体，你们没有给我穿的；我患病，我在监里，你们没有来探望我。}\BibleRef{玛 25:42} 祂先前又对他们说：\BibleQuote{可咒骂的，离开我，到那为魔鬼和他的使者预备的永火里去吧！}\BibleRef{玛 25:41} 请看，他们绝未被告知犯过抢劫或任何暴力行为，却仍被投入革赫纳的永火。由此可以推断，那些夺取他人财物的人将受到何等严厉的惩罚，倘若那些不明智地守着自己财物的人竟遭受如此大的惩罚。让他们想想，夺取财物必会带来何等罪责，若不分施财物竟招致如此刑罚。让他们想想，不义之举该当何等处分，若吝于施恩竟值得如此重罚。\par

当他们一心夺取他人财物时，让他们听听经上记载的：\BibleQuote{祸哉，那增加不属于自己之物的人！他堆积厚泥自害，要到几时呢？}\BibleRef{哈 2:6} 的确，贪婪之人堆积厚泥自害，就是将地上的收益堆成罪恶的重担。当他们渴望大大扩充居住的空间时，让他们听听经上记载的：\BibleQuote{祸哉，你们这些使房屋连接房屋，使田地连接田地，直到再无空地的人！难道你们要独自住在大地中间吗？}\BibleRef{依 5:8} 这仿佛在明白地说：你们这些不能容忍在共同世界中有同伴的人，要伸展到多远呢？你们压制那些与你们相连的人，并总是找到一些你们可以伸展自己的对象。当他们一心增加钱财时，让他们听听经上记载的：\BibleQuote{贪爱钱财的人，不因金钱而满足；爱财富的人，也不得其收益。}\BibleRef{训 5:9} 因为他们本可得到收益，如果他们存心不贪爱钱财，好好分施。但谁若在情感上执着于钱财，必将在此地留下它们，毫无收益。当他们急于一下子拥有各种财富时，让他们听听经上记载的：\BibleQuote{急于致富的人，不能无罪。}\BibleRef{箴 28:20} 因为谁若谋求增加财富，必会疏忽避免罪恶；如同鸟儿一样，贪看地上的诱饵时，不知不觉被罪恶的圈套勒住。当他们贪求现世任何利益，却不知来世将受损失时，让他们听听经上记载的：\BibleQuote{起初急于获得的产业，终必不蒙祝福。}\BibleRef{箴 20:21} 因为我们从今生开始，为的是在终结时能到达蒙福的份。因此，那些在起初急于获得产业的人，就在终结时断绝了自己蒙福的份；因为当他们因贪婪的不义而渴望在此地增加财富时，他们就在那里失去了永恒的基业。当他们或过分恳求，或成功地获得所恳求的一切时，让他们听听经上记载的：\BibleQuote{人纵然赚得了全世界，却赔上了自己的灵魂，为他有什么益处？}\BibleRef{玛 16:26} 这仿佛真理在明白地说：人即使聚集了身外的一切，却毁灭了自身，这对他有何益处呢？但大多数情况下，掠夺者的贪财若能经由劝诫者的话指出现世生命的短暂，就较易改正；若能提及那些长久力图在此世致富，却不能长久留在所获财富中的人，从他们那里，突然的死亡将一切（既非一下子也非突然聚集的）瞬间夺走；他们不仅在此地留下了所夺取的，还带着掠夺的控告到了审判台前。因此，应向他们讲述这类人的例子，他们自己无疑也会在言语上加以谴责；这样，他们在言语之后回到自己心中，也许会因羞耻而不去效法他们所审判的人。\par

% structure: p0138 source=h2 target=subsection reason=relative-heading-rank; base=h2

\subsection{第二十一章}

\Emph{如何劝诫那些不贪图他人财物而守自己财物的人，以及那些施舍自己财物却夺取他人财物的人。}\par

(\Emph{劝谕} 22.) 应当以不同的方式告诫那些既不想拥有他人的财物，也不肯施舍自己财物的人，以及那些施舍自己财物却仍不停止夺取他人财物的人。那些既不想拥有他人财物、也不肯施舍自己财物的人，应当被告诫要仔细考虑：他们从中被取出的土地是众人共有的，因此它也为众人共同带来养料。因此，那些将天主的公共恩赐归于自己私用的人，虚妄地自以为无罪；那些不肯分施所得的人，在邻人的杀戮中行走；因为他们几乎每天杀害的人，等于那因他们紧握自己手中而死亡的穷人的数目。因为当我们向贫困者提供任何必需品时，我们并非施舍自己的东西，而是归还属于他们的东西；我们宁可说是偿还正义的债务，而非完成慈悲的工作。因此，真理本身在谈论施舍慈悲时所需的谨慎时也说：\BibleQuote{你们应当小心，不要在人前施行你们的义德}\BibleRef{玛 6:1}；圣咏作者也同意这句话说：\BibleQuote{他散施钱财，他赒济贫苦，他的义德永远常存}\BibleRef{咏 111:9}。\par

因为，他先提到对穷人的慷慨施舍，却不愿意称此为慈悲，而宁可称之为正义：因为凡接受由共同的主所赐予之物的人，应当共同使用，这确实是正义的。因此，\Person{撒罗满}也说：\BibleQuote{正义的人必施舍，绝不吝惜}\BibleRef{箴 21:26}。他们还应被告诫要急切地注意，那严厉的农夫如何抱怨那棵不结果的无花果树，说它甚至连地也占用了。\par

因为一棵不结果的无花果树占用了土地，当吝啬的灵魂无益地扣留着可能有利于许多人的东西时。一棵不结果的无花果树占用了土地，当愚昧的人在懒惰的荫庇下使一块本可在善工之阳光下被他人耕种的土地保持荒芜时。\par

但这些人与常会说：\Quote{我们使用所赐予我们的东西；我们不追求他人的财物；即使我们没有做任何配得上慈悲赏报的事，我们也没有犯任何错。}他们这样想，是因为他们实在闭起了心灵之耳，不听来自天上的话语。因为福音中那位身穿紫红细麻衣、天天奢华宴乐的富人，并没有被说成夺取了别人的东西，而是无益地使用了属于自己的东西；报应的地狱在他死后接纳了他，不是因为他做了任何非法的事，而是因为过度放纵，他完全沉溺于合法的事物中。\par

应当告诫吝啬的人注意，他们首先对天主行了不义：他们从赐予他们一切的天主那里，什么都没有回报，没有慈悲的祭献。因为圣咏作者说：\BibleQuote{他不把赎价献给天主，也不付出他灵魂的代价}\BibleRef{咏 48:9}。\Person{若翰}也高声喊道：\BibleQuote{斧子已经放在树根上了；凡不结好果子的树，都要被砍倒，投入火中}\BibleRef{路 3:9}。因此，那些自以为无罪，因为不夺取他人财物的人，应当期待近在眼前的斧头一击，并放下不谨惧安全的麻木，免得他们因忽略结出善工的果子，而从现在的生命中被完全砍除，如同从树根的青翠中被砍掉一样。\par

然而，另一方面，应当告诫那些既施舍自己所有、又不停止夺取他人之物的人，不要渴望显得过分慷慨，因而因外在的善行而变得更糟。因为这些人，不明智地施舍自己的东西，不仅如前所述陷入不耐烦的抱怨，而且当贫穷迫使他们时，甚至被卷向贪婪。那么，还有什么比这些人的心态更悲惨呢？在他们心中，贪婪生于慷慨，罪孽的收成仿佛从美德中播下。因此，首先应当告诫他们学习如何合理地保有自己之物，然后最终不再去获取他人之物。因为，如果错误的根源不在浪费本身中被烧尽，那么贪婪的荆棘，通过枝条繁茂，就永远不会干枯。因此，如果先妥善调整所有权，夺取的理由就被消除了。然而，进一步地，应当告诉那些被告诫的人，当他们学会不因将抢劫的邪恶混入慈悲的善行之中而混淆慈悲的善时，如何慈悲地施舍自己所拥有的。因为他们强取豪夺，然后慈悲地施舍。因为，因我们的罪而显示慈悲是一回事；因显示慈悲而犯罪是另一回事；这确实不能再称为慈悲，因为它无法结出甜美的果实，反而被其有害之根的毒药所苦。因此，主通过先知甚至拒绝祭献本身，说：\BibleQuote{我上主喜爱正义，憎恨不义之祭}\BibleRef{依 61:8}。祂又说：\BibleQuote{不虔敬之人的祭献是可憎恶的，因为他们用邪恶献祭}\BibleRef{箴 21:28}。这些人也常常从穷人那里取走他们献给天主的东西。\par

但主以何等严厉的谴责表明祂不认他们，藉某位智者说：\BibleQuote{凡以穷人的财物献祭的，无异于在父亲眼前杀其儿子}\BibleRef{德 34:20}。因为还有什么比儿子在父亲眼前死去更难以忍受的呢？由此可见，这种祭献被视为何等大的忿怒，竟被比作丧子父亲的悲伤。然而，人们多半仔细衡量自己奉献了多少，却忽略考虑自己夺取了多少。他们算自己的工资，却不肯考虑自己的过犯。因此，让他们听听所记载的：\BibleQuote{那积聚工资的，将工资放入有洞的钱袋}\BibleRef{盖 1:6}。的确，钱放入有洞的钱袋，放入时看得见，丢失时却看不见。所以，那些只盯着自己施舍了多少，而不考虑自己夺取了多少的人，正是将工资放入有洞的钱袋，因为实际上，他们看着工资积聚起来，以为稳妥，却不知不觉地失去了它们。\par

% structure: p0147 source=h2 target=subsection reason=relative-heading-rank; base=h2

\subsection{第二十二章}

\Emph{如何劝诫争吵者与和睦者}\par

(\Emph{劝谕} 23.) 应不同方式劝诫那些彼此不和的人与那些和睦相处的人。因为那些彼此不和的人应受劝诫，要最确切地知道，无论他们拥有何等丰富的德行，若忽略与邻人通过和睦结合，便绝不可能成为属神的人。因为经上记载：\BibleQuote{但圣神的效果是仁爱、喜乐、平安}（\BibleRef{迦 5:22}）。因此，那不留意保持和平的人，就是拒绝结出圣神的果实。为此 \Person{保禄} 说：\BibleQuote{你们中间既然有嫉妒和纷争，你们岂不是属血肉的吗？}（\BibleRef{格前 3:3}）。为此他又说：\BibleQuote{你们要竭力与众人和平相处，务要圣洁，因为没有圣洁，谁也见不到主}（\BibleRef{希 12:14}）。为此他又劝诫说：\BibleQuote{尽力以和平的联系，保持心神的合一，因为只有一个身体和一个圣神，正如你们蒙召，同有一个希望一样}（\BibleRef{弗 4:3}）。因此，我们蒙召的唯一希望永远无法达到，如果我们不与邻人心意合一地奔向它。但常常发生的是，有些人因对自己某些特别拥有的恩赐感到骄傲，而失去了那更大的和睦恩赐；例如，一个人若因节制情欲比他人更能克制肉身，便不屑与那些在禁欲上不如他的人和睦相处。但谁若将禁欲与和睦分开，就请留意圣咏作者的劝诫：\BibleQuote{用鼓和舞蹈赞美他}（\BibleRef{咏 150:4}）。因为在鼓中，干枯被打的兽皮发出声响，而在合唱中，声音在和谐中结合在一起。因此，谁若刻苦自己的身体，却抛弃和睦，他确是用鼓赞美天主，却不是用合唱赞美祂。然而，当优越的知识提升某些人时，常常使他们脱离与其他人的交往；他们越有智慧，在和睦的德行上反而越不智慧。因此，这些人应听取真理本身所说的话：\BibleQuote{你们当中有盐，并彼此和睦相处}（\BibleRef{谷 9:50}）。因为确实，没有和睦的盐不是德行的恩赐，而是定罪的证据。因为一个人越有智慧，他的过失就越严重，他将因此无可推诿地应受惩罚，因为他本可用自己的明智避免罪恶。对此，通过 \Person{雅各伯} 恰当地对他们说：\BibleQuote{如果你们心中怀有苦毒的嫉妒和纷争，就不必夸耀，你们不可说谎违背真理。这种智慧不是从上而来的，而是属于地、属于血肉、属于魔鬼的。但从上而来的智慧，先是纯洁的，其次是和平的}（\BibleRef{雅 3:14}）。纯洁，就是说，因为它的思想是贞洁的；同时也是和平的，因为它绝不因高傲而脱离与邻人的交往。那些彼此不和的人应受劝诫，要留意，只要他们不与邻人存有爱德，他们向天主献上的善工之祭便不被悦纳。因为经上记载：\BibleQuote{所以，你若在祭台前献你的礼物，在那里想起你的兄弟有什么怨你的事，就把你的礼物留在祭台前，先去与你的兄弟和好，然后再来献你的礼物}（\BibleRef{玛 5:23}）。由此，我们得以看出人的罪是多么不可容忍，竟使他们的祭献被拒绝。因为既然一切罪恶在善行后继时都被洗涤，我们当思想，纷争的罪恶是何等重大，除非它被彻底消灭，否则就不容许任何善行随之而来。那些彼此不和的人应受劝诫，若他们不愿侧耳听从天主的命令，就该睁开心灵的眼睛，去观察最低等受造物的方式；常常，同一种类的鸟类在其群体飞行中互不离弃，野兽也成群结队地觅食。因此，若我们明智地观察，无理性的本性借着一致，表明有理性的人因不一致而犯下多大的恶；后者因运用理性而失去了前者凭自然本能所保持的。但另一方面，那些和睦相处的人应受劝诫，要警惕，以免他们过度喜爱自己所享有的和平，而不渴望达到那永久的和平。因为通常，平静的环境更严厉地考验心灵的倾向，以至于它们所占据的事物越不麻烦，吸引它们的事物就显得越不可爱，现时的事物越令人喜悦，永恒的事物就越不被追求。因此，真理本身在亲自说话时，为了区分地上的和平与天上的和平，并激励祂的门徒从现在转向将来，说：\BibleQuote{我把平安留给你们，我将我的平安赐给你们}（\BibleRef{若 14:27}）。也就是说，我留下暂时的平安，我赐予长存的平安。因此，若心灵执着于那被留下的，就永远达不到那将被赐予的。所以，现世的平安应被视为既可爱又可轻，以免若过度爱慕它，那爱者的心便陷入过错。因此，那些和睦相处的人也应受劝诫，以免在过于渴望人的和平时，完全不去指责人的恶行，因附和悖逆之人而脱离造物主的平安；以免他们在外面畏惧人类的争吵，内心却被内在盟约的破裂所击打。因为暂时的和平若非永恒和平的足迹，又是什么？因此，有什么比爱慕印在尘土上的足迹，却不爱那印下足迹的人更疯狂的呢？为此，当 \Person{达味} 要完全投身于内在和平的盟约时，他见证自己不与恶人妥协，说：\BibleQuote{上主，憎恨你的人，我怎能不憎恨他们？攻击你的人，我怎能不厌恶他们？我彻底憎恨他们，把他们当作我的仇敌}（\BibleRef{咏 138:21}）。因为以完全的恨憎恨天主的仇敌，就是爱他们受造的本性，并指责他们的行为，严厉对待恶人的作风，并有益于他们的生命。因此，需要好好思量，当从指责中得享安息时，与最坏的人保持和平是何等有罪，如果如此伟大的先知将此作为祭献献给天主，就是他为主激起了恶人对自己的仇恨。为此，肋未支派当他们拿起刀剑穿过营中，因不肯宽恕当受击打的罪人时，被称为将自己的手奉献给天主（\BibleRef{出 32:27}）。为此，\Person{丕乃哈斯} 藐视同胞犯罪时的偏爱，击杀了那些与米德杨人结合的人，用他的愤怒平息了天主的愤怒（\BibleRef{户 25:9}）。为此，真理本身亲自说：\BibleQuote{你们不要以为我来，是为把平安带到地上；我来不是为带平安，而是带刀剑}（\BibleRef{玛 10:34}）。因为，当我们不谨慎地与恶人结为朋友时，我们就与他们同负罪债。因此，\Person{约沙法特} 虽被以往生活诸多赞颂所称赞，却因与 \Person{阿哈布} 王的友谊而几乎被毁灭，当先知对他说：\BibleQuote{你岂是帮助恶人，爱那憎恨上主的人？为此，上主的忿怒临于你。不过你还是行了善事，因为你从国内铲除了木偶，专心致志寻求了上主}（\BibleRef{编下 19:2}）。因为，当我们在悖逆者的友谊中一致时，我们的生活已经与那至公义者相违背。那些和睦相处的人应受劝诫，不要畏惧扰乱他们暂时的和平，如果他们说出责备的话。他们也应受劝诫，要在内心以不减的爱德保持那他们在外面因责备的声音而扰乱的和平。这二者，\Person{达味} 宣称自己已明智地遵守了，说：\BibleQuote{我愿与憎恨和平的人和平相处，但我一讲话，他们就无故攻击我}（\BibleRef{咏 119:7}）。看，当他讲话时，他受到攻击；然而，当他受到攻击时，他仍然是和平的，因为他既没有停止责备那些对他发怒的人，也没有忘记爱那些被责备的人。为此，\Person{保禄} 也说：\BibleQuote{如若可能，应尽力与众人和平相处}（\BibleRef{罗 12:18}）。因为，在劝诫门徒与众人和平相处之前，他先说：\BibleQuote{如若可能}，又加上：\BibleQuote{尽力}。因为对他们来说，如果他们责备恶行，就难以与众人和平相处。但是，当我们因责备而在恶人心中扰乱了暂时的和平时，有必要在我们的心中保持它不受损害。因此，他恰当地说：\BibleQuote{尽力}。确实，他好像在说：既然和平在于双方的同意，如果它被受责备的人赶走，那么它仍应毫无减损地保留在你们这些责备的人心中。为此，同一宗徒再次劝诫他的门徒，说：\BibleQuote{如果有人不听从我们这封信上的话，要记下这人，不要与他交往，好叫他羞愧}（\BibleRef{得后 3:14}）。他立刻又加上：\BibleQuote{但不要把他当作仇敌，而要规劝他如兄弟}（\BibleRef{得后 3:15}）。仿佛在说：与他断绝外在的和平，但在你内心深处，要为他保持内在的和平；这样，你与他的不和可以打击罪人的心灵，而和平即使不给予他，也不会从你们心中消失。\par

% structure: p0150 source=h2 target=subsection reason=relative-heading-rank; base=h2

\subsection{第二十三章}

\Emph{如何劝诫播散纷争者与缔造和平者。}\par

（\Emph{劝诫}第二十四。）要辨别加以劝诫的是播散纷争者与缔造和平者。因为播散纷争者应被劝诫认识自己是谁的追随者。关于那叛变的天使，经上记载，当莠子被撒在好庄稼中时，\BibleQuote{这是仇人做的}（\BibleRef{玛 13:28}）。关于他的一个成员，经上借撒罗满的口说：\BibleQuote{无赖之徒，奸诈的人，行走口舌乖张，他以眼传神，以脚示意，以手指点，心藏欺诈，常谋恶事，播散纷争}（\BibleRef{箴 6:12}）。看，他愿意称之为播散纷争的人，首先称他为叛徒；因为除非他像那骄傲的天使一样，先因心思远离造物主的面而内心堕落，就不会后来向外播散纷争。他也被恰当地描述为以眼传神、以手指点、以脚示意。因为内在的警醒使肢体在外部秩序井然。那么，失去了内心稳定的人，便在外部陷入动作的轻浮不定，其外在的浮动表明他里面没有扎根。让播散纷争者听到经上所说：\BibleQuote{缔造和平的人是有福的，因为他们要称为天主的子女}（\BibleRef{玛 5:9}）。另一方面，让他们明白，如果缔造和平的人被称为天主的子女，无疑那些混淆和平的人就是撒殚的子女。此外，凡因纷争而与仁爱的绿意分离的人，都枯干了；虽然他们在行为上结出善行的果实，但实际上没有果实，因为它们不是出自爱德的合一。因此，让播散纷争者思想他们犯的罪是多么众多；因为当他们犯一个不义时，同时从人心拔除了所有德行。因为在一个恶事上，他们成就了无数恶事，因为播散纷争，他们熄灭了爱德，而爱德实际上是所有德行之母。但是，既然在天主眼中没有比仁爱的德行更宝贵的，在魔鬼眼中也没有比爱德的熄灭更可取的。那么，凡通过播散纷争摧毁邻人仁爱的人，是作为密友事奉天主的仇敌；因为他们夺走了邻人的这一点，而魔鬼正是因失去这一点而堕落，他们便切断了邻人上升的道路。\par

但另一方面，缔造和平者应被劝诫，不可因不知道应在谁之间建立和平而减损如此伟大事业的效果。因为，正如善人缺乏合一有很大害处，同样，恶人不缺乏合一也有极大的害处。如果悖逆者的不义在和平中联合，那么他们的恶行必然会增添力量，因为他们越在邪恶中彼此一致，就越猛烈地冲击善人，折磨他们。因此，关于那注定灭亡的器皿，即假基督的传道者，天主的声音对可敬的约伯说：\BibleQuote{牠的肉块互相联结}（\BibleRef{约 41:14}）。因此，在鳞甲的比喻下，关于牠的随从说：\BibleQuote{它们彼此相连，甚至气也不能透入其间}（\BibleRef{约 41:7}）。的确，牠的追随者之间没有纷争分裂，因此他们在杀害善人时结合得更加紧密。那么，谁使不义者在和平中联合，就是给不义增添力量，因为他们更加一致地迫害善人，压垮他们。因此，那位杰出的宣讲者，当受到法利塞人和撒杜塞人的暴力迫害时，他试图分裂那些他看见激烈联合起来反对自己的人，他大声喊道：\BibleQuote{诸位弟兄！我是法利塞人，是法利塞人的儿子；我是为了希望和死人的复活而受审}（\BibleRef{宗 23:6}）。而由于撒杜塞人否认希望和死人的复活，法利塞人则按照圣经的诫命相信这一点，于是在迫害者的一致中产生了分歧；保禄得以从分裂的人群中安全逃脱，而之前当这些人联合时曾凶猛地攻击他。因此，那些致力于缔造和平的人，应被劝诫：他们首先应当将内在和平的爱注入悖逆者的心中，以使外在的和平以后能有益于他们；这样，当他们的心悬于认识前者时，就绝不会因感知后者而被卷入邪恶；当他们看到面前超性的事物时，就不会将世俗的事物转而用来损害自己。但是，如果有任何悖逆之人，即使他们想要伤害善人也做不到，那么在他们之间，甚至在上升到认识超性和平之前，就应当建立地上的和平；甚至使那些被不虔敬的邪恶激怒而反对天主仁爱的人，至少能因爱邻人而软化，并好似从邻近的位置过渡到更好的位置，从而上升到尚远离他们的、造物主的和平。\par

% structure: p0154 source=h2 target=subsection reason=relative-heading-rank; base=h2

\subsection{第二十四章}

\Emph{如何劝诫在神圣学问上粗浅者与有学问但不谦卑者。}\par

(\Emph{劝言} 25.) 那些不能正确理解圣经之言的人，与那些虽能正确理解却不谦卑地宣讲的人，应以不同的方式加以劝诫。因为那些不能正确理解圣经之言的人，应当被劝诫去思考：他们为自己将最有益健康的美酒变成了毒杯，并用医疗之刀给自己造成致命创伤，因为那本应为了健康而切除病变之处的工具，他们却用它摧毁了自己健全的部分。他们应当被劝诫去思考：圣经犹如一盏明灯，放在我们今生的黑夜中；若他们不能正确理解其中的话，便会从光明中得到黑暗。但实际上，若非骄傲先使他们膨胀，扭曲的心志不会促使他们错误地理解。因为他们自以为比众人聪明，不屑于跟随他人理解更好的事物；为了从无知的群众中为自己强夺一个博学的名声，他们竭力推翻他人的正确见解，并坚持自己的谬见。因此，先知说得对：\BibleQuote{他们剖开了基肋阿得怀胎的妇人，为要扩大自己的疆界} \BibleRef{亚 1:13}。因为基肋阿得按解释是见证之堆 \BibleRef{创 31:47}。而且，既然教会的全体会众通过其宣认为真理作证，那么用基肋阿得来表达教会并非不当，教会通过所有信友的口为关于天主的一切真理作证。此外，当灵魂因天主的爱而怀上对圣言的理解时，她们被称为怀胎；这样，若她们足月，就能在工作的彰显中生出所孕育的智慧。再者，扩大疆界就是扩展他们声誉的名声。因此，他们剖开了基肋阿得怀胎的妇人为要扩大疆界，因为异端者确实通过他们邪恶的讲道杀害了那些已经孕育了一些真理理解的信友灵魂，并为自己扩展了博学的名声。他们用错误的剑劈开那些已因对圣言的理解而怀胎的小子们的心，并借此为自己博得教师的声誉。因此，当我们努力教导他们不要居心邪僻时，必须首先劝诫他们躲避虚荣。因为如果昂首自傲的根被切断了，错误论断的枝子也就会随之枯萎。他们还应被劝诫要留心，以免因制造错误与纷争，而把天主赋予为了阻止祭献撒但的同一法律，变成了祭献撒但的祭献。因此，主通过先知抱怨说：\BibleQuote{我赐给他们五谷、酒和油，并增加了他们的金银，他们却把这些祭献给巴耳} \BibleRef{欧 2:8}。因为当我们从更隐晦的言语中，剥去字面的外壳，在圣神之精髓中领悟法律的内在意义时，我们确实从主那里领受了五谷。主赐给我们祂的酒，当祂用圣经崇高的宣讲使我们陶醉时。祂也赐给我们祂的油，当祂通过更明确的诫命，温和而顺畅地规范我们的生活时。祂增加银子，当祂供给我们充满真理之光的雄辩言辞时。祂还用金子使我们富足，当祂用对至高光辉的理解照亮我们的心时。所有这一切，异端者都献给了巴耳，因为他们通过腐败的理解，在听众的心中扭曲了这一切。并且，他们将天主的五谷、酒、油，以及祂的银子与金子，作为祭献献给撒但，因为他们把和平之言转向助长纷争的错误。因此，他们应被劝诫去思考：当他们以邪僻之心从和平的诫命中制造纷争时，他们自己，在天主公正的审判中，必因生命之言而死亡。\par

但另一方面，那些确实正确理解法律之言，却不谦卑宣讲的人，应受劝诫：在神圣的讲论中，在他们向别人宣讲之前，应先省察自己；免得在追随他人的行为时，把自己抛在后面；免得在对圣经其余所有部分持正确见解时，唯独不注意其中对骄傲者所说的话。因为那想治好别人疾病，却不知自己正受何苦的医生，实在是一个贫乏且无能的医生。因此，那些不谦卑宣讲天主之言的人，当然应受劝诫：当他们为病人敷药时，要留心自己感染的毒害，免得在治愈他人时自己却死去。他们应受劝诫，要当心，免得他们说话的方式与所说之事的高贵相悖，免得他们口说一套，举止又一套。因此，请听经上所记载的：\BibleQuote{若有人讲论，就该讲论如天主的话}\BibleRef{伯前 4:11}。那么，如果他们所说的话不是出于自己的东西，为什么还要为此自傲，仿佛是他们自己的一样呢？请听经上所记载的：\BibleQuote{我们是出于天主，在天主台前，在基督内讲论}\BibleRef{格后 2:17}。因为那出于天主、在天主台前讲论的人，既明白他领受了宣讲之言来自天主，也藉此寻求取悦天主，而非取悦人。请听经上所记载的：\BibleQuote{凡心中骄傲的，为上主所憎恶}\BibleRef{箴 16:5}。因为，当他在天主圣言中寻求自己的荣耀时，他侵犯了赐予者的权利；他毫不畏惧地将赞美归于自己，而他所受赞美的正是从赐予者领受的。请听上主藉撒罗满对宣讲者所说的话：\BibleQuote{你要喝自己池中的水，饮自己井里的活水。让你的泉源向外涌流，将你的水分散在各街市上。但唯独归你独享，不要让外人与你同享}\BibleRef{箴 5:15}。的确，宣讲者喝自己池中的水，是当他返回自己的心，首先自己聆听他必须说的话。他饮自己井里的活水，是当他被自己的言语所滋润。接着，经上恰当地加上：\BibleQuote{让你的泉源向外涌流，将你的水分散在各街市上}。因为，他首先应当自己喝，然后才能藉宣讲流到别人身上。因为使泉源向外涌流，就是将宣讲的力量向外倾注给他人。此外，将水分散在各街市上，就是按照每个人的品质，将神圣的言语分施给众多听众。并且，因为当天主圣言自由通向许多人的认识时，虚荣的欲望多半会悄然潜入，所以在说了\BibleQuote{将你的水分散在各街市上}之后，又恰当地加上：\BibleQuote{唯独归你独享，不要让外人与你同享}。他在这里称恶灵为外人，关于他们，经上藉受诱惑者的话说：\BibleQuote{外人起来攻击我，强暴者寻索我的性命}\BibleRef{咏 53:3}。因此他说：\BibleQuote{既要将你的水分散在各街市上，又要唯独归你独享}；好像更明白地说：你们必须在宣讲上向外服务，但不可因自傲而与不洁之灵联合，免得在神圣言语的职务中，容让仇敌与你们同享。因此，我们将我们的水分散在各街市上，却又唯独占有它们，当我们同时向外广泛倾注宣讲，却又绝不藉此寻求人的赞美时，便是如此。\par

% structure: p0158 source=h2 target=subsection reason=relative-heading-rank; base=h2

\subsection{第二十五章}

\Emph{应如何劝诫那些因过分谦卑而推辞宣讲职务的人，以及那些轻率急切抓住它的人。}\par

（\Emph{劝诫}第二十六。）应不同地劝诫那些虽然能妥善宣讲，却因过分谦卑而退缩的人，以及那些因不成熟或年龄所限而不能宣讲，却被轻率所驱使的人。因为那些虽能有益地宣讲，却因过分谦卑而退缩的人，应受劝诫，从较小的事例中认识到他们在更大事上的过失。因为，如果他们向贫困的邻人隐藏自己所拥有的金钱，无疑会成为其灾难的推手。那么，他们应当明白，当自己向犯罪的弟兄隐瞒宣讲之言时，便是从垂死的灵魂中藏起了生命之药，这是何等的罪过。因此，一位智者说得对：\BibleQuote{隐藏的智慧，和看不见的宝藏，两者有何益处？}\BibleRef{德 20:32}。倘若饥荒使百姓衰亡，而他们却隐藏谷物，他们无疑就是死亡的制造者。因此，要想想当灵魂因饥渴圣言而丧亡时，他们不拿出自己所领受的恩宠之粮，将受到何等的惩罚。因此，经上藉撒罗满说得恰当：\BibleQuote{隐藏谷物的，必受民众诅咒}\BibleRef{箴 11:26}。因为隐藏谷物，就是自己保留神圣宣讲的言语。凡这样做的人，必受民众诅咒，因为他仅仅是因沉默之过，就受那许多本可得矫正之人的惩罚。倘若某个绝不缺乏医术之人看到一个需要切开疗治的伤口，却拒绝切开，那么仅凭他的不作为，他对兄弟的死亡就难逃罪责。那么，他们当明白，当自己明知灵魂的伤口，却疏于用言语的刀来医治时，是何等大的罪过。因此，经上藉先知说得恰当：\BibleQuote{凡阻止自己的刀剑不沾血的，应受诅咒}\BibleRef{耶 48:10}。因为阻止刀剑不沾血，就是阻止宣讲之言去杀戮肉体的生命。关于这把剑，经上又说：\BibleQuote{我的刀剑要吞噬血肉}\BibleRef{申 32:42}。\par

因此，当这些人将宣讲之言缄默自守时，当畏惧地聆听天主对他们的判决，好让恐惧从他们心中驱走恐惧。让他们听到那不愿动用自己\Quote{塔冷通}的人如何失去它，并加上定罪的判决，\BibleRef{玛 25:24}。让他们听到保禄如何认为自己与邻人的血无干，正是因为他毫不姑息那应受责打的恶习，说：\Quote{\Emph{我今天向你们作证：我是无罪的，因为我没有避讳将天主的全部计划宣布给你们}}\BibleRef{宗 20:26}。让他们听到若望如何受天使之声的劝诫，当祂说：\Quote{\Emph{让听见的人说：\InnerQuote{来吧！}}}\BibleRef{默 22:17}；无疑是为了让那内心已容纳内在声音的人，藉高声呼喊将他人引向他所前往之处，免得他即使被召，若空空地走近那召叫他的人，却要发现门已关闭。让他们听到依撒意亚，当他被天上的光照亮后，因在圣言的职务中缄默不语，便以大声的忏悔自责说：\Quote{\Emph{我有祸了！因为我缄默不语}}\BibleRef{依 6:5}。让他们听到撒罗满应许说，那不被怠惰之习束缚于已达之境的人，宣讲的知识将倍增。因为他说：\Quote{\Emph{祝福的灵魂必得丰盈；那灌溉的人，也必被灌溉}}\BibleRef{箴 11:25}。因为那藉宣讲向外祝福的人，领受内在增长的丰盈；当他不断用口才之酒灌溉听众的心灵时，他也因倍增恩赐的畅饮而愈发沉醉。让他们听到达味如何将此作为礼物献给天主，即他不隐藏所领受的宣讲恩宠，说：\Quote{\Emph{看，我必不缄默我的唇舌；上主，你知道，我没有将你的正义藏于心中；我宣告了你的真理和你的救恩}}\BibleRef{咏 39:10}。让他们听到新郎与新妇对话时所说的话：\Quote{\Emph{你住在花园中，你的朋友听闻：让我听到你的声音}}\BibleRef{歌 8:13}。因为教会住在花园中，因为她将培育的美德苗圃保持内在的葱绿。而她的朋友听闻她的声音，是指所有蒙选者渴望她宣讲的话语；新郎也渴望听到这声音，因为他通过蒙选者的灵魂渴慕她的宣讲。让他们听到梅瑟如何看到天主对祂的子民发怒，并命令拿起刀剑执行报复，便宣布那些毫不犹豫击打犯罪者罪行的人属于天主一边，说：\Quote{\Emph{凡属上主的，到我这里来！各人把刀佩在腰间，从这门到那门，穿越营中，击杀自己的兄弟、朋友和邻人}}\BibleRef{出 32:27}。因为把刀佩在腰间，是将宣讲的热诚置于肉体的享乐之前；这样，当一个人热衷于说出圣言时，他必须留意制服不当的暗示。但从这门到那门，是指从一个恶习到另一个恶习来回奔走责备，直至每一个使死亡进入灵魂的恶习。而穿越营中，是指在教会内如此公正地生活，以至责备罪人罪行时不偏袒任何人。因此也正确加上：\Quote{\Emph{击杀自己的兄弟、朋友和邻人}}。他真正击杀兄弟、朋友和邻人，是当他发现应受惩罚之事时，连因亲情所爱的人也不从责备之剑下放过。因此，若那因神圣之爱的热忱而奋起击打恶习的人被称为属于天主，那么那拒绝尽己所能责备肉性之人的生活的人，便是否认自己属于天主。\par

但是，另一方面，那些因不完美或年龄而被排除于宣讲职务之外，却又因急躁而冲动去宣讲的人，应受劝诫，以免他们冒然将这些重大职务的重担加于自身，反而断绝了自己后来进步的道路；并因不适时地夺取他们力不能及的事，而连他们或许能在适合时期完成的事也失去了；并且因为他们企图不适当地展示自己的知识，而显得理应失去知识。他们应受劝诫，要思考：幼鸟若在翅膀尚未长成之前尝试飞翔，就会从它们本希望高飞之处重重坠落。他们应受劝诫，要思考：如果在新建筑尚未坚固之前就铺上木料的重担，所建造的不是居所，而是废墟。他们应受劝诫，要思考：如果妇女在怀孕的胎儿尚未完全成形之前就生产，她们所充满的不是房屋，而是坟墓。因此，真理本身（祂本可瞬间坚固祂所愿意的人）为了给祂的跟随者一个榜样，使他们不敢在不完善时擅自宣讲，在祂完全教导了门徒关于宣讲的能力之后，立即加上：\BibleQuote{但你们要留在城中，直到领受从上而来的能力}\BibleRef{路 24:49}。事实上，我们若约束自己，不向外言语，只留在我们灵魂的围墙内，那么当我们完全被赋予神圣的能力时，我们就可以从自己里面出来，也教导他人。因此，通过一位智者说：\BibleQuote{青年，在你自己的案件中少说话；若有人问你两次，你的回答应有开头}\BibleRef{德 32:10}。因此，我们同一的救赎主，虽然在天上原是创造者，甚至在祂大能的显示中是天使的老师，却不愿在三十岁之前成为地上人的老师，这是为了给急躁的人注入一种最有益的敬畏之情，因为连祂自己，这位不可能跌倒的，也要等到完美年龄才宣讲完美生活的恩宠。因为经上记载：\BibleQuote{当他十二岁时，孩童耶稣留在耶路撒冷}\BibleRef{路 2:42}。稍后，当祂的父母寻找祂时，又说：\BibleQuote{他们发现他在圣殿里，坐在经师中间，听他们，也问他们}\BibleRef{路 2:46}。因此，需要警醒地思考：当提到耶稣十二岁时坐在经师中间，我们发现祂不是教导，而是提问。这个例子清楚表明，任何软弱的人都不敢冒险教导，因为连那个孩子，凭祂神性的能力向祂的老师本人提供知识的话语，也愿意通过提问来受教。但是，当保禄对他的弟子说：\BibleQuote{这些事你要命令和教导：不要让人轻视你的青春期}\BibleRef{弟前 4:11}时，我们必须理解，在圣经的语言中，青年有时被称为青春期。这点在引用撒罗满的话时更加明显，他说：\BibleQuote{青年，在你青春期要快乐}\BibleRef{训 11:9}。因为除非他用这两个词指同一个意思，否则他不会称一个他正劝诫在青春期的人为青年。\par

% structure: p0163 source=h2 target=subsection reason=relative-heading-rank; base=h2

\subsection{第二十六章}

\Emph{如何劝诫那些凡事顺遂的人和那些事事不顺的人。}\par

(\WorkTitle{劝言} 27.) 对于在现世事物中顺遂如意的那些人，以及虽然贪恋此世之物却被逆境的辛劳所困的那些人，应当以不同的方式加以劝诫。因为对于在现世事物中顺遂如意的那些人，当一切顺遂其心愿时，应当劝诫他们，免得他们将心固定在所赐予的恩惠上，而忽略了寻求赐予者；免得他们爱慕流徙之地胜过爱慕故乡；免得他们将旅途的补给转变为抵达终点的障碍；免得他们沉醉于夜间月光的愉悦，而畏于观望太阳的光辉。因此，应当劝诫他们，将他们在此世所获得的一切视为逆境中的安慰，而非报偿的赏报；反而应使他们以心神对抗世间的恩惠，免得他们完全沉溺于其中而心思全然愉悦。因为凡是在自己内心的判断中，不以对更美好生命的爱慕来克制自己所享有的顺遂的人，就会将此短暂生命的恩惠转变为永久死亡的机会。故此，在\Place{厄东}人的形象下，那些让自己被自己的顺遂所征服的人，那些因此世的成功而欢喜的人，受到谴责，正如经上所说：\BibleQuote{他们以欢乐并全心、全意，把我的地据为己有}\BibleRef{则 36:5}。在这些话中，值得留意的是，他们受到严厉的谴责，不仅是因为他们欢喜，而是因为他们全心、全意地欢喜。因此，\Person{撒罗满}说：\BibleQuote{愚蒙人的背弃必杀其身，愚昧人的顺遂必灭其命}\BibleRef{箴 1:32}。故此，\Person{保禄}劝诫说：\BibleQuote{购买的，要像一无所有；享用这世界的，要像不享用的}\BibleRef{格前 7:30}。因此，凡供应给我们的外在事物，应当只对我们如此有用，以致不将我们的心意由对至高喜乐的渴望转移；这样，那些在我们流徙状态中给我们救助的事物，就不会减少我们灵魂朝圣之旅的哀恸；而我们这些自视因与永恒事物分离而可怜的人，就不会因暂时的事物而欢喜如同幸福。故此，教会藉着蒙选者的声音说：\BibleQuote{他的左手在我头下，他的右手必拥抱我}\BibleRef{歌 2:6}。天主的左手，即现世生命的顺遂，她放在头下，因为她以最高之爱的专注将其压制。但天主的右手拥抱她，因为她以全部的热诚被祂永恒的福乐所环绕。所以，又藉着\Person{撒罗满}说：\BibleQuote{她右手有长寿，左手有富贵荣华}\BibleRef{箴 3:16}。因此，当他说富贵荣华放在左手时，就显示了应以何种态度看待它们。故圣咏作者说：\BibleQuote{求用你的右手拯救我}\BibleRef{咏 107:6}。因为他没有说\Quote{用你的手}，而是说\Quote{用你的右手}；也就是说，在说\Quote{右手}时，表明他所求的是永远的救恩。所以又记着说：\BibleQuote{上主啊，你的右手击碎了敌人}\BibleRef{出 15:6}。因为天主的敌人，虽然在他的左手下顺遂，却被他的右手击碎；因为现世生命大多抬高恶人，而永恒福乐的到来则定他们的罪。\par

应当劝诫那些在此世顺遂的人，要明智地思量：现世生命的顺遂有时是为了诱人进入更美好的生命，但有时却是为了使他们永远受更重的审判。因此，天主的子民被应许了客纳罕地，为使他们有时能被激发去盼望永恒之事。因为那粗野的民族，若没有从应许者那里也获得一些近在眼前的事物，就不会相信天主的遥远应许。因此，为使他们在信德上更坚定地仰望永恒之事，他们不仅因盼望而被引向现实，也因现实而被引向盼望。圣咏作者明确为此作证说：\BibleQuote{他赐给他们外邦人的土地，他们获得了万民的劳苦，为使他们遵守他的法令，寻求他的法律}\BibleRef{咏 104:44}。但是，当人的心灵不以善行回应天主的慷慨恩赐时，它就因被认为曾受良好养育而更公正地被定罪。故此，圣咏作者又说：\BibleQuote{他们被高举时，你便推倒他们}\BibleRef{咏 72:18}。的确，当被抛弃的人不以正义的行为回报天主的恩赐，当他们在此世完全放纵自己，沉溺于丰厚的顺遂中时，他们在外在方面所获得的，正是他们从内在方面堕落的开始。因此，在阴间受痛苦的富翁被说：\BibleQuote{你活着的时候享尽了你的福}\BibleRef{路 16:25}。因此，他虽然是个恶人，却在此世享福，为的是在那里更充分地受苦，因为即使在此世，他在幸福中也没有悔改。\par

但是，另一方面，那些确实贪恋世俗事物，却又因逆境的辛劳而感到疲惫的人，应当受到劝诫，要深思熟虑地思考：万物的创造者与管理者是何等慈爱地眷顾那些祂不交给他们自己欲望的人。因为，医生对之绝望的病人，他允许他随意取用所渴望的一切；但被认为可以治愈的人，却被禁止许多他所贪求的事物；我们也从孩子们那里收回钱财，但同时我们却将全部遗产留给他们作为继承人。那么，那些被暂时的逆境所谦卑的人，应当从永恒产业的希望中获得喜乐，因为天主眷顾若不是为了以纪律的规条来教育他们，就不会约束他们，除非祂计划要他们永远得救。因此，那些在贪恋的世俗事物方面，因逆境的辛劳而感到疲惫的人，应当被劝诫要仔细思考：大多数情况下，即使是义人，当世俗的权力使他们高升时，也会像陷入网罗一样被罪所捕获。因为，正如我们在本书前一部分已经说过的，天主所爱的\Person{达味}，在受奴役时比在他来到王位时更为正直 \BibleRef{撒上24:18}。因为，当他为仆时，出于对正义的爱，他不敢攻击那被交到他手中的仇敌；但当他为王时，因好色的诱惑，他用诡计甚至处死了一位忠心的士兵 \BibleRef{撒下11:17}。那么，谁能无害地寻求财富、权力或荣耀，如果这些东西连那位不寻求而得的人也受到了伤害？如果在这些事物中，连那位由天主选择为之预备的人，也因罪的介入而受扰，那么谁能在不经过巨大斗争的情况下得救呢？他们应当被劝诫思考：\Person{撒罗满}，那位如此大智慧后竟堕入偶像崇拜的人，据记载在他跌倒之前并未遭受任何现世逆境；但那赐予他的智慧完全离开了他的心，因为连最小的磨练纪律也没有保护它。\par

% structure: p0168 source=h2 target=subsection reason=relative-heading-rank; base=h2

\subsection{第二十七章}

\Emph{如何劝诫已婚者和独身者。}\par

（\Emph{第二十八种劝诫}）对已婚者和未婚者的劝诫方式不同。因为，那些受婚姻束缚的人应当被劝诫，在彼此顾念对方的好处时，双方都应当努力取悦伴侣，以免使他们的创造者不悦；他们应当如此处理现世的事务，同时仍不忽略渴望天主的事物；他们应当如此为现世之善而喜乐，同时仍以恳切的忧虑惧怕永罚；他们应当如此为暂时的祸患而悲伤，同时仍以完全的安慰将希望寄托于永恒之善；这样，当他们知道自己所从事的是短暂的，而所渴望的是永恒的时，现世的祸患不会使他们灰心，因他们被天上之善的希望所坚强；现世的美善也不会欺骗他们，因他们因惧怕将来的审判而被忧伤所困扰。因此，基督徒婚姻的心意既是软弱的，又是坚定的：因为它不能完全藐视一切暂时的事物，却又能以渴望联系于永恒的事物。虽然此时它耽于肉身的快乐，但愿它因至高希望的安慰而坚强；如果它拥有现世之物以供旅程之需，愿它希望天主之物作为旅程终点的果实；也不要把自己完全投入目前所从事的，以免它完全失落于所应坚定希望的事物。\Person{保禄}宗徒对此简洁地表达说：\BibleQuote{有妻子的，要像没有一样；哭泣的，要像不哭泣的；欢乐的，要像不欢乐的} \BibleRef{格前7:29}。因为，那有妻子却像没有一样的人，是能通过她享受肉身的安慰，却从未因对她的爱而偏离更好意向的正直，去行恶事。那有妻子却像没有一样的人，是看到万物皆为短暂，不得已忍受肉身的挂虑，却以渴望期盼着灵魂的永恒喜乐。此外，哭泣像不哭泣，是如此哀悼外在的逆境，却仍知道在永恒希望的安慰中喜乐。再者，欢乐像不欢乐，是如此因低下之物而鼓舞，却仍永不停止对天上之物的畏惧。在同一处稍后他又恰当地补充说：\BibleQuote{因为这世界的格局正在逝去} \BibleRef{格前7:31}；好像他明确地说：不要持久地爱世界，因为你们所爱的世界本身不能持久。你们枉然将情感寄托于它，如同它是永恒的，而你们所爱的本身正在消逝。丈夫和妻子应当被劝诫，他们彼此有时使对方不悦的那些事，要以相互的忍耐承受，并以相互的劝勉来补救。因为经上记载：\BibleQuote{你们要彼此分担重担，如此你们就满全了基督的法律} \BibleRef{迦6:2}。因为基督的法律就是爱德；因为祂已丰沛地将祂的美善赐予我们，并忍耐地承受了我们的恶。因此，当我们以效法来履行基督的法律时，就是既仁慈地施予我们的美善，又虔诚地承受朋友们的恶。他们也应受劝诫，每人都要注意，并不在于自己从对方受了什么，而在于对方从自己受了什么。因为如果一个人想到自己使对方承受了什么，就更能轻省地忍受自己从对方所受的。\par

丈夫与妻子应受劝诫，要记住他们是为生育子女而结合的；当他们沉迷于过度的房事，将生育的机会转为享乐的服务时，要想到即便他们没有越出婚姻之外，却仍在婚姻本身内超过了婚姻的正当义务。因此，他们需要藉着频繁的祈求，清除他们因掺杂了享乐而玷污的婚姻结合的美善形式。因为正是为此，擅长天上医药的宗徒，当他说\BibleQuote{男人不亲近女人倒好，但为了避免淫乱，男人当各有自己的妻子，女人当各有自己的丈夫}\BibleRef{格前 7:1}时，与其说是为众人规定了一种生活方式，不如说是为软弱者指出了补救方法。因为他预先提出对淫乱的戒惧，确实不是给站得住的人一条诫命，而是给要跌倒的人指出一张床，免得他们或许跌倒在地。因此他给仍然软弱的人加上说：\BibleQuote{丈夫当将妻子应得的给予妻子，妻子也同样给予丈夫}\BibleRef{格前 7:3}。并且，在至为尊贵的婚姻身份中容许他们一些享乐的同时，他加上说：\BibleQuote{但是我说这话，是出于迁就，并不是出于命令}\BibleRef{格前 7:6}。那么，哪里提到迁就，自然就暗示了过失；但这过失之所以更易被宽恕，在于它不在于做那非法的事，而在于没有把合法的事控制在限度内。此事由\Person{罗特}其人亲身表达得很好，当他逃离燃烧的\Place{索多玛}，却找到\Place{左哈尔}，未再登上高山之巅。因为逃离燃烧的索多玛，就是要避免肉体之非法欲火。但高山之巅，是节制者的纯洁。或者，至少可以说，那些人仿佛在山上，他们虽然黏着于肉体的交合，但仍没有超出为生育子女而应有的结合，并未放纵地迷失在肉体的享乐中。因为站在山上，就是在肉体中除了生育的果实别无所求。站在山上，就是不使自己以属肉体的方式黏着于肉体。但是，由于有许多人确实放弃了肉体的罪，然而在婚姻状态中，却没有只遵守应尽的房事义务，所以\Person{罗特}确实走出了\Place{索多玛}，却尚未立即抵达高山之巅；因为可诅咒的生活已经放弃了，但婚姻节制的崇高境界仍未完全达到。但中间有\Place{左哈尔}城，以拯救软弱的逃亡者；也就是说，当已婚者彼此有房事甚至不加节制时，他们仍避免陷入罪，仍因仁慈而得救。因为他们仿佛找到了一座小城，得以在火中受保护；因为这种婚姻生活固然在德行上并不惊人，但尚可免受惩罚。因此同一位\Person{罗特}向天使说：\BibleQuote{看，这座城离的很近，可以逃到那里，那是一座小城，请容我逃到那里；那里不是一座小城吗？我可保全性命}\BibleRef{创 19:20}? 因此，它被称为近，却又被说成安全的避难所，因为婚姻生活既没有远离世界，也没有与救恩的喜乐隔绝。但是，婚姻中的人，在这种行为方式中，当他们藉着不断的祈求彼此代祷时，他们就如在一座小城中保存了自己的性命。因此天使对同一位\Person{罗特}也正确地说：\BibleQuote{好罢！我也在这事上俯听你，不消灭你提及的这座城}\BibleRef{创 19:21}。因为实在，当祈求向天主倾注时，这样的婚姻生活绝不至受谴责。关于这种祈求，\Person{保禄}也劝诫说：\BibleQuote{你们不要彼此亏负，除非两相情愿，暂时分房，为专务祈祷}\BibleRef{格前 7:5}。\par

但是，另一方面，那些不受婚姻束缚的人应受劝诫，因为他们既无肉身结合的轭使他们俯就世俗的挂虑，就应更加密切地遵守天上的规诫；既然他们免去了婚姻的合法负担，就绝不要让世俗忧虑的非法重量压住他们；愿末日来临之时，他们因较少累赘而更加准备好；以免他们因自由而有能力行更善之事却忽略不做，因而受更重惩罚。让他们听宗徒如何培养某些人去领受独身的恩宠时，并未轻视婚姻，而是防范婚姻所产生的世俗挂虑，说：\BibleQuote{我说这话是为你们的益处，并不是要给你们设下圈套，而是为那合宜的事，并使你们能无阻碍地事奉主}\BibleRef{格前 7:3}。因为婚姻产生世俗的焦虑；因此外邦人的导师劝勉他的听众趋向更善之事，以免他们被世俗焦虑束缚。那么，那独身而受世俗挂虑阻碍的人，即使没有使自己受婚姻约束，仍未逃脱婚姻的重担。应劝诫独身者不要以为与无配偶的女子交往而无受定罪判决之虞。因为当保禄将邪淫的恶行列在那么多可憎的罪恶中时，指出了这罪过，说：\BibleQuote{无论是邪淫的、崇拜偶像的、通奸的、作娈童的、好男色的、偷窃的、贪婪的、酗酒的、辱骂的、勒索的，都不能承继天主的国}\BibleRef{格前 6:9}。又说：\BibleQuote{因为邪淫的和通奸的，天主必要审判}\BibleRef{希 13:4}。因此应劝诫他们，如果他们在诱惑的风暴中危及得救，就应寻求婚姻的港口。因为经上记载：\BibleQuote{结婚比燃烧更好}\BibleRef{格前 7:9}。事实上，他们若不曾许下更善的愿，便无可指责地进入婚姻。因为凡是立志追求更大善功的人，便使先前合法的较小善功成为非法。因为经上记载：\BibleQuote{手扶着犁而往后看的，不适于天主的国}\BibleRef{路 9:62}。因此，凡曾专心致志于更坚决决意的人，若离开更大的善而退回到最小的善，便被证为往后看。\par

% structure: p0173 source=h2 target=subsection reason=relative-heading-rank; base=h2

\subsection{第二十八章}

\Emph{如何劝诫那些曾经历肉身之罪的人，以及那些未曾经历的人。}\par

（\Emph{劝诫} 29.）应以不同方式劝诫那些自知有肉身之罪的人，与那些不知有罪的人。因为那些曾经历肉身之罪的人应受劝诫，至少要在船难之后畏惧海洋，并在得知丧亡危险后感到恐怖；以免在天主仁慈地保存他们于已行恶事之后，因邪恶地重复同样的事而死去。因此，对那犯罪且不停止犯罪的灵魂说：\BibleQuote{娼妓的额头在你身上，你拒绝羞惭}\BibleRef{耶 3:3}。因此应劝诫他们留心，好使若他们拒绝保全所领受的自然恩惠之完整，至少在被撕裂后予以修补。他们务必思考，在如此众多的信徒中，有多少人不仅保持自身纯洁，还感化他人离开邪道。那么，当别人站立于正直中时，他们自己即使在失落之后也不归向更好的心志，他们能说什么？当许多人带着别人一同进入天国时，他们甚至连自己也不带回给等待他们的主，他们能说什么？应劝诫他们思念过去的过犯，并逃避将要来到的过犯。因此，主藉先知以犹太人的形象，将已往的罪恶记忆唤醒那些在这世上被败坏了的灵魂，好使她们羞耻于在未来的罪中受玷污，说：\BibleQuote{她们在埃及行淫，在幼年时行淫，那里有人压迫她们的乳房，揉搓她们处女的乳头}\BibleRef{则 23:3}。因为乳房在埃及受压，是指人灵魂的意志向此世的卑劣欲望卖淫。处女的乳蒂在埃及被揉搓，是指自然的感觉仍完整时，被来袭的贪欲之败坏所腐化。\par

那些曾经验过肉身的罪的人，应当被劝诫，要警醒地观察天主以何等伟大的仁慈向我们敞开祂怜悯的怀抱，如果我们过犯之后归向祂，当祂通过先知说：\BibleQuote{\Emph{如果一个男子休了他的妻子，她离开他而成了别人的妻子，他还能回到她那里去吗？那女子难道不是被玷污、被污染了吗？但你和许多爱人行了淫；仍要回到我这里来，主说。}}\BibleRef{耶 3:1}。所以，关于那行了淫并被离弃的妻子，公义的论证被提出；然而对我们，跌倒后归向的，所展示的不是公义，而是怜悯。由此我们当然应该明白我们的邪恶是多么大，如果我们即使过犯之后也不归向，因为我们犯罪时，竟被如此大的怜悯所宽恕；或者对于恶人，将从那位在我们犯罪后仍不停召唤我们的主那里得到怎样的宽恕。确实，这种在过犯后召唤的慈悲，通过先知得到了很好的表达，当时对背离天主的人说：\BibleQuote{\Emph{你的眼睛必看见你的导师，你的耳朵必听见背后有人劝诫你。}}\BibleRef{依 30:20}。的确，主曾当面劝诫人类，当祂对那受造于乐园、立于自由意志中的人，宣告祂当作和不当作的事。但人在骄傲中轻视祂的命令时，就把背转向了天主的面。然而天主并未在人的骄傲中抛弃他，因为祂赐下法律是为召回人，并差遣劝诫的天使，且祂亲自以我们有死性的肉身显现。因此，站在我们背后，祂劝诫我们，因为即使被轻视，祂仍召唤我们恢复恩宠。所以，凡能普遍适用于所有人的，也必须特别适用于每个人。因为每个人听到天主劝诫的话语，好像是摆在他面前，当他在犯罪之前，知道祂旨意的规诫。因为仍站在祂面前还不是犯罪轻慢祂。但是，当一个人舍弃无辜之善，且故意选择邪恶时，他就把背转向了天主的面。但是看，甚至在他背后，天主跟随并劝诫他，因为甚至在犯罪之后，祂仍说服他归向祂。祂召回那背离的人，祂不念旧恶，祂为归向者敞开怜悯的怀抱。因此，我们听从背后劝诫我们的声音，如果我们至少在罪过后归向那邀请我们的主。我们应当为那召叫我们的主的怜悯感到羞愧，如果我们不畏惧祂的公义：因为轻慢祂是更严重的邪恶，因为虽然被轻慢，祂仍不嫌弃继续召唤我们。\par

但是，另一方面，那些不熟悉肉身的罪的人，应当被劝诫，要更加焦虑地恐惧骤然堕落，因为他们站在更高的地方。他们应当被劝诫：他们站立的地方越突出，埋伏者的箭矢也就越频繁地攻击他们。因为埋伏者通常更加激烈地兴起，因为他看到自己被更坚强地击败；并且他越是轻视和感到被击败难以忍受，因为他看到软弱肉身的破营地摆阵对抗他。他们应当被劝诫，要不断仰望赏报，那么他们无疑会欣然践踏他们所忍受的诱惑的劳苦。因为，如果注意力固定在已经获得的幸福上，而不顾通往幸福的过程，那么过程的艰辛就变得轻省了。让他们听先知所说的话：\BibleQuote{\Emph{主这样对阉人说：凡守我的安息日，选择我所喜悦的事，坚守我的盟约的，我必在我的殿内，在我的墙内，赐给他们比子女更好的地方和名号。}}\BibleRef{依 56:4}。因为他们确实是阉人，他们抑制肉身的冲动，在自己内割断了对邪行的迷恋。此外，他们在圣父面前被看重的地位被显示出来，因为在圣父的家中，也就是圣父的永恒住所里，他们甚至被置于子女之前。让他们听通过若望所说的话：\BibleQuote{\Emph{这些人没有与女人玷污，因为他们是童贞者，并且跟随羔羊，无论祂到哪里去。}}\BibleRef{默 14:4}；以及他们如何唱一首歌，除了那十四万四千人以外，没有人能唱。因为单独向羔羊唱一首歌，就是因肉身的朽坏不败而永远超越所有忠信者与祂一同喜乐。然而其余的选民可以听见这首歌，虽然他们不能唱，因为通过爱德，他们为那些人的提升而喜乐，尽管他们没有上升至他们的赏报。让那些不熟悉肉身的罪的人听真理本人论及这贞洁所说的话：\BibleQuote{\Emph{不是所有人都接受这话。}}\BibleRef{玛 19:11}。祂称这事为至高的，因为祂说它不属于所有人；并且，在预言领受这话很困难时，祂向听众表明，当领受后，应以何等谨慎持守它。\par

因此，那些不熟悉肉身的罪的人，应当被劝诫，既要认识到童贞超越婚姻，却又不要把自己抬高到已婚者之上：为的是，当他们把童贞置于首位，而把自己置于末位时，他们既能持守自己所珍视为最好的，同时也能谨守自己，不虚妄地自高自大。\par

应劝诫他们思考：通常节制者的生活会被世俗之人的行为置于羞愧之中，因为后者承担了超出其本分的工作，而前者却不激励自己的心志去达到自身秩序的标竿。因此，先知说得很好：\BibleQuote{漆冬，你要羞愧，因为大海说}\BibleRef{依 23:4}。因为漆冬仿佛被大海的声音羞辱了，那坚固而稳定之人的生活，与世俗而漂泊于世之人相比，受到了责斥。因为时常有些人，在肉身的罪过后回归主，看到自己因恶行更易受罚，便在善工上显得更为热切；而另一些坚持肉身纯洁的人，因过去少有可哀叹之事，就认为生命的纯正足以自足，不激发自己以任何热忱的激励去追求精神的热火。大多数情况下，罪后燃烧着爱的生活，比安全中麻木的纯洁更令天主喜悦。因此，审判者的声音也说：\BibleQuote{她的那许多罪得了赦免，因为她爱得多}\BibleRef{路 7:47}；以及\BibleQuote{对于一个人悔改，在天上所有的欢乐，甚于对那九十九个无须悔改的义人}\BibleRef{路 15:7}。如果我们衡量自己内心的判断，就更能从经验本身领悟这一点。因为我们爱一块地，它在除掉荆棘后产生丰硕果实，胜过一块从未长过荆棘但耕作后收成荒芜的地。那些不知肉身之罪的人应受劝诫，不要因自身高级秩序的崇高而自认为优于他人，因为他们不知道比自己低微的人所做的事比他们更好。因为在公义审判者的审问中，行为的品质会改变秩序的功绩。因为谁在思考事物外表时能不知道，在宝石的本性中，红宝石优于风信子石？但蔚蓝色的风信子石仍优于苍白的红宝石；因为前者之美为其提供了自然秩序所拒绝的东西，而后者虽被自然秩序优先，却被颜色品质所贬低。因此，在人类中，有的在较好秩序中反而较差，有的在较差秩序中反而较好；因为这些人通过善生超越了低等状态的命运，而那些人的行为未能达到高等位置而减损了其功绩。\par

% structure: p0180 source=h2 target=subsection reason=relative-heading-rank; base=h2

\subsection{第二十九章}

\Emph{如何劝诫那些哀叹行为之罪的人，以及那些仅哀叹思想之罪的人。}\par

（第三十项劝诫）哀叹行为之罪的人与哀叹思想之罪的人，应被不同地劝诫。因为哀叹行为之罪的人应受劝诫：完美的哀号应洗除掉已完成的恶事，免得他们所负的已完成行为的债，大于他们在补赎中因眼泪所偿还的。因为经上记载：\BibleQuote{他按量给我们喝下眼泪}\BibleRef{咏 78:6}；这就是说，每个灵魂在悔改时，应依自己记得曾因罪而干涸于天主之前，饮下多少痛悔之泪。应劝诫他们不断将过去的过犯置于眼前，如此生活，以使这些过犯不必被严厉的审判者看见。\par

因此，\Person{达味}祈祷时说：\BibleQuote{求你转面不看我的罪过}\BibleRef{咏 50:3}，他稍前也说：\BibleQuote{我的罪愆常在我眼前}\BibleRef{咏 50:3}；仿佛在说：我恳求你不要看我的罪，因为我自身不停止地看它。同样，主通过先知说：\BibleQuote{我必不再记念你的罪过，但你却要记念它们}\BibleRef{依 43:25}。应劝诫他们细数所有过去的过犯，在哀悼从前流浪的每个污点时，同时用眼泪洁净整个自我。因此，先知\Person{耶肋米亚}在考虑\Place{犹大}的种种过犯时，说得很好：\BibleQuote{我的眼睛流下溪水般的眼泪}\BibleRef{哀 3:48}。因为我们从眼中流下分流的泪水，就是将分别的眼泪赋予我们各自的罪过。因为心灵不能同时等同地为一切事物悲伤；但每当记忆如今更尖锐地触动此罪或彼罪，在每一个罪中为所有罪过所动，便同时从所有罪过中被洁净。\par

应劝戒他们，要依靠他们所渴求的仁慈，以免因过度痛苦而灭亡。因为主若愿意亲自以严厉的惩罚击打他们，就不会将罪摆在犯罪者眼前令其哀哭。显然，祂愿意将那些祂预先使其成为自身审判者的人，隐藏在祂的审判之外。因此经上记载：\BibleQuote{让我们在告解中提前来到主面前}\BibleRef{咏 94:2}。因此通过保禄说：\BibleQuote{如果我们自己审判自己，就不致受审判}\BibleRef{格前 11:31}。同样，应劝戒他们，要满怀希望地信赖，以免因粗心大意而陷入麻木的安全感。因为狡猾的仇敌常见被他用罪绊倒的灵魂因堕落而痛苦，便用有害的安全感的甜言蜜语诱惑它。这事在狄纳的史事中形象地表达了出来。因为经上记载：\BibleQuote{狄纳出去观看那地的女子；当地的首领、希未人哈摩尔的儿子舍根看见她，就爱上她，抓住她，与她同寝，玷辱了她；他的心依恋着她，当她忧愁时，就用温和的甜言蜜语安慰她}\BibleRef{创 34:1}。的确，狄纳出去观看外邦女子，代表任何灵魂忽略自己的事，留意他人的行为，从而背离自身的状态和秩序。而当地的首领舍根制服了她，因为魔鬼在灵魂忙于外在挂虑时发现并败坏它。\BibleQuote{他的心依恋着她}，因为他视灵魂因不义而与他结合。因为当灵魂意识到自己的罪时，它被定罪，并渴望哀叹自己的过犯；但败坏者却将虚空的希望和安全的理由摆在它眼前，以剥夺它悲伤的益处，因此经文中正确添加了：\BibleQuote{当她忧愁时，就用甜言蜜语安慰她}。因为他时而提及他人更重的过犯，时而说所行之事无关紧要，时而又说天主仁慈；或者又承诺日后尚有悔改的时间；这样，灵魂被这些欺骗所迷惑，便悬而不定，不能坚定悔改的意向，以致日后一无所获，如今也不因恶而悲伤，到那时因如今甚至以过犯为乐，便更将被惩罚彻底淹没。\par

但另一方面，应劝戒那些哀哭思想之罪的人，要在灵魂深处焦虑地思考：他们只是在快感中犯了罪，还是在同意中也犯了罪。因为内心通常受到诱惑，在肉体的罪恶中体验到快感，然而在其判断中却抵抗这同样的罪恶；因此，在思想的隐秘中，他既因所悦之事而忧伤，又因所忧之事而喜悦。但有时灵魂如此被诱惑的深渊淹没，以致完全不抵抗，反而故意随从那借着快感攻击它的事物；如果有外在的机会，它很快就将内在的愿望实现为行动。当然，若按照严厉审判者的公正惩罚来看，这罪不是思想的罪，而是行为的罪；因为，即使环境的迟缓在外部推迟了犯罪，意志却通过同意的行为在内部完成了犯罪。\par

此外，我们从原祖的事例中得知，我们以三种方式犯下各种罪的不义：即，在诱惑、快感和同意中。第一种通过仇敌施行，第二种通过肉体，第三种通过精神。因为埋伏者暗示错误之事；肉体顺从于快感；最后，被快感战胜的精神表示同意。因此，那蛇暗示了错误之事；然后厄娃，如同肉体，顺从于快感；但是亚当，如同精神，被诱惑和快感征服，表示了同意。因此，通过诱惑我们得知罪，通过快感我们被征服，通过同意我们也被捆绑。因此，应劝戒那些哀哭思想之罪的人，要焦虑地思考他们堕入罪恶到了何种程度，以便他们能根据内心所意识到的堕落程度，以相应的哀哭程度被提升；免得，如果所思想的罪恶过于轻微地折磨他们，便会引导他们甚至去付诸行动。但在这一切中，他们应受到如此的警醒，以致仍应毫不沮丧。因为仁慈的天主常常更迅速地赦免内心的罪，因祂不允许它们化为行动；所思想的不义之所以被更迅速地解除，是因为没有被付诸行为所过于紧紧地束缚。因此圣咏作者正确地说：\BibleQuote{我说：我要向上主承认我的过犯；你就赦免了我内心的不敬}\BibleRef{咏 30:5}。因为他添加了\Quote{内心的不敬}，表明他所要承认的是思想的过犯；而他说\Quote{我说：我要承认}，随即补充\Quote{你就赦免了}，显示了在这种情况下宽恕是何等容易。因为，当他只承诺要请求时，就已获得了他承诺要请求的；这样，既然他的罪没有进展到行为，他的悔改也不必进展到折磨；那思想的忧伤应洁净灵魂，其实不过是思想的不义玷污了它。\par

% structure: p0187 source=h2 target=subsection reason=relative-heading-rank; base=h2

\subsection{第三十章}

\Emph{如何劝戒那些不戒除自己所痛哭的罪的人，以及那些戒除了罪却不痛哭的人。}\par

(\(\\Emph{规劝}\) 31.) 那些哀叹自己的过犯却不肯改过的人，与那些改过却不肯哀叹的人，应用不同的方式加以规劝。因为那些哀叹自己的过犯却不肯改过的人，应当被规劝要焦虑地学习：他们若在哭泣中洁净自己，却在生活中邪恶地玷污自己，那他们的哭泣就是枉然的，因为他们流泪洗刷自己的目的，本是在洁净后不再回到污秽中。因此经上记载：\(\\BibleQuote{狗转身吃自己所吐的，猪洗净了，又回到污泥中打滚}\) \(\\BibleRef{伯后 2:22}\)。狗呕吐时，确实吐出压迫胃的食物；但当它回到呕吐物时，又再次被已经摆脱的东西所压。同样，那些哀叹自己过犯的人，确实通过告解吐出充满罪恶的邪恶，这邪恶曾压迫他们灵魂的深处；然而，告解后又回到其中，便再次把它吞下。而猪洗净后在污泥中打滚，变得更脏。一个哀叹过去的过犯却不肯改过的人，使自己陷入更严重的罪的惩罚中，因为他既轻视那本可由哭泣获得的宽恕，又好像在脏水中打滚；因为他不让自己的哭泣带来纯洁的生活，甚至使自己的眼泪在天主眼中成为污秽。因此经上又记载：\(\\BibleQuote{在祈祷中不要重复言词}\) \(\\BibleRef{德 7:14}\)。在祈祷中重复言词，就是在哀悼之后，又犯下需要重新哀悼的事。因此通过依撒意亚说：\(\\BibleQuote{你们要洗濯，自洁}\) \(\\BibleRef{依 1:16}\)。因为谁在洗濯后保持不了生活的纯洁，就是忽略了自洁。因此，那些不断哀悼自己所犯的罪，却又重复犯下需要哀悼的事的人，他们虽被洗濯，却绝不洁净。因此通过一位智者说：\(\\BibleQuote{谁因触摸死者身体而受洗，后又触摸它，他的洗濯有何益处？}\) \(\\BibleRef{德 34:25}\)。的确，那因触摸死者身体而受洗的人，是通过哭泣从罪中洁净的人；但他在受洗后又触摸死者身体，就是在哭泣之后重复犯罪。\par

那些哀叹过犯却不肯改过的人，应当被规劝要承认自己在严格的审判者眼中，如同某些人来到某人面前时，极谦卑地奉承他们，但离开后却恶毒地尽力对他们施加敌意和伤害。因为为罪哭泣，不就是向天主表现自己虔诚的谦卑吗？而哭泣后又作恶，不就是对那被恳求的主施行傲慢的敌意吗？\(\\Person{雅各伯}\) 证实了这一点，他说：\(\\BibleQuote{谁若愿意作世界的朋友，就成了天主的仇敌}\) \(\\BibleRef{雅 4:4}\)。那些哀叹自己的过犯却不肯改过的人，应当被规劝要焦虑地思考：通常坏人被痛悔所吸引，却无益于正义，正如通常好人受诱惑犯罪，却不受伤害。这里确实发现了内心倾向与功过要求相应的奇妙尺度：坏人在做某些善事却未完成时，在完全做出邪恶当中却骄傲地自信；而好人受到邪恶诱惑却毫不赞同，因软弱而动摇，却因此通过谦卑更稳固地迈向正义。\(\\Person{巴郎}\) 望着义人的帐幕说：\(\\BibleQuote{愿我的灵魂像义人一样死去，愿我的结局像他们一样}\) \(\\BibleRef{户 23:10}\)。但当痛悔的时刻过去后，他却对那些他曾希望自己甚至在死时也像他们的人的生命出谋献策；当他找到满足贪婪的机会时，他立刻忘记了他曾为自己祈求的一切纯洁。因此，外邦人的师傅与宣讲者 \(\\Person{保禄}\) 说：\(\\BibleQuote{我发觉在我的肢体中另有一条法律，与我理智的法律交战，并使我隶属于那在我肢体中的罪恶的法律}\) \(\\BibleRef{罗 7:23}\)。他受诱惑确实正是为了这个目的：为能因认识自己的软弱而更坚定地被确认在善中。那么，为何一个被痛悔触动，却不接近正义；另一个受诱惑，却不受玷污？岂不显然是因为：未完成的善事无助于恶人，而未完成的恶事也不定好人的罪吗？\par

另一方面，那些改过却不肯哀叹的人，应当被规劝不要以为罪已经赦免了，因为他们虽不再以行动增加罪，却仍不以任何哀哭来清除它们。写字的人停止书写，并没有擦去已写的字，因为他没有再添写；侮辱他人的人仅仅沉默，并不能满足对方，因为他必须用后来谦卑的话来驳斥先前骄傲的话；欠债的人不增加债务并不等于被豁免，除非他也偿还所欠的。同样，当我们得罪天主时，绝不因停止作恶而满足，除非我们也以相应的哀叹来追补我们所贪恋的快乐。因为如果今生没有行为上的罪玷污我们，那么我们的清白本身也绝不能保证我们在此世的安全，因为仍有许多不法之事攻击我们的心。既然如此，一个犯过不义、自己见证自己不无罪的人，怎能安心呢？\par

因为天主并非因我们的痛苦而得到饱饫；而是他以相反的药医治我们过犯的疾病，使我们这些因享乐而离开他的人，能因眼泪而苦涩地回到他身边；使我们因放纵于不法之事而跌倒，却能因在合法之事上也克制自己而站起；使被狂喜淹没的心，能因健康的悲伤而被烧净；使骄傲的昂首所伤的，能被谦卑生活的沮丧所医治。因此经上记载：\BibleQuote{我曾对恶人说：不要行恶；对悖逆的人说：不要高举角。}\BibleRef{咏 74:4}。因为悖逆的人若在知道自己的不义之后仍不谦卑悔改，就是高举角。因此又说：\BibleQuote{天主不轻视破碎和谦卑的心。}\BibleRef{咏 50:17}。因为凡哀悼自己的罪却不肯离弃的人，确实破碎了心，却轻视谦卑；但离弃了罪却不为罪哀悼的人，的确已经谦卑了心，却拒绝破碎它。因此保禄说：\BibleQuote{你们从前也确是那样，但你们已经洗净了，已经圣化了。}\BibleRef{格前 6:11}；因为，事实上，改过的生命圣化了那些借着悔改，以泪水的痛苦洗涤所洁净的人。因此，伯多禄看到有些人因想起自己的恶行而惊惶时，劝诫他们说：\BibleQuote{你们悔改吧，并每人受洗。}\BibleRef{宗 2:38}。因为，当他即将谈论圣洗时，他先提到了悔改的哀恸；他们应先在自己痛苦的泪水中沐浴，然后再在圣洗圣事中洗洁。那么，那些忽略为自己过去的恶行哭泣的人，怎能心安理得地生活，确信罪得赦免呢？连教会的大牧者本人也认为，即使是在这主要消除罪恶的圣事中，也必须加上悔改。\par

% structure: p0193 source=h2 target=subsection reason=relative-heading-rank; base=h2

\subsection{第三十一章}

\Emph{如何劝诫那些称赞自己明知不合法的行为的人，以及那些谴责不法行为却毫不防备的人。}\par

（\Emph{劝诫} 32.）当以不同的方式劝诫那些甚至称赞自己所行不法之事的人，以及那些谴责错误之事却不去避免的人。因为那些甚至称赞自己所行不法之事的人，应当被劝诫去思考，他们大多时候用口舌比用行为更得罪人。因为行为只在他们自己身上行恶；但口舌却将邪恶传播到所有听众的灵魂中，他们借着称赞将恶事教导听众。因此，应当劝诫他们，如果他们逃避根除恶事，至少应当害怕播种恶事。应当劝诫他们，让他们自己的个别沉沦足以满足他们。又应当劝诫他们，如果他们不怕作恶，至少应羞于被人看到他们是怎样的人。因为通常罪被隐藏时，就会被避免；因为当灵魂羞于被人看到它本不怕成为的样子时，它迟早会羞于成为它怕被人看到的样子。但是，当任何恶人无耻地招引注意时，他越是自由地行一切恶事，就越是进而认为它是合法的；在他认为合法的事上，他无疑越陷越深。因此经上记载：\BibleQuote{他们宣告自己的罪如索多玛，并未隐藏。}\BibleRef{依 3:9}。因为，如果索多玛隐藏了她的罪，她仍会犯罪，但在恐惧中。但她完全失去了恐惧的约束，因为她甚至不为自己的罪寻求黑暗。因此又记载：\BibleQuote{索多玛和蛾摩拉的呼喊增加了。}\BibleRef{创 18:20}。因为带着声音的罪是在行为中的罪责；但甚至带着呼喊的罪是在放纵中的罪责。\par

但是，另一方面，那些谴责错误之事却不去避免的人，应当被劝诫认真思考，在天主严格的审判面前，他们能为自己说什么辩护，既然他们自己作为审判官也不能免除他们罪行的罪责。那么，这些人不是自己的控告者吗？他们发声反对恶行，却把自己作为有罪者交出。应当劝诫他们去察觉，即使在现在，审判的隐蔽报应如何临到：他们的心智被光照，看见自己所行的恶，却不愿努力克服它；以致于他们看得越清楚，就败亡得越惨；因为他们既感知了理解的光明，又不肯放弃行恶的黑暗。因为，当他们忽略了为帮助他们而赐予的知识，他们就把它变成了反对自己的证据；并且从他们确实领受了以消除自己罪过的理解之光中，他们加重了自己的惩罚。事实上，他们的这种邪恶，行出了自己所谴责的恶，在此世已经尝到了未来审判的滋味；因此，当他们仍要受永罚时，在此世他们自己的审判中却不得释放；以至于他们在彼世要经历更严厉的折磨，因为他们在此世不放弃连自己也谴责的恶。因此，真理说：\BibleQuote{那个知道主人旨意，却未预备，也不照他旨意行的仆人，将受许多鞭打。}\BibleRef{路 12:47}。因此，圣咏作者说：\BibleQuote{愿他们活生生地坠入地狱。}\BibleRef{咏 54:15}。因为活人知道并感觉周围发生的事；死人却毫无感觉。因为如果他们在无知中行恶，就会死人般地坠入地狱。但当他们知道什么是恶，却仍去行时，他们就活生生地、悲惨地、有感觉地坠入罪恶的地狱。\par

% structure: p0197 source=h2 target=subsection reason=relative-heading-rank; base=h2

\subsection{第三十二章}

\Emph{如何劝诫那些因一时冲动而犯罪的人和那些蓄意犯罪的人。}\par

（\Emph{劝诫}33）. 对受突发情欲所胜者与存心犯罪者应作不同的劝诫。受突发情欲所胜者应受劝诫，要视自己每日处于现世生活的战斗中，并以忧虑恐惧的盾牌保护心，因为后者无法预见创伤；要警惕埋伏之敌的暗箭，在如此黑暗的争战中，以不断的注意守护灵魂的内营。因为，若心灵缺乏谨慎的关注，便暴露于创伤之下；诡诈的敌人攻击胸膛，越是自由，越是他发现胸膛没有预谋的护甲。受突发情欲所胜者应受劝诫，不要过于关心尘世事物；因为当他们过度纠缠于短暂事物时，便意识不到刺穿他们的罪的箭。因此，所罗门也表达了一个被打伤却仍沉睡者的言语，他说：\BibleQuote{他们打我，我却未觉疼痛；他们拖我，我竟未察觉。我几时醒来，再寻酒喝？}\BibleRef{箴 23:35}？因为灵魂因疏忽关注而沉睡，被打伤而不觉疼痛，因为它既不能预见将来的灾祸，也不觉察自己已犯的罪。它被拖走，却毫无感觉，因为它被恶习的诱惑引导，却未惊醒守护自己。但它又想醒来，好再寻酒喝；因为尽管被麻木的沉睡压得无法自守，它仍努力醒着去关心世事，好常常沉醉于享乐；当它应当明智醒着的事情上沉睡时，却渴望醒着去做另一件本可赞美地沉睡的事情。因此前面写道：\BibleQuote{你必如在海中沉睡，又如掌舵者被摇撼，松开了舵柄。}\BibleRef{箴 23:35}。因为他在海中沉睡，意味着置身于今世诱惑之中，却忽略警戒如浪涛般涌来的恶习冲动。而掌舵者仿佛松开了舵柄，当心灵失去了引导身体这艘船的恳切关注时。因为，在海上松开舵柄，就是在今世的风暴中放弃预谋的专心。因为，如果掌舵者焦虑地紧握舵柄，他有时在波涛中直冲向前，有时斜切突来的大风。同样，当心灵警醒地引导灵魂时，它有时克服并践踏某些事物，有时谨慎地避开其他事物，以便既因艰苦努力克服眼前的危险，又因预见为未来的斗争积蓄力量。因此，关于天国刚勇的战士又说：\BibleQuote{每人腰间佩刀，因夜间有惊恐。}\BibleRef{歌 3:8}。因为刀佩在腰间，意味着肉体的邪恶建议被圣道宣讲的利刃所制服。而夜间表达了我们软弱的盲目；因为在夜间看不到迫在眉睫的反对。因此，每人腰间佩刀，因夜间有惊恐；即，因为圣人惧怕他们看不见的事物，总是准备好斗争的紧张。因此，又对新娘说：\BibleQuote{你的鼻子如同黎巴嫩的高塔。}\BibleRef{歌 7:4}。因为我们通常用嗅觉预知眼睛看不见的事物。用鼻子，我们也能辨别香味和臭味。那么，教会的鼻子除了圣人的先见之明，还能象征什么呢？它也被说成像黎巴嫩的高塔，因为他们的辨别性预见如此高耸，以至于甚至能看见将来临的诱惑斗争，并在它们来临时坚立防御。因为，预见到未来的事情，在它们成为现实时力量较小；因为当每个人都更加准备好面对打击时，原以为出奇不意的敌人，因被预料而削弱了。\par

但另一方面，那些存心自陷于罪的人，应受劝诫，要谨慎地思考：当他们凭自己的判断作恶时，他们为自己燃起了更严厉的审判；且他们越是因深思熟虑而牢牢被罪捆绑，将来受的惩罚就越沉重。假如他们只是因仓促而犯罪，或许能更快地藉补赎洗清过犯。因为存心而牢牢绑住的罪，不易迅速解开。若非灵魂完全轻视永恒之事，它就不会存心在罪中丧亡。由此可见，存心丧亡的人与因仓促跌倒的人的区别在于：前者因犯罪而失去正义状态时，多半也会陷入绝望的陷阱。因此，主藉先知责备的，与其说是仓促的恶行，不如说是犯罪的意图，说：\BibleQuote{免得我的忿怒如火烧起，燃烧不止，无人能熄灭，因为你们的恶计}\BibleRef{耶 4:4}。又，祂在怒气中说：\BibleQuote{我必照你们计谋的果实惩罚你们}\BibleRef{耶 23:2}。既然存心所犯的罪与其他罪不同，主责备的是恶计而非恶行。因为行为上的罪常出于软弱或疏忽，但计谋中的罪则常出于恶意。相反，先知描述有福的人时说得很好：\BibleQuote{他不坐在瘟疫的座位上}\BibleRef{咏 1:1}。因为座位通常是审判官或主席所坐的。坐在瘟疫的座位上，就是以审判的姿态行不义；坐在瘟疫的座位上，就是以理智辨别邪恶，却故意去行。凡因极大的罪恶傲慢而甚至试图用计谋去行恶的人，就是坐在悖逆的座位上。并且，正如因座位尊严而受支持的人，凌驾于站立的人群之上；同样，存心寻求的罪，也超越那些因仓促而跌倒的人的过犯。因此，那些甚至藉计谋自陷于罪的人，应受劝诫：要由此意识到，他们将来必定会受到怎样的惩罚，因为如今他们已成为恶行的首领，而非同伴。\par

% structure: p0201 source=h2 target=subsection reason=relative-heading-rank; base=h2

\subsection{第三十三章}

\Emph{如何劝诫那些犯很小但频繁的过失的人，以及那些虽然避免很小的过失，但有时陷入严重过失的人。}\par

（\Emph{第三十四种劝诫}）那些虽行非法之事却很小、但屡次为之的人，与那些避免小罪、却有时陷入大罪的人，应以不同方式劝诫。屡次犯小过的人，应受劝诫：绝不可只考虑所犯之罪的质，而要考虑其量。因为，如果他们权衡行为后不屑于害怕，那么他们应当在数算时感到惊恐；因为巨大的江河水坝是由无数小水滴汇成的；舱底水悄然上涨，其后果与公开的风暴相同；身上发疹的溃疡虽小，但当无数疹子蔓延时，它和胸部所受的重伤一样能毁灭身体的生命。因此，经上记载：\BibleQuote{轻视小事的人，必渐渐堕落}\BibleRef{德 19:1}。因为那忽略哀叹和避免最小罪恶的人，不是突然，而是一点一点地完全失去正义状态。屡次犯小过的人，应受劝诫要忧虑地思考：有时在小罪中犯罪可能比在大罪中更严重。因为大罪因其很快被承认是罪，所以更快被改正；而小罪因其被看作不算罪，所以更被忽略，以更坏的效果保留在习惯中。因此，多半发生这样的情况：习惯于轻微邪恶的心灵，对重大邪恶也不再恐惧，且因罪而饱足，最终达到一种对不义的认可，并且因学会在无恐惧中犯小罪，而在大事上更不屑于害怕。\par

但是，另一方面，那些远离小罪，却有时陷于重罪的人，应被恳切地劝告认识到自己所处的状态：当他们因防范极小的事而心高气傲时，竟被自己的高傲深渊所吞没，以至于犯下更严重的事；并且，当他们外表上克制了小恶，内心却被虚荣骄傲所膨胀，他们在灵魂内部被骄傲的病态所征服，而外在则陷入更大的恶中。那么，那些远离小过，却有时陷入重罪的人，应被劝告要小心，免得他们在自以为外表站立之处，内心却跌倒；也免得按照严格审判者的报应，因较小的义而骄傲，成了陷入更重罪孽的途径。因为那些虚妄高傲、将守最小的善归于自己力量的人，被正义地遗弃，便被更大的罪所压倒；借由跌倒，他们学到自己的站立并非出于自己，这样，那被最小善举所高举的心，便可能因无边的恶而谦卑。他们应被劝告思考：当他们在更大的过犯中束缚自己于深重的罪债时，他们在所防范的小过犯中往往犯罪更重；因为在前面，他们行恶事，而在后面，他们向人隐藏自己是恶的。因此便发生：当他们在天主面前行更大的恶时，是公开的不义；而当他们在人前谨慎地遵守小小的善事时，是伪装的圣德。因为为此关于法利塞人曾这样说：\BibleQuote{滤出蚊蚋，却吞下骆驼}\BibleRef{玛 23:24}。仿佛清楚地说：最小之恶你辨别，更大之恶你吞噬。为此，他们又被真理之口所责备，当他们对他们说：\BibleQuote{你们将薄荷、茴香、莳萝献上十分之一，却忽略了律法上更重的事，就是公义、怜悯、信实}\BibleRef{玛 23:23}。因为也不可轻忽地听：当祂说最小的事物被献上十分之一时，祂特意提及最卑微的香草，却是芬芳的；这无疑是为了显示：当伪善者遵守小事时，他们为自己追求圣洁声誉的香气；尽管他们忽略了履行最大的事，他们仍遵守那些在人的判断中远近散发芬芳的最小事物。\par

% structure: p0205 source=h2 target=subsection reason=relative-heading-rank; base=h2

\subsection{第三十四章}

\Emph{如何劝告那些甚至不开始善行的人，以及那些开始善行却不完成的人。}\par

（\Emph{劝告} 35.）那些甚至不开始善行的人，与那些开始善行却绝不完成的人，应以不同的方式加以劝告。因为对于甚至不开始善行的人，首要的需要不是建立他们可能健康地爱慕的事物，而是拆除他们错误从事的事物。因为他们不会遵循他们所听来的未经尝试的事，除非他们首先感受到他们所尝试的事物是多么有害；因为一个不知道自己已经跌倒的人不会渴望被扶起；一个感受不到伤口疼痛的人也不会寻求任何医治。首先，应当向他们表明他们所爱的事物是多么虚妄，然后才小心地让他们知道他们所忽略的事物是多么有益。让他们先看到他们所爱的是应逃避的，然后毫无困难地感知到他们所逃避的是应爱的。因为他们更快地接受他们未曾尝试的事物，如果他们确认他们关于所尝试事物所听到的任何论述是真实的。那么，当他们以确定的判断学到他们曾错误地坚持虚妄事物时，他们就学会以充沛的渴望寻求真正的善。因此，要告诉他们：当前的善事将很快从享受中消逝，而关于它们的账目仍将存留，永不消逝，以待惩罚；因为现在既令他们欢喜的事违背他们意愿被夺走，而令他们痛苦的事也违背他们意愿为他们留存，以待惩罚。这样，他们在有害地取乐的事物中，可能被健康地充满恐惧；以致当受创的灵魂，在看见其跌倒的深重毁灭时，感知到自己已到达悬崖边缘，它可能后退一步，并害怕它曾爱慕的，从而学会珍视它曾鄙视的。\par

因为正因此，耶肋米亚被派遣去宣讲时，主对他说：\BibleQuote{看，我今天立你于列国和邦国之上，为要拔除、拆毁、毁灭、分散、建造和栽植}（\BibleRef{耶 1:10}）。因为，除非他先毁灭错误的事物，他不能有效地建造正确的事物；除非他从听者心中拔出虚妄爱慕的荆棘，他必然徒然栽植神圣宣讲的话语。因此，伯多禄首先摧毁，为的是随后建造；他绝没有劝告犹太人他们现在该做什么，反而谴责他们所做的，说：\BibleQuote{诸位以色列人！请听这些话：纳匝肋的耶稣是天主用德能、奇迹和征兆向你们所证明的人，正如你们自己所知道的；祂按天主的既定计划和预知被交付，你们却借着不法之人的手钉死杀害了祂；天主却解除了地狱的痛苦，使祂复活了}（\BibleRef{宗 2:22}）；这样，为的是让他们因承认自己的残忍而被推倒后，能更加殷切地寻求神圣的建造，从而更受益地听见。因此他们也立刻回答说：\BibleQuote{诸位仁人弟兄，我们该作什么？}接着便对他们说：\BibleQuote{你们每人要悔改，并领受圣洗}（同上，37-38节）。这些建造的话语，他们必定会轻视，倘若他们没有首先有益地意识到自己被推倒的败坏。因此，扫禄当天上的光照射他时，并没有立刻听到他该做什么正确的事，而是听到他做错了什么。因为，当他跌倒在地，询问说：\BibleQuote{主，你是谁？}立刻有回答说：\BibleQuote{我是你所迫害的纳匝肋的耶稣。}当他立刻回答说：\BibleQuote{主，你要我做什么？}随即加上：\BibleQuote{起来，进城去，在那里会告诉你该做什么}（\BibleRef{宗 9:4}）。看，主从天上说话，谴责了迫害者的行为，却没有立刻指示他该做什么。看，他骄傲的整个结构已被推倒，然后，跌倒后谦卑了，他寻求被建造；而当骄傲被推倒时，建造的话语仍被保留；即是说，那残忍的迫害者可能长久躺卧在地，之后在善中更坚固地兴起，正如他从前的错误中彻底倾倒。因此，那些尚未开始行善的人，首先要被矫正的手从其不义的顽梗中推翻，以便之后被提升到行善的状态。出于这个原因，我们也砍伐森林中高大的树木，为能在建筑屋顶上竖起它；但它并不立即被放置在结构中；而是为了首先使其有害的湿气干燥：因为其湿气在最底部渗出得越多，它就被越坚固地提升到最高处。\par

但是，另一方面，那些完全没有完成他们已经开始之善事的人，应当受到劝诫，要谨慎地思考：当他们没有实现自己的目的时，他们如何连同已经开始了的事情也一并摧毁。因为，如果那被视为应当完成的事情，没有通过勤勉的实践而前进，那么即使是已经做得好的事情也会退步。因为人的灵魂在今世，就像一艘逆流而上的船：它决不能停留在原地，因为除非它奋力向上，否则就会漂回最下游。因此，如果工人的强手没有把已经开始之善事带向完美，那么工作中的懈怠本身就会反对已经完成的工作。因此，经上通过\Person{撒罗满}说：\BibleQuote{工作懈怠懒惰的人是浪费自己工作的弟兄}\BibleRef{箴 18:9}。的确，凡不努力执行他已开始之善事的人，在其疏忽的懈怠中，效法了破坏者的手。因此，天使给\Place{撒尔德}教会的信中说：\BibleQuote{你要警醒，坚固那剩下将要衰微的；因为我见你的行为，在我天主面前没有一样是完成的}\BibleRef{默 3:2}。因此，因为行为在他的天主面前没有被发现是完成的，他预言那些剩下的，甚至已经完成的，也快要死了。因为，如果我们内那死了的不被点燃进入生命，那被保留如同仍然活着的，也会被熄灭。他们应当受到劝诫：对他们来说，没有走上正路可能比走上后又转身离开更能容忍。因为他们若不回头看，就不会因任何麻木而对自己承担的目的变得软弱。那么，让他们听听经上所说的：\BibleQuote{他们认识正义之道以后，再背转身去，倒不如不认识它更好}\BibleRef{伯后 2:21}。让他们听听经上所说的：\BibleQuote{我巴不得你或冷或热；但是你既然温吞，也不冷也不热，我要把你从我口中吐出去}\BibleRef{默 3:15}。因为热的人是既承担又完成善事的人；冷的人是根本还没有开始任何需要完成之事的人。而正如通过温吞从冷过渡到热，同样，通过温吞也会从热返回到冷。因此，凡失去了不信之冷而得以存活，却丝毫没有越过温吞而继续燃烧的人，他无疑绝望于热，当他停留在有害的温吞中时，他正走上变冷的道路。但是，正如在温吞之前，冷中尚有希望，同样，在冷之后，温吞中只有绝望。因为尚在罪中的人并没有失去对悔改的信赖；但悔改后变得温吞的人，却收回了本来可以为罪人提供的希望。因此，每个人都必须是或冷或热，免得因温吞而被吐出来：也就是说，要么尚未皈依，仍能提供皈依的希望；要么已经皈依，就在德行中炽热；免得因温吞而被吐出，因为他从立定的热忱因麻木退回到有害的寒冷。\par

% structure: p0210 source=h2 target=subsection reason=relative-heading-rank; base=h2

\subsection{第三十五章}

\Emph{应如何劝诫那些暗中做坏事而公开做好事的人，以及那些相反行事的人。}\par

（\Emph{劝诫} 36.）应不同地劝诫那些暗中做坏事、公开做好事的人，与那些隐藏自己所做的好事，却因某些公开的行为让人以为他们是坏的人。因为那些暗中做坏事、公开做好事的人，应当受到劝诫，要思考人的判断何其迅速地消失，而天主的判断却何其恒久不变。他们应当受到劝诫，要定睛于事物的结局；因为人的赞许的见证会过去，而穿透隐密之事的属天判决，却坚固以达永远的报应。因此，当他们将隐藏的恶行置于天主的审判前，将正当的行为置于人眼前时，他们公开做的好事既没有见证，他们隐秘的过犯也并非没有永恒的见证。这样，通过向人隐藏他们的过失，而展示他们的德行，他们既在隐藏中暴露了他们应受惩罚的事，又在暴露中隐藏了他们可能获得赏报的事。这样的人，真理称之为\BibleQuote{粉刷的坟墓，外表好看，里面却装满死人的骨骸}\BibleRef{玛 23:17}；因为他们掩盖了内在的邪恶，却透过展示某些行为，仅以正义的外表颜色来取悦人的眼睛。因此，他们应当受到劝诫，不要轻视自己所行的正事，而要相信它们有更好的功绩。因为那些以为人的恩宠足以作为其赏报的人，极大地误判了他们自己的善事。因为当寻求易逝的赞美作为正确行为的回报时，一件值得永恒赏报的事就被贱卖了。关于这种回报，真理确实说：\BibleQuote{我实在告诉你们，他们已经获得了他们的赏报}\BibleRef{玛 6:2}。他们应当受到劝诫，要思考：当他们在隐密的事上显明自己是坏的，却在好行为上公开将自己作为榜样时，他们表明自己所逃避的是应当效法的；他们大声宣称自己所憎恨的是应当喜爱的；总之，他们为别人而活，却为自己而死。\par

但另一方面，那些在暗处行善，却在某些公开的事上容许人对自己有恶评的人，应当被劝告：在他们里面善的东西固然以行善的德行激励他们自己，但他们自己不应通过恶名声的榜样杀害他人；他们不应爱近人少于爱自己，也不应在自己汲取有益健康的酒时，向那些观察他们的心灵递出一杯毒害的毒酒。这些人确实一方面对近人的生命几乎无益，另一方面则大大加重其负担，因为他们既努力暗中行善，又在某些作为榜样的事上，从自身播下恶的种子。因为凡已能践踏称赞之欲的人，若隐瞒自己所做的善事，便是欺骗建树；他就如同撒种后窃取了发芽的根，因为他没有将可效仿的工作显示出来。因此，福音中真理说：\BibleQuote{好使他们看见你们的善行，光荣你们在天之父}\BibleRef{玛 5:16}。但随后又有这句话，看似吩咐了非常不同的事，即：\BibleQuote{你们应当小心，不要在人前表演你们的正义，为叫他们看见}\BibleRef{玛 6:1}。\par

那么，既命令我们的工作应做得不被人看见，又命令它应被看见，这是什么意思呢？不就是我们做的事要隐藏，以免我们自己受称赞，又要显扬，以增加我们在天之父的光荣吗？因为，当主禁止我们在人前表演我们的正义时，祂立刻加上：\BibleQuote{为叫他们看见}。又，当祂命令我们的善行应在人前被看见时，祂立刻附加说：\BibleQuote{好使他们光荣你们在天之父}\BibleRef{玛 5:16}。那么，祂在命令的结尾指明了它们应以何种方式被看见，以何种方式不被看见，即：工作者不应为自己的缘故寻求工作被看见，但因在天之父的光荣，他也不应隐瞒它。由此常发生：一件善工，当公开进行时，也是在暗中；当秘密进行时，又是在明处。因为那在公开善工中不寻求自己的，只求在天之父的光荣的人，隐藏了所做的事，因为他只以他所愿悦纳的那一位为见证。而那在秘密善工中贪求被观看和称赞的人，即使没有人看见他所做的事，也是在众人面前做的，因为他心里寻求多少人间的称赞，就为他的善工引进了多少证人。但是当恶名声，只要在未犯罪的情况下传开，没有被观望者从心中抹去时，就如同榜样，将罪的杯递给所有存恶念的人。因此也常发生：那些漫不经心让人对自己有恶评的人，自己固然不亲身作恶，但通过那些可能以他们为榜样的人，他们以更多的方式犯罪。因此，保禄对那些没有受玷污而吃某些不洁之物，但他们的吃却在软弱者面前放置了诱惑的绊脚石的人说：\BibleQuote{你们要谨慎，免得你们这自由竟成了软弱者的绊脚石}\BibleRef{格前 8:9}；又说：\BibleQuote{因你的良心，那软弱的弟兄，基督为他死了，就要丧亡。但你们这样得罪了弟兄们，并伤了他们软弱的良心，就是得罪基督}\BibleRef{格前 2:12}。因此，梅瑟说：\BibleQuote{不可咒骂聋子}，随即加上：\BibleQuote{也不可将绊脚石放在瞎子面前}\BibleRef{肋 19:14}。因为咒骂聋子就是毁谤一个不在场且听不见的人；而将绊脚石放在瞎子面前，则是虽然确有辨识地行事，却给没有辨识之光的人提供冒犯的理由。\par

% structure: p0215 source=h2 target=subsection reason=relative-heading-rank; base=h2

\subsection{第三十六章}

\Emph{关于对多人同时进行的劝告，要使劝告有助于每个人的德行，以免通过它使与这些德行相反的恶习滋长。}\par

身为灵魂牧者，在多样训导中应留意这些事，以便为不同听众的不适对症下药。然而，个别劝诫时依据各自需要提供服务，固然是极大的劳苦——因为要顾及每人的实况，给予应有的考虑，实属不易——但更困难的是，面对众多怀有各种情欲的听众，同时以同一劝言加以训诫。因为此时言辞需如此巧妙调配，使听者虽罪过各异，却都能各得其所，而言辞本身又不失其一致；它如同一柄双刃剑，一击穿透情欲的中心，从两侧削断肉身思虑的肿胀；这样，对骄傲者宣讲谦逊，却不致增加胆小者的恐惧；对胆小者注入信心，却不致助长骄傲者的纵肆；对怠惰麻木者宣讲行善的勤奋，却不致增加躁动者过度行为的放纵；对躁动者设定界限，却不致在怠惰者心中生出疏忽的麻木；对急躁者熄灭怒火，却不致在柔顺温和者心中生出疏忽；对柔顺温和者激发热心，却不致给易怒者火上浇油；对吝啬者灌输施舍的慷慨，却绝不放松挥霍者的缰绳；对挥霍者宣讲节俭，却不致增加吝啬者对易朽之物的执守；对不洁者赞美婚姻，却不致将已守贞者召回复归纵欲；对守贞者赞美身体的童贞，却不致让已婚者轻视肉体的生育之能。善事应如此宣讲，以免旁助恶事。最高的善应受赞美，以免最卑微者绝望。最卑微者应受珍视，以免因满足于最低限度而停止追求最高者。\par

% structure: p0218 source=h2 target=subsection reason=relative-heading-rank; base=h2

\subsection{第三十七章}

\Emph{对受相反情欲困扰之人的劝言。}\par

宣讲者在公开讲道中留意个人隐藏的情感和动机，并如角力者般灵活转向两侧，这诚然是艰巨的劳作；但当他被迫向一个受相反情欲困扰之人宣讲时，其辛劳更为沉重。因为常有这样的情况：某人天性过于欢愉，而突发的悲伤又不合时宜地压抑他。因此，宣讲者必须留意：当暂时之悲伤被消除时，不可因此增加其天性的欢愉；当天性之欢愉被抑制时，也不可加重暂时之悲伤。此人被过度的急躁习惯所困，但有时突然生出的恐惧又阻碍他急于行事。彼人被过度的恐惧习惯所困，但有时又被过度的急躁鲁莽所驱使，急欲求成。因此，在前者，应如此抑制突生的恐惧，以免其长期滋养的急躁加剧；在后者，应如此抑制突生的急躁，以免其天性植入的恐惧增强。的确，灵魂的医师在这些事上谨慎，又有何稀奇呢？何况那些医治身体而非心灵的人，也以极大的辨别智慧自我约束。常有这样的情况：虚弱的身体被极度的昏厥压倒，本应以强力药物应对；但虚弱的身体却无法承受强力药物。因此，治疗者留心：既要驱散新生的疾病，又不增加身体原有的虚弱，以免昏厥随生命一同消逝。他如此调和药物，同时兼顾昏厥与虚弱。那么，若不加区分使用的身体药物能分途发挥作用，为何同一讲道中应用的灵魂药物，就不能在多种方向上有能力应对道德疾病呢？这药物作用于无形之物，其运作更为精妙。\par

% structure: p0221 source=h2 target=subsection reason=relative-heading-rank; base=h2

\subsection{第三十八章}

\Emph{有时更轻微的罪过可暂置之不理，以便除去更严重的罪过。}\par

当两种恶习的疾病攻击一个人时，一种较轻，另一种也许更重，那么无疑应赶快支援那致使更快趋向死亡的一种。若不能阻止一种导致迫近的死亡，除非另一种相反恶习增长，那么宣讲者必须巧妙地管理他的劝诫，容忍一种增长，以便阻止另一种造成迫近的死亡。他这样做时，并非加重疾病，而是保全他所施药者的生命，以便寻找时机探求康复的方法。常有这样的人：他一方面毫不节制\OriginalTerm{贪饕}{gluttony}饮食，另一方面立即受到几乎要征服他的\OriginalTerm{迷色}{lechery}刺痛的逼迫；当他因害怕这场战斗而恐惧，努力通过\OriginalTerm{节制}{abstinence}约束自己时，却受\OriginalTerm{虚荣}{vain-glory}诱惑的折磨。在这种情况下，除非培养另一恶习，否则一个恶习绝不能被消灭。那么，应当更热切地攻击哪一种瘟疫呢？岂不是那更危险地压迫人的那一种？因为要容忍，借着节制的德行，暂时让骄傲在活人身上增长，以免通过贪饕，迷色将他完全从生命中剪除。因此，\Person{保禄}当他考虑到他那软弱的听众要么继续作恶，要么为行善而欢喜于人的称赞时，他说：\BibleQuote{你岂不怕那权柄么？你只管行善，就必得他的称赞。}\BibleRef{罗 13:3}\par

% structure: p0224 source=h2 target=subsection reason=relative-heading-rank; base=h2

\subsection{第三十九章}

\Emph{深奥之事决不可向软弱的灵魂宣讲。}\par

但宣讲者应知道如何避免将听众的心灵拉伸超过其力量，以免，可以说，灵魂的弦被拉得太紧而断裂。因为一切\OriginalTerm{深奥之事}{deep things}都应在众多听众面前隐藏，只向少数人稍作揭示。因此真理亲自说：\BibleQuote{你认为谁是那忠信及精明的管家，主人派他管理他的家仆，按时配给食粮呢？}\BibleRef{路 12:42}。所谓\Quote{一升麦子}是表达圣言的一份，以免当给予狭窄心灵超其容量时，被泼洒出去。因此\Person{保禄}说：\BibleQuote{我以前对你们说话，还不能把你们当作属神的人，只能当作属血肉的人，当作在基督内的婴孩。我给你们喝的是奶，不是饭。}\BibleRef{格前 3:1}。因此梅瑟从天主的圣所出来时，在百姓面前用帕子蒙上发光的脸；因为事实上他不向群众显示内在光明的奥秘\BibleRef{出 34:33}。因此天主的声音命令他：若有人挖掘蓄水池，而不盖好，牛或驴掉进去，他应赔偿\BibleRef{出 21:33}。因为当一个人达到了知识的深流，却不在愚钝的听众面前遮盖它们时，他应承担刑罚的责任，如果因他的话，一个灵魂（无论是洁净的还是不洁的）被绊脚石绊倒。因此对圣约伯说：\BibleQuote{谁给鸡有聪明？}\BibleRef{约 38:36}。因为圣洁的宣讲者在黑暗中大声呼喊，如同夜间啼叫的公鸡，当他说：\BibleQuote{现在是该由睡梦中醒来的时候了}\BibleRef{罗 13:11}。又说：\BibleQuote{你们应醒悟行义，不要犯罪}\BibleRef{格前 15:34}。但公鸡习惯在深夜高声啼唱；但当早晨临近时，它发出细小柔和的声音；因为事实上，正确宣讲的人向仍然在黑暗中的心灵清晰地呼喊，不向他们显示隐藏的奥秘，以便当他们接近真理之光时，能听到更精微的关于天上事物的教导。\par

% structure: p0227 source=h2 target=subsection reason=relative-heading-rank; base=h2

\subsection{第四十章}

\Emph{论宣讲的事工与声音}\par

但在这些事情中，我们被爱德的迫切渴望带回到上面已经说过的话：每位宣讲者应更多以行为而非言语发出声音，应更以善行留下可追随的足迹，而非以说话指示当行的道路。因为那只公鸡，主用比喻将其视为善宣讲者的代表，当它准备啼叫时，先抖动翅膀，拍打自己使自己更清醒；因为宣讲圣言的人必须先在善行的热忱中完全清醒，以免他们用声音唤醒别人，而自己在实行上却麻木不仁；他们应先以崇高的行为抖动自己，然后使他人关心善行；他们应先以思想的翅膀拍打自己；凡在自己内无益麻木之处，他们应以醒悟的调查发现，以严格的训诫改正，然后才能以言语规范他人的生活；他们应注意先以哀叹惩罚自己的过失，然后谴责他人应受惩罚之事；并且在发出劝诫言语之前，先以行为宣告他们即将要说的一切。\par

\SourceRecord{Translated by James Barmby.；Nicene and Post-Nicene Fathers, Second Series；Vol. 12.；Edited by Philip Schaff and Henry Wace.；Buffalo, NY: Christian Literature Publishing Co.,；1895.；<http://www.newadvent.org/fathers/36013.htm>.}

% structure: p0001 source=h1 target=section reason=page-title

\section{\WorkTitle{司牧典范（第四卷）}}

讲道者既妥善完成一切，应如何回归自身，以免生活或讲道使他自高自大。\par

然而，由于讲道常以适宜的方式大量倾泻时，讲者之心会因隐藏的自我炫耀之乐而自高，故需极大谨慎，使他以恐惧的撕咬来啃噬自己，免得他以药石治愈他人疾病，却因忽略自身健康而肿胀；免得他在帮助他人时遗弃自己，在扶起他人时跌倒。因为对某些人而言，美德的伟大往往成为他们丧亡的契机；使他们过度信赖自己的力量而失去警觉，因疏忽而意外死亡。美德与恶习相争；心灵因美德而自得其乐；于是行善者的灵魂抛却谨慎之惧，安于自信；这昏沉之魂，狡猾的诱惑者一一列举它所行的善事，以膨胀的思绪将它抬举，仿佛超乎一切。由此，在义德的审判者眼前，美德的记忆成了灵魂的陷阱；因为回想所行之善时，若在自视中抬高自己，便在谦逊之作者前跌倒。因此，骄傲之魂被说：\Quote{因你更美，下去，与未受割礼者同眠。}\BibleRef{则 32:19}；仿佛明言：因你凭美德之荣美自高，你便由这荣美被驱而跌倒。故此，以耶路撒冷为喻，那骄傲于美德之魂受责，当它说：\Quote{我在所赐予你的荣美中完全了，上主说，你倚靠你的美貌，因你的名声而行淫。}\BibleRef{则 16:14,15}。心灵因信赖己美而自高，即当它为美德的功绩而欣喜，在安然中自夸。但通过这信赖，它被引向行淫；因为灵魂被自己的思念欺骗时，恶灵占有它，以无数恶习的诱惑玷污它。但应注意，经上说：\Quote{你因你的名声而行淫。}因为当灵魂放弃对至高治理者的关注时，它立即寻求自身的赞美，开始将一切为彰显施予者赞美而领受的善物归于自己；它渴望散布自身声誉的光荣，忙于成为众人钦佩的对象。因此，它在自己的名声中行淫，因它离弃合法婚床的契约，在赞美的欲望中将自己卖与那玷污的灵。因此，达味说：\Quote{他把他们的美德交给俘虏，把他们的荣美交在敌人手中。}\BibleRef{咏 67:61}。美德被交付俘虏，荣美交在敌人手中，即当老仇敌因行善中的得意而控制了受骗的灵魂。然而，这种美德的得意虽不完全征服，却也试探选民之心。\par

但这心灵被高举时便遭遗弃，遭遗弃后便被召回恐惧。因此，达味又说：\Quote{我在富足中说过：我永不动摇。}\BibleRef{咏 29:6}。但稍后他补充了因信赖美德而自傲所受的苦：\Quote{你转面不顾我，我便惊慌。}\BibleRef{咏 5:8}。仿佛明言：我自以为在美德中强壮，但被遗弃时，我才知我的软弱何其大。因此他又说：\Quote{我发誓并坚定立志，要遵守你正义的典章。}\BibleRef{咏 118:106}。但因他无力持续所誓之守，立刻惊慌，发现了自己的软弱。于是，他立刻投靠祈祷的帮助，说：\Quote{我全然卑微；上主，求你照你的话使我存活。}\BibleRef{咏 118:107}。但有时，神圣治理在借恩赐提升灵魂之前，先唤起它对软弱的记忆，免得它因所领受的美德而自傲。因此，先知厄则克耳在受引导默观天上之事前，先被称为人子；好像上主明言劝告他：免得你因所见之事心中得意，当谨慎思量你是什么；好使你虽洞悉至高之事，却记得你是人，为叫你在被提至自我之外时，因软弱的缰绳而忧惧地回归自身。因此，当美德丰盈使我们愉悦时，灵魂之眼当回归自身的弱点，有益地自卑；当看那未行之事，而非所行之善；好使心灵因软弱之思而破碎，在谦逊之作者前更坚立于美德。为此，全能的天主虽在极大程度上成全治理者之心，仍常在某些微小部分留下不完全；使他们虽以奇妙美德闪耀，却因自身不完全而苦恼，绝不因大事自高，既在最低微之事上挣扎无法得胜，便不敢以主要成就自夸。\par

请看，善人，我如何因你申斥所加于我的必要，为显示牧者当如何，如画匠状美男子而拙劣无能；我指引他人至完美的岸，自己却仍在过犯的波浪中颠簸。但在此生沉船中，我恳求你，以你祈祷的木板托住我，使我的重量虽将我沉下，你功绩之手能将我举起。\par

\SourceRecord{Translated by James Barmby.；Nicene and Post-Nicene Fathers, Second Series；Vol. 12.；Edited by Philip Schaff and Henry Wace.；Buffalo, NY: Christian Literature Publishing Co.,；1895.；<http://www.newadvent.org/fathers/36014.htm>.}
